%!TEX TS-program = xelatex 
%!TEX encoding = UTF-8 Unicode






% Options for packages loaded elsewhere
% \PassOptionsToPackage{unicode}{hyperref}
% \PassOptionsToPackage{hyphens}{url}


%

% \documentclass[
% % % % % % ]{report}

% PEARSON
\documentclass[twoside,10pt,letterpaper,usenames]{newstyle-PearsonGeneric-7-38-ROMEO}
\usepackage[twoside]{geometry} % to set the dimensions of the page

% PEARSON
\geometry{                                  % setting all the necessary dimensions...
paperwidth=7.375in,                  % Do not change these settings.
paperheight=9.125in,                 % These set up the trim size for the book.
lmargin=1in,
rmargin=.875in,
% bmargin=.650in,  %% to allow space for the drop folio
bmargin=.9in,
tmargin=.95in,
width=5.5in,
% height=7.4in,  %% stealing space for the drop folio
height=7.15in,
marginparwidth=0in,
marginparsep=0in,
headheight=0.2in,
headsep=.25in,
footskip=.025in}





% PEARSON
%%%%%%%% Skip over this part.  These lines should remain unaltered.
%
%%    These lines input packages that are used for this template.
\usepackage[letter,cam,center]{crop-AlteredCamMarks-Romeo}  % to add trim marks; updated Nov 2015 to allow more space between marks and trim edge
\usepackage{graphicx} % for graphics layout

%%%%%%%%%%%%%%%%% ROMEO- Running Heads
% Set up RHs for frontmatter and Ch0; they will have to be altered again for Chs 1-3 and then revert to this again for Ch4
\usepackage{fancyhdr}
\usepackage{extramarks}
\newcommand{\cppxx}{C++xx}
\newcounter{cppxx}[chapter]
\pagestyle{fancy}
% for fancy pages, i.e., most pages
\rfoot[\fancyplain{}{\hspace{\fill}\\ \sfnine\thepage\hfill}]{\fancyplain{\hspace{\fill}\\ \sfnine\thepage}{\hspace{\fill}\\ \sfnine\thepage}}
\lfoot[\fancyplain{\bf\thepage}{}]{\fancyplain{}{}}
\cfoot[\fancyplain{}{}]{\fancyplain{}{}}
%\renewcommand{\sectionmark}[1]{\markright{{\sfnine\thesection\ #1}}}
%\renewcommand{\sectionmark}[1]{\markright{{\sfnine\ #1}}} % to remove section number
\renewcommand{\sectionmark}[1]{\markright{{\sfRHRomeo #1}}} % to remove section number 
%% ROMEO - making feature short title larger and bolder
%\renewcommand{\chaptermark}[1]{\markboth{{\sfb Chapter~\thechapter\quad {\sfb #1}}}{}}
\renewcommand{\chaptermark}[1]{\markboth{{\sfelevenRomeo Chapter~\thechapter\quad {\sfelevenRomeo #1}}}{}}
\renewcommand{\headrulewidth}{0pt}
%\renewcommand{\chaptersmark}[1]{\markboth{}{\sfb #1}}
\renewcommand{\chaptersmark}[1]{\markboth{}{\sfRHRomeo #1}}
% for plain pages, like chapter openers
\fancypagestyle{plain}{%
\fancyhf{} % clear all header and footer fields
\fancyfoot[RO]{\hspace*{\fill}\\\sfnine{\thepage}} 
\renewcommand{\headrulewidth}{0pt}
\renewcommand{\footrulewidth}{0pt}
}






%%%% CHANGED FOR ROMEO
% Changing code font to Cousine
\usepackage{mathspec}
\setmonofont[Ligatures=NoCommon,Scale=.85,SlantedFont={Cousine Italic},BoldFont={Cousine Bold},BoldSlantedFont={Cousine Bold Italic}]{Cousine Regular} % to set our code font
\newfontface{\SecCode}[Ligatures=NoCommon,Scale=1.3]{Cousine Bold}
%\newfontface{\SecCode}[Ligatures=NoCommon,Scale=1.105]{Cousine Bold}
\newfontface{\SubsecCode}[Ligatures=NoCommon,Scale=1.2]{Cousine Bold}
%\newfontface{\SubsecCode}[Ligatures=NoCommon,Scale=1.02]{Cousine Bold}
\newfontface{\SubsubsecCode}[Ligatures=NoCommon,Scale=1]{Cousine Bold}
%\newfontface{\SubsubsecCode}[Ligatures=NoCommon,Scale=.935]{Cousine Bold}
\newfontface{\ParaCode}[Ligatures=NoCommon,Scale=.9]{Cousine Bold}
\newfontface{\SubparaCode}[Ligatures=NoCommon,Scale=.9]{Cousine Regular}
%\newfontface{\RHCode}[Ligatures=NoCommon,Scale=.81]{Cousine Regular}
%\newfontface{\RHCode}[Ligatures=NoCommon,Scale=.85]{Cousine Bold}
%\newfontface{\RHCode}[Ligatures=NoCommon,Scale=1.17]{Cousine Bold} % set at scaled .9 of 13pts
\newfontface{\RHCode}[Ligatures=NoCommon,Scale=1.35]{Cousine Bold} % set at scaled .9 of 15pts
%\newfontface{\RHCode}[Ligatures=NoCommon,Scale=1.53]{Cousine Bold} % set at scaled .9 of 17pts
%\newfontface{\RHCode}[Ligatures=NoCommon,Scale=1.71]{Cousine Bold} % set at scaled .9 of 19pts
\newfontface{\RHCodeEleven}[Ligatures=NoCommon,Scale=.935]{Cousine Bold}
%\newfontface{\RHCodeTwelve}[Ligatures=NoCommon,Scale=1.02]{Cousine Bold}
\newfontface{\RHCodeTwelve}[Ligatures=NoCommon,Scale=1.1]{Cousine Bold}
%\newfontface{\TOCCode}[Ligatures=NoCommon,Scale=.85]{Cousine Bold}
%\newfontface{\TOCCode}[Ligatures=NoCommon,Scale=1.105]{Cousine Bold} % set at scaled .85 of 13pts
%\newfontface{\TOCCode}[Ligatures=NoCommon,Scale=1.0625]{Cousine Bold} % set at scaled .85 of 12.5pts
\newfontface{\TOCCode}[Ligatures=NoCommon,Scale=1.02]{Cousine Bold} % set at scaled .85 of 12pts
%\newfontface{\TOCCode}[Ligatures=NoCommon,Scale=.9]{Cousine Bold}
\newfontface{\TOCCodeR}[Ligatures=NoCommon,Scale=.85]{Cousine Regular}
\newfontface{\TOCCodeEleven}[Ligatures=NoCommon,Scale=.935]{Cousine Regular}
\newfontface{\TOCCodeElevenB}[Ligatures=NoCommon,Scale=.935]{Cousine Bold}
%\newfontface{\TOCCodeElevenB}[Ligatures=NoCommon,Scale=1.1]{Cousine Bold}
\newfontface{\TOCCodeTwelve}[Ligatures=NoCommon,Scale=1.02]{Cousine Regular}
\newfontface{\TOCCodeThirteen}[Ligatures=NoCommon,Scale=1.105]{Cousine Regular}

% Changing sans serif font to Gill Sans Nova
\setsansfont[Ligatures=Common,Scale=1,SlantedFont={Gill Sans Nova Medium Italic},BoldFont={Gill Sans Nova Bold},BoldSlantedFont={Gill Sans Nova Bold Italic}]{Gill Sans Nova Medium} % to set our sans font
%%
\newfontface{\sfbHugeRomeo}[Ligatures=Common,Scale=2.5]{Gill Sans Nova Bold}
\newfontface{\cmssbxparttocRomeo}[Ligatures=Common,Scale=1.2]{Gill Sans Nova Bold}
\newfontface{\cmssbxsectionRomeo}[Ligatures=Common,Scale=1.4]{Gill Sans Nova Bold}
\newfontface{\cmssbxelevenRomeo}[Ligatures=Common,Scale=1.1]{Gill Sans Nova Bold}
\newfontface{\cmssbxchaptocRomeo}[Ligatures=Common,Scale=1]{Gill Sans Nova Bold}
\newfontface{\cmssbxchaptitleRomeo}[Ligatures=Common,Scale=3.0]{Gill Sans Nova Bold}
\newfontface{\cmsschapnameRomeo}[Ligatures=Common,Scale=2.6]{Gill Sans Nova Medium}
\newfontface{\cmssbxpartRomeo}[Ligatures=Common,Scale=4.4]{Gill Sans Nova Bold}
\newfontface{\cmssparttitleRomeo}[Ligatures=Common,Scale=3.4]{Gill Sans Nova Medium}
\newfontface{\sfiHuge}[Ligatures=Common,Scale=2.5]{Gill Sans Nova Medium Italic}
\newfontface{\sfititleRomeo}[Ligatures=Common,Scale=2.4]{Gill Sans Nova Medium Italic}
\newfontface{\sfihalftitleRomeo}[Ligatures=Common,Scale=2.0]{Gill Sans Nova Medium Italic}
\newfontface{\sfiauthorRomeo}[Ligatures=Common,Scale=1.6]{Gill Sans Nova Medium Italic}
\newfontface{\sfblargeRomeo}[Ligatures=Common,Scale=1.2]{Gill Sans Nova Bold}
\newfontface{\sfbelevenRomeo}[Ligatures=Common,Scale=1.1]{Gill Sans Nova Bold}
\newfontface{\sfelevenRomeo}[Ligatures=Common,Scale=1.1]{Gill Sans Nova Medium}
\newfontface{\sftwelveRomeo}[Ligatures=Common,Scale=1.2]{Gill Sans Nova Medium}
\newfontface{\sftwelvefiveRomeo}[Ligatures=Common,Scale=1.25]{Gill Sans Nova Medium}
\newfontface{\sfthirteenRomeo}[Ligatures=Common,Scale=1.3]{Gill Sans Nova Medium}
\newfontface{\sfbRomeo}[Ligatures=Common,Scale=1]{Gill Sans Nova Bold}
\newfontface{\sfbiRomeo}[Ligatures=Common,Scale=1]{Gill Sans Nova Bold Italic}
\newfontface{\sfeightRomeo}[Ligatures=Common,Scale=.8]{Gill Sans Nova Medium}
\newfontface{\sfnineRomeo}[Ligatures=Common,Scale=.9]{Gill Sans Nova Medium}
\newfontface{\sfbnineRomeo}[Ligatures=Common,Scale=.9]{Gill Sans Nova Bold}
\newfontface{\defnemRomeo}[Ligatures=Common,Scale=1]{Gill Sans Nova Bold}
\newfontface{\sfbitwelveRomeo}[Ligatures=Common,Scale=1.2]{Gill Sans Nova Bold Italic}
%%%%%
\newfontface{\sfbsectionRomeo}[Ligatures=Common,Scale=1.3]{Gill Sans Nova Bold}
\newfontface{\sfbsectionitalRomeo}[Ligatures=Common,Scale=1.4]{Gill Sans Nova Bold Italic}
\newfontface{\sfbsubsecRomeo}[Ligatures=Common,Scale=1.2]{Gill Sans Nova Bold}
\newfontface{\sfbsubsubRomeo}[Ligatures=Common,Scale=1.1]{Gill Sans Nova Bold}
\newfontface{\sfbparaRomeo}[Ligatures=Common,Scale=1]{Gill Sans Nova Bold}
\newfontface{\sfbsubsecitalRomeo}[Ligatures=Common,Scale=1.2]{Gill Sans Nova Bold Italic}
\newfontface{\sfbsubsubsecitalRomeo}[Ligatures=Common,Scale=1.1]{Gill Sans Nova Bold Italic}
%%%
% scaled at 13pts
%\newfontface{\sfbiRHRomeo}[Ligatures=Common,Scale=1.3]{Gill Sans Nova Bold Italic}
%\newfontface{\sfiRHRomeo}[Ligatures=Common,Scale=1.3]{Gill Sans Nova Medium Italic}
%\newfontface{\sfRHRomeo}[Ligatures=Common,Scale=1.3]{Gill Sans Nova Medium}
% scaled at 15pts
\newfontface{\sfbiRHRomeo}[Ligatures=Common,Scale=1.5]{Gill Sans Nova Bold Italic}
\newfontface{\sfiRHRomeo}[Ligatures=Common,Scale=1.5]{Gill Sans Nova Medium Italic}
\newfontface{\sfRHRomeo}[Ligatures=Common,Scale=1.5]{Gill Sans Nova Medium}
%%%%%
% scaled at 17pts
%\newfontface{\sfbiRHRomeo}[Ligatures=Common,Scale=1.7]{Gill Sans Nova Bold Italic}
%\newfontface{\sfiRHRomeo}[Ligatures=Common,Scale=1.7]{Gill Sans Nova Medium Italic}
%\newfontface{\sfRHRomeo}[Ligatures=Common,Scale=1.7]{Gill Sans Nova Medium}
%%%%%% 
%% scaled at 19pts
%\newfontface{\sfbiRHRomeo}[Ligatures=Common,Scale=1.9]{Gill Sans Nova Bold Italic}
%\newfontface{\sfiRHRomeo}[Ligatures=Common,Scale=1.9]{Gill Sans Nova Medium Italic}
%\newfontface{\sfRHRomeo}[Ligatures=Common,Scale=1.9]{Gill Sans Nova Medium}


%% redefining existing sans serif font commands
\renewcommand{\sfbHuge}{\sfbHugeRomeo}       
\renewcommand{\cmssbxparttoc}{\cmssbxparttocRomeo}
\renewcommand{\cmssbxsection}{\cmssbxsectionRomeo}
\renewcommand{\cmssbxeleven}{\cmssbxelevenRomeo}
\renewcommand{\cmssbxchaptoc}{\cmssbxchaptocRomeo}    
\renewcommand{\cmssbxchaptitle}{\cmssbxchaptitleRomeo}
\renewcommand{\cmsschapname}{\cmsschapnameRomeo}
\renewcommand{\cmssbxpart}{\cmssbxpartRomeo}
\renewcommand{\cmssparttitle}{\cmssparttitleRomeo} 
\renewcommand{\sfiHuge}{\sfiHugeRomeo}
\renewcommand{\sfititle}{\sfititleRomeo}
\renewcommand{\sfihalftitle}{\sfihalftitleRomeo} 
\renewcommand{\sfiauthor}{\sfiauthorRomeo} 
\renewcommand{\sfblarge}{\sfblargeRomeo} 
\renewcommand{\sfbeleven}{\sfbelevenRomeo}  
\renewcommand{\sfb}{\sfbRomeo} 
\renewcommand{\sfeight}{\sfeightRomeo}
\renewcommand{\sfnine}{\sfnineRomeo}
\renewcommand{\defnem}{\defnemRomeo}                  
%%%%
\renewcommand{\sfbsection}{\sfbsectionRomeo}
\renewcommand{\sfbsubsec}{\sfbsubsecRomeo}
\renewcommand{\sfbsubsub}{\sfbsubsubRomeo}
\renewcommand{\sfbpara}{\sfbparaRomeo}


%%%%%%% strikeout text function
\usepackage[normalem]{ulem}  % used for a meredith citation
%%%%%%%%%%%%%%

%%%%%%%%%% getting picky about hyphenation
\hyphenation{name-space}
%%%%%%%%%%%%%

\newcommand\subsubemdash{\sfbsubsubRomeo\char"2014} %% for whatever reason, emdashes are not displaying properly in subsubsection heads. This creates a direct command to create them. 

\newcommand\romeovalue[1]{\textit{#1}}  % for normal text usage
\newcommand\romeovalueinside{\itshape} % for inside other environments, like table cells
\newcommand\romeogloss[1]{\textbf{#1}}
\newcommand\intraref[2]{\textit{\titleref{#1}~\intrarefdelim~\titleref{#2}} on page~\pageref{#2}}
\newcommand\intrarefsimple[1]{\textit{\titleref{#1}} on page~\pageref{#1}}
\newcommand\seealsoref[2]{``\titleref{#1}"~({#2}, p.~\pageref{#1})~\seealsodelim~}
\newcommand\featureref[2]{{#1}.``\titleref{#2}" on page~\pageref{#2}}
\newcommand{\intrarefdelim}{---}
\newcommand{\seealsodelim}{~$\blacklozenge$~}
\usepackage{ifthen}

\usepackage{cite}
%%%%%%%%%%%%%%%%% to set citations and biblabels bold 
\makeatletter
\renewcommand\@cite[2]{%
{\bfseries #1}}\ifthenelse{\boolean{@tempswa}}{,\nolinebreak[3] #2}{}
%\makeatletter
\renewcommand\@biblabel[1]{\hspace*{7pt}\MakeLowercase{\bfseries #1}}  % sets bib label citations are bold
\makeatother
%%%%%%%%%%%%%%%%

% Because the cppxx in the RH takes the moniker of "section" per the Au's requirement even though the Au requires that it not exist on the page and because each "feature" is a section and is unnumbered and because the Au decided when the design was almost complete that a section xref was needed for each feature, these manual commands will take the place of a working \ref for each cppxx fake section. Readdress this for the second edition and come up with a better solution; that will probably involve redefining the section command, so cppxx can be an actual section, or creating a new section-like environment for cppxx.   
\newcommand{\locationa}{Section~1.1}
\newcommand{\locationb}{Section~1.2}
\newcommand{\locationc}{Section~2.1}
\newcommand{\locationd}{Section~2.2}
\newcommand{\locatione}{Section~3.1}
\newcommand{\locationf}{Section~3.2}


%%%%%%%%%%CHANGED FOR ROMEO
\setcounter{secnumdepth}{0}




%%% LORI
\usepackage{cprotect}

\usepackage{amsfonts} % to ensure attractive mathematics
\usepackage{amssymb} % to ensure attractive mathematics
\usepackage{float}      % to allow things to float even if they normally don't
\usepackage{wrapfig} % to allow text to wrap around small figures
\usepackage{outline} % to allow outline creation
\usepackage{framed-PearsonGeneric} % to create the leftside vertical bar for the example environment

\usepackage{listings}  % for setting attractive computer code



%\definecolor{maroon}{cmyk}{.26,1,.63,.17}
%\definecolor{skyblue}{cmyk}{1,.69,0,0}  %% sky blue is now bright blue
%\definecolor{forestgreen}{cmyk}{.9,0,.9,0}
%\definecolor{mustard}{cmyk}{.0,.45,1,0}  %% mustard is now really orange
%\definecolor{purple}{cmyk}{.83,.81,0,0}
%\definecolor{codegray}{cmyk}{0,0,0,.5}  %% based on bright blue, not black 

%%%%% CHANGED FOR ROMEO
% changed font to Cousine, changed to 2ex indent

%% PRINT BOOK SETTINGS  % PRINT BOOK AND EBOOK WILL BE SET THE SAME
%\usepackage[xetex]{color}
%\usepackage{everypage-1x}
\usepackage{xespotcolor-Romeo}
\NewSpotColorSpace{PANTONE}
\AddSpotColor     {PANTONE}  {PANTONE661PC} {PANTONE\SpotSpace 661\SpotSpace PC} {1 0 0.05 0.72}
\SetPageColorSpace{PANTONE}
\definecolor{PANTONE661PC}  {spotcolor} {PANTONE661PC,1.0}
\definecolor{PANTONE661PCA} {spotcolor} {PANTONE661PC,0.5}
\definecolor{grey}           {gray}      {0.5}
\definecolor{grey}           {gray}      {0.5}


\makeatletter
\lstdefinestyle{Romeo}{
language=C++,
 framerule=0.5pt,
    keywordstyle={\ttfamily\bfseries},
    commentstyle={\color{PANTONE661PCA}\ttfamily},
    xleftmargin=2ex,
    escapeinside={(ù}{ù)},
    showstringspaces=false, 
  basicstyle=\ttfamily%
   \lst@ifdisplaystyle\small\fi
 }
\makeatother

\lstset{style=Romeo}
%{
%language=C++,
% framerule=0.5pt,
%    keywordstyle={\ttfamily\bfseries},
%    commentstyle={\color{PANTONE661PCA}\ttfamily},
%    xleftmargin=2ex,
%    escapeinside={(ù}{ù)},
%    showstringspaces=false, 
%}

%\lstset{style=Romeo,language=C++,
%basicstyle=\ttfamily\small,framerule=0.5pt,
%%emph=[1]{const},
%%  emphstyle=[1]{\color{PANTONE661PC}\ttfamily},
%%  moreemph=[2]{int,void,char},
%%  emphstyle=[2]{\color{PANTONE661PC}\ttfamily},
%%    moreemph=[3]{decltype},
%%      emphstyle=[3]{\color{PANTONE661PC}\ttfamily},
%%    keywordstyle={\color{PANTONE661PC}\ttfamily\bfseries},
%        keywordstyle={\ttfamily\bfseries},
%%  stringstyle={\color{PANTONE661PC}},
%%commentstyle={\color{PANTONE661PCA}\itshape},
%%commentstyle={\color{PANTONE661PCA}\itshape},
%commentstyle={\color{PANTONE661PCA}\ttfamily},
%%numberstyle=\color{PANTONE661PC},
%xleftmargin=2ex,
%    escapeinside={(ù}{ù)},
%    showstringspaces=false, 
%} % setting the code in small sans serif font and setting the margin for the code.  Nov 2015: setting the rules to be 0.5pt
%%%%%%%%%%%%
\newcommand{\codeincomments}{\color{PANTONE661PCA}\ttfamily}
\newcommand{\keyincomments}{\color{PANTONE661PCA}\ttfamily\bfseries}
\newcommand{\emphincomments}{\color{PANTONE661PCA}\ttfamily\itshape}
%%%%%%%%%%%%%%%%%%%%%%%%%%%%%

\lstdefinestyle{plain}{basicstyle=\ttfamily\small,framerule=0.5pt,
%emph=[1]{const},
%  emphstyle=[1]{\ttfamily},
%  moreemph=[2]{int,void,char},
%  emphstyle=[2]{\ttfamily},
%    moreemph=[3]{decltype},
%      emphstyle=[3]{\ttfamily},
    keywordstyle={\ttfamily},
%  stringstyle={},
commentstyle={\color{grey}\ttfamily},
%numberstyle={},
xleftmargin=2ex,
    escapeinside={(ù}{ù)},
    showstringspaces=false, 
}  

\lstdefinestyle{footcode}{language=C++,
basicstyle=\ttfamily\footnotesize,framerule=0.5pt,
keywordstyle={\ttfamily\bfseries\footnotesize},
commentstyle={\color{grey}\ttfamily},
xleftmargin=2ex,
escapeinside={(ù}{ù)},
showstringspaces=false, 
}  
%%%%%%%%%%%%%%%%%%%%%%%%%%%%%%%%%%%%%%%%
\newcommand{\fncodeincomments}{\color{grey}\ttfamily}
\newcommand{\fnkeyincomments}{\color{grey}\ttfamily\bfseries}
%%%%%%%%%%%%%%%%%


%%% EBOOK SETTINGS
%\lstset{language=C++,basicstyle=\ttfamily\small,framerule=0.5pt,
%emph=[1]{const},
%  emphstyle=[1]{\color{mustard}\ttfamily},
%  moreemph=[2]{int,void,char},
%  emphstyle=[2]{\color{maroon}\ttfamily},
%    moreemph=[3]{decltype},
%      emphstyle=[3]{\color{forestgreen}\ttfamily},
%    keywordstyle={\color{forestgreen}\ttfamily\bfseries},
%  stringstyle={\color{purple}},
%commentstyle={\color{codegray}\itshape},
%numberstyle=\color{purple},
%xleftmargin=2ex,
%    escapeinside={(ù}{ù)},
%    showstringspaces=false, 
%} %%%% setting the code in small sans serif font and
%%% setting the margin for the code.  Nov 2015: setting the rules to be 0.5pt
%
%\lstdefinestyle{plain}{basicstyle=\ttfamily\small,framerule=0.5pt,xleftmargin=2ex,
%emph=[1]{const},
%  emphstyle=[1]{\ttfamily},
%  moreemph=[2]{int,void,char},
%  emphstyle=[2]{\ttfamily},
%    moreemph=[3]{decltype},
%      emphstyle=[3]{\ttfamily},
%    keywordstyle={\ttfamily},
%  stringstyle={},
%commentstyle={\ttfamily},
%numberstyle={},
%    escapeinside={(ù}{ù)},
%    showstringspaces=false, 
%}
%
%%%%%%%%%%%%%%%%
%\newcommand{\codeincomments}{\color{codegray}\ttfamily}
%%%%%%%%%%%%%%%%%%%%%%%%%%%%%%%%%


%%%%%%%%%%%%%%%%%%%%%

\newcommand{\authorsnote}{Authors' Note: }
%\newcommand{\intraref}[1]{{\textit{#1}}}
\newcommand{\intrarefsub}[1]{\textit{#1}}  % needed for V's markdown conversion


\usepackage[stable,bottom]{footmisc} % to stabilize the footnote environment and to ensure footnotes
							  % appear at the bottom of the page, even if a figure or table
							  % appears at the bottom  


%\renewcommand{\thetable}{\thesection.\arabic{table}} 

					  
\usepackage[flushleft]{threeparttable} % to set attractive footnotes within table environments
\usepackage{titleref} % to allow cross-referencing by title
\usepackage{xr} % to allow cross-referencing between files
\usepackage{colortbl} % to allow the use of color/shading in tables
\usepackage{caption} % to set the appearance of captions
\captionsetup[table]{singlelinecheck=off,labelfont={sf,bf,small},font={sf,bf,small}}
\captionsetup[figure]{labelfont={sf,bf,small},font={sf,bf,small}}
\captionsetup[lstlisting]{labelfont={sf,bf,small},font={sf,bf,small}}
\usepackage{tocloft-hacked-PearsonGeneric-ROMEO} % to allow manipulation of the TOC
\frenchspacing

%%%%%%%%%% Creating C++xx box in the outside margin
%\usepackage{chngpage} %% no longer being used. 
%%%%%%%%%%%



%%%%%%%% Determine your indexing plan.
%
% Most of the time, the publisher will hire an independent indexing expert to create an
%    index for your book.  Your production editor can let you know the plan for your book.
% Use the line below only if your production plan includes the creation of a LaTeX index.
% Note that this situation is uncommon.
%\makeindex           % if your book will include a LaTeX-generated index
%\usepackage{See-makeidx}

%%%%%%%%%% These lines set up the Preamble of the document. You can skip over this part.
%% Preamble:
\input{definitions-PearsonGeneric-ROMEO}  % Loading the definitions needed for this book:



% Required for PANDOC conversion
\newcommand{\passthrough}[1]{#1}
%
%% Hack to allow PANDOC to parse Markdown in inline LaTeX environments
\let\Begin\begin
\let\End\end

% PANDOC stuff
%\IfFileExists{upquote.sty}{\usepackage{upquote}}{}
%\IfFileExists{microtype.sty}{% use microtype if available
%  \usepackage[]{microtype}
%  \UseMicrotypeSet[protrusion]{basicmath} % disable protrusion for tt fonts
%}{}
%\makeatletter
%\@ifundefined{KOMAClassName}{% if non-KOMA class
%  \IfFileExists{parskip.sty}{%
%    \usepackage{parskip}
%  }{% else
%    \setlength{\parindent}{0pt}
%    \setlength{\parskip}{6pt plus 2pt minus 1pt}}
%}{% if KOMA class
%  \KOMAoptions{parskip=half}}
%\makeatother
%\usepackage{xcolor}
% \IfFileExists{xurl.sty}{\usepackage{xurl}}{} % add URL line breaks if available

% \IfFileExists{bookmark.sty}{\usepackage{bookmark}}{\usepackage{hyperref}}
% \hypersetup{
% % % % % % %   hidelinks,
% %   pdfcreator={LaTeX via pandoc}}


% \urlstyle{same} % disable monospaced font for URLs


% 
% 


%\setlength{\emergencystretch}{3em} % prevent overfull lines
%\providecommand{\tightlist}{%
%\setlength{\itemsep}{0pt}\setlength{\parskip}{0pt}}

% 




%%%%%%%%% Now the document begins

\begin{document}
%%%%%% commands to fine-tune to table of contents
%\cftnodots
%\renewcommand{\@dotsep}{10000}


\renewcommand{\cfttoctitlefont}{\cmsschapname}
\renewcommand{\cftpartfont}{\cmssbxparttoc}
\renewcommand{\cftchappagefont}{\cmssbxparttoc}
\renewcommand{\cftchapfont}{\cmssbxchaptoc}
\renewcommand{\cftchapnumwidth}{5.5em}
\renewcommand{\cftchappagefont}{\cmssbxchaptoc}
\renewcommand{\cftsecfont}{\sf}
\renewcommand{\cftsecindent}{1.25em}
%\renewcommand{\cftsecnumwidth}{0em}
\renewcommand{\cftsecpagefont}{\sf}
\renewcommand{\cftsubsecfont}{\sf}
\renewcommand{\cftsubsecindent}{4em}
%\renewcommand{\cftsubsecnumwidth}{0em}
\renewcommand{\cftsubsecpagefont}{\sf}
\renewcommand{\contentsname}{{Contents}} 
\setcounter{tocdepth}{1}



\setcounter{secnumdepth}{0} 
%\renewcommand{\numberline}[2]{} 
%%%%% end toc commands % Loading a few commands related to the tocloft package



\renewcommand{\thelstlisting}{\arabic{lstlisting}}  %%% Romeo - reset listings per section

%%%%%%%%%%% Creating C++xx box in the outside margin  -- no longer being used
%\newcommand{\Cbox}{%\checkoddpage
%\ifcpoddpage\begin{marginpar}
%{{\hspace*{.5em}\color{forestgreen}\fbox{\sfb\cppxx}}}
%\end{marginpar}
%\else\begin{marginpar}
%{{\hspace{-4.25em}\color{forestgreen}\fbox{\sfb\cppxx}}} 
%\end{marginpar}\fi}
%%%%%%%%%%%%%%%%%%%%%%%%%


%%%%% The front matter begins:
\pagenumbering{roman}             % Roman numbering
\setcounter{page}{1}                    % Starting page number (specified by the publisher)
\cleardoublepage                         % This command is defined in our documentclass file. It tells 
                                                      % our compiler that we want the next item 
                                                       % to start on a right-hand page, and it inserts a blank left-hand 
                                                       % page if needed.    
\thispagestyle{empty}
\begin{center}
{\sfititle Title of the Book}
\end{center}                        % The half-title page
%\include{series}                             % A placeholder for the series page. The publisher normally 
%                                                       % creates this page and just adds it into the final PDF. 
\cleardoublepage
\thispagestyle{empty}
\begin{center}
{\sfititle Embracing Modern C++ Safely}\\
\vspace*{54pt}
%{\sfihalftitle Subtitle of the Book}\\
\vspace*{\fill}    % This space will automatically adjust
\noindent{\sfiauthor John Lakos}\\[3ex]
\noindent{\sfiauthor Vittorio Romeo}\\[3ex]        % Add or delete lines as necessary for author names
\noindent{\sfiauthor Rostislav Klebnikov}\\[3ex]
\noindent{\sfiauthor Alisdair Meredith}\\[3ex]

\vspace*{\fill}    % This space will automatically adjust
% Authors:  For the next line, use either the AW line or the PH line per your 
%   production editor's instructions. You will find both PC and Mac versions of the logos 
%   in the template folder. Use the one appropriate for your system.
%\includegraphics{PH_title_page_Mac} \\   %PH Mac
%\includegraphics{PH_stack_BLACK-PC} \\  %PH PC
\includegraphics{AWlogoTitlePage-Mac}\\   % AW Mac
%\includegraphics{AWlogoTitlePage_PC}  \\  %AW PC
{\footnotesize 
Boston~$\bullet$~Columbus~$\bullet$~New York~$\bullet$~San Francisco~$\bullet$~Amsterdam~$\bullet$~Cape Town\\[1ex]
Dubai~$\bullet$~London~$\bullet$~Madrid~$\bullet$~Milan~$\bullet$~Munich~$\bullet$~Paris~$\bullet$~Montreal~$\bullet$~Toronto~$\bullet$~Delhi~$\bullet$~Mexico City\\[1ex]
S\~{a}o Paulo~$\bullet$~Sydney~$\bullet$~Hong Kong~$\bullet$~Seoul~$\bullet$~Singapore~$\bullet$~Taipei~$\bullet$~Tokyo\\[1ex]
}
\end{center}                               % The title page
\thispagestyle{empty}
\vspace*{\fill}
{\footnotesize\raggedright 
\noindent Many of the designations used by manufacturers and sellers to distinguish their products are claimed as trademarks. Where those designations appear in this book, and the publisher was aware of a trademark claim, the designations have been printed with initial capital letters or in all capitals.\\
\hspace{\fill}\\
\noindent The authors and publisher have taken care in the preparation of this book, but make no expressed or implied warranty of any kind and assume no responsibility for errors or omissions. No liability is assumed for incidental or consequential damages in connection with or arising out of the use of the information or programs contained herein.\\   % change to authors (plural) if this book has multiple authors
\hspace{\fill}\\
\noindent For information about buying this title in bulk quantities, or for special sales opportunities (which may include electronic versions; custom cover designs; and content particular to your business, training goals, marketing focus, or branding interests), please contact our corporate sales department\linebreak[4] at \mbox{corpsales}@\mbox{pearsoned.com} or (800) 382-3419.\\
\hspace{\fill}\\
\noindent For government sales inquiries, please contact governmentsales@pearsoned.com.\\
\hspace{\fill}\\
\noindent For questions about sales outside the U.S., please contact intlcs@pearson.com.\\
\hspace{\fill}\\
\noindent Visit us on the Web: informit.com/aw\\     % use this line if you are publishing with AW
%\noindent Visit us on the Web: informit.com/ph\\   % use this line if you are publishing with PH
\hspace{\fill}\\
\noindent Library of Congress Control Number: TO COME FROM ITP\\  % update this with info provided by your production editor
\hspace{\fill}\\
\noindent Copyright~\copyright~2021 Pearson Education, Inc.\\
\hspace{\fill}\\
\noindent Cover image: [TO COME FROM ITP]\\
\hspace*{\fill}\\
\noindent All rights reserved. This publication is protected by copyright, and permission must be obtained from the publisher prior to any prohibited reproduction, storage in a retrieval system, or transmission in any form or by any means, electronic, mechanical, photocopying, recording, or likewise. For information regarding permissions, request forms and the appropriate contacts within the Pearson Education Global Rights \& Permissions Department, please visit www.pearson.com/permissions.\\
\hspace*{\fill}\\
\noindent ISBN-13: 978-0-13-738035-0\\[-.5ex] % update this with info provided by your production editor 
\noindent ISBN-10: 0-13-738035-6 \\ % update this with info provided by your production editor
\hspace*{\fill}\\
\noindent{\sf ScoutAutomatedPrintCode}
%\noindent Text printed in the United States on recycled paper at PRINTER INFO HERE.\\[-.5ex] % update this with info provided by your production editor
%\noindent First printing, MONTH YEAR % update this with info provided by your production editor
}                                % The cataloging-in-publication page
\thispagestyle{empty}
\begin{center}The lines of text of the dedication will be broken\\ 
with the help of the production team\\ 
to ensure an attractive\\
layout.\\[1ex] 
Author Initials\\
\vspace*{3pc}
Should there be a second author, that author's lines of\\  % These lines are 
dedication text will go here and will also be broken\\        % commented out 
in whatever way is attractive and logical.\\[1ex]                % if unneeded.
Author Initials\\
\end{center}                      % The dedication page
\cleardoublepage
\parskip .5ex                       % Add space between Contents items
\tableofcontents                  % Make the Table of Contents 
%\input{EMCS-testfile-replacetoc}
%\input{EMCS-testfileX.tex.toc}  % For testing Slava's magic toc fix
\cleardoublepage                
%\parskip 0in                         % Reset to zero interpar spacing
\setlength{\parindent}{0pt}  %% ROMEO
\parskip 1ex  %%%% ROMEO
%\cleardoublepage
\chapter*{Foreword}
\markboth{Foreword}{Foreword}

The text of the foreword will go here.                 % Add foreword
\include{preface}                 % Add preface 
\cleardoublepage
\chapter*{Acknowledgments}
\markboth{Acknowledgments\hfill}{\hfill\/Acknowledgments}

\noindent The text of the author's acknowledgements will go here.                         % Add acknowledgements 
\include{author-LakosRomeo}                   % Add about the author page
\cleardoublepage                  
%%%%%%%The front matter ends

\pagenumbering{arabic}		  % Arabic numbering
%%\part{This is the Part Title}
\vspace*{4pt}
\sffamily % this sets this page in sans serif font

\noindent Chapter \ref{ExampleChap}, "\titleref{ExampleChap},"  is followed by a one-sentence description of each chapter. 

Chapter \ref{samplechap2}, "\titleref{samplechap2},"  is followed by a one-sentence description of each chapter and if the sentence is longer what happens when it runs to two lines.

Chapter \ref{samplechap3}, "\titleref{samplechap3}," is followed by a one-sentence description of each chapter and if the sentence is longer what happens when it runs to two lines. 

\normalfont % this resets the fonts to our normal settings




            % Part I
\cleardoublepage
%\addtocontents{\makeatletter\renewcommand{\cftsubsecindent}{8em}\makeatother} % placement here actually puts the command just before Ch1 in toc file
\setcounter{chapter}{-1}
\chapter[Introduction]{Introduction}\label{ch-intro}
% Andrei's version

Welcome! \textit{Embracing Modern C++ Safely} is a \textit{reference book} dedicated to professionals who want to leverage modern C++ features in the development and maintenance of large-scale, complex C++ software systems.

This book deliberately concentrates on the productive value afforded by each new language feature added by C++ starting with C++11, particularly when the systems and organizations involved are considered at scale. We left aside ideas and idioms, however clever and intellectually intriguing, that could hurt the bottom line when applied at large. Instead, we focus on what is objectively true and relevant to making wise economic and design decisions, with an understanding of the inevitable tradeoffs that arise in any engineering discipline. In doing so, we do our best to steer clear of subjective opinions and recommendations.

Richard Feynman famously said: ``If it disagrees with experiment, it’s wrong. In that simple statement is the key to science."\footnote{Richard Feynman, lecture at Cornell University, 1964. Video and commentary available at https://\linebreak[3]fs.\linebreak[3]blog/\linebreak[3]2009/12/mental-model-scientific-method.} There is no better way to experiment with a language feature than letting time do its work. We took that to heart by dedicating \textit{Embracing Modern C++ Safely} to only the features of Modern C++ that have been part of the Standard for at least five years, which grants enough perspective to properly evaluate its practical impact. Thus, we are able to provide you with a thorough and comprehensive treatment based on practical experience and worthy of your limited professional development time. If you're out there looking for tried and true ways to better use modern C++ features for improving your productivity, we hope this book will be the one you'll reach for.

What's missing from a book is as important as what's present. \textit{Embracing Modern C++ Safely} is not a tutorial on programming, on C++, or even on new features of C++. We assume you are an experienced developer, team lead, or manager, that you already have a good command of ``classic" C++98/03, and that you are looking for clear, goal-driven ways to integrate modern C++ features within your and your team's toolbox.

\section{What Makes This Book Different}

The book you're now reading aims very strongly at being objective, empirical, and practical. We simply present features, their applicability, and their potential pitfalls as reflected by the analysis of millions of human-hours of using C++11 and C++14 in the development of varied large-scale software systems; personal preference matters have been neutralized to our, and our reviewers', best ability. We wrote down the distilled truth that remains, which should shape your understanding of what modern C++ has to offer to you without being skewed by our subjective opinions or domain-specific inclinations.

The final analysis and interpretation of what is appropriate for your context is left to you, the reader. Hence, this book is, by design, not a C++ style or coding-standards guide; it would, however, provide valuable input to any development organization seeking to author or enhance one.

Practicality is a topic very important to us, too, and in a very real-world, economic sense. We examine modern C++ features through the lens of a large company developing and using software in a competitive environment. In addition to showing you how to best utilize a given C++ language feature in practice, our analysis takes into account the costs associated with having that feature employed routinely in the ecosystem of a software development organization. (We believe that costs of using language features are sadly neglected by most texts.) In other words, we weigh the benefits of successfully using a feature against the risk of its widespread ineffective use (or misuse) and/or the costs associated with training and code review required to reasonably ensure that such ill-conceived use does not occur. We are acutely aware that what applies to one person or small crew of like-minded individuals is quite different from what works with a large, distributed team. The outcome of this analysis is our signature categorization of features in terms of safety of adoption --- namely \textit{safe}, \textit{conditionally safe}, or \textit{unsafe} features.

We are not aware of any similar text amid the rich offering of C++ textbooks; in a very real sense, we wrote it because we needed it.

\section{Scope for the First Edition}

Given the vastness of C++'s already voluminous and rapidly growing standardized libraries, we have chosen to limit this book’s scope to just the language features themselves. A companion book, \textit{Embracing Modern C++ Standard Libraries Safely}, is a separate project that we hope to tackle in the future. However, to be effective, this book must remain small, concise, and focused on what expert C++ developers need to know well to be successful right now.

In this first of an anticipated series of periodically extended volumes, we characterize, dissect, and elucidate most of the modern language features introduced into the C++ Standard starting with C++11. We chose to limit the scope of this first edition to only those features that have been in the language Standard and widely available in practice for at least five years. This limited focus enables us to more fully evaluate the real-world impact of these features and to highlight any caveats that might not have been anticipated prior to standardization and sustained, active, and widespread use in industry.

We assume you are quite familiar with essentially all of the basic and important special-purpose features of classic C++98/03, so in this book we confined our attention to just the subset of C++ language features introduced in C++11 and C++14. This book is best for you if you need to know how to safely incorporate C++11/14 language features into a predominately C++98/03 code base, today.

Over time, we expect, and hope, that practicing senior developers will emerge entirely from the postmodern C++ era. By then, a book that focuses on all of the important features of modern C++ would naturally include many of those that were around before C++11. With that horizon in mind, we are actively planning to cover pre-C++11 material in future editions. For the time being, however, we highly recommend \textit{Effective C++} by Scott Meyers\footnote{\textbf{meyers05}} as a concise, practical treatment of many important and useful C++98/03 features.

\section[The {\sffamily\itshape EMC++S} White Paper]{The {\sfbsectionitalRomeo EMC++S} White Paper}

Throughout the writing of \textit{Embracing Modern C++ Safely}, we have followed a set of guiding principles, which collectively drive the style and content of this book.

\subsection{Facts (Not Opinions)}

This book describes only beneficial uses and potential pitfalls of modern C++ features. The content presented is based on objectively verifiable facts, either derived from standards documents or from extensive practical experience; we explicitly avoid subjective opinion such as our evaluation of the relative merits of design tradeoffs (restraint that admittedly is a good exercise in humility). Although such opinions are often valuable, they are inherently biased toward the author's area of expertise.

Note that \textit{safety} --- the rating we use to segregate features by chapter --- is the one exception to this objectivity guideline. Although the analysis of each feature aims at being entirely objective, its chapter classification --- indicating the relative safety of its quotidian use in a large software-development environment --- reflects our combined accumulated experience totaling decades of real-world, hands-on experience with developing a variety of large-scale C++ software systems.

\subsection{Elucidation (Not Prescription)}

We deliberately avoid prescribing any cut-and-dried solutions to address specific feature pitfalls. Instead, we merely describe and characterize such concerns in sufficient detail to equip you to devise a solution suitable for your own development environment. In some cases, we might reference techniques or publicly available libraries that others have used to work around such speed bumps, but we do not pass judgment about which workaround should be considered a best practice.

\subsection{Brevity (Not Verbosity)}

\textit{Embracing Modern C++ Safely} is neither designed nor intended to be an introduction to modern C++. It is a handy reference for experienced C++ programmers who may have a passing knowledge of the recently added C++ features and a desire to perfect their understanding. Our writing style is intentionally tight, with the goal of providing you with facts, concise objective analysis, and cogent, real-world examples. By doing so we spare you the task of wading through introductory material. If you are entirely unfamiliar with a feature, we suggest you start with a more elementary and language-centric text such as \textit{The C++ Programming Language} by Bjarne Stroustrup.\footnote{\textbf{stroustrup13}}

\subsection{Real-World (Not Contrived) Examples}

We hope you will find the examples in this book useful in multiple ways. The primary purpose of examples is to illustrate productive use of each feature as it might occur in practice. We stay away from contrived examples that give equal importance to seldom-used aspects of the feature, as to the intended, idiomatic uses. Hence, many of our examples are based on simplified code fragments extracted from real-world codebases. Though we typically change identifier names to be more appropriate to the shortened example (rather than the context and the process that led to the example), we keep the code structure of each example as close as possible to its original real-world counterpart.

\subsection{At Scale (Not Overly Simplistic) Programs}

By scale, we attempt to simultaneously capture two distinct aspects of size: (1) the sheer product size (e.g., in bytes, source lines, separate units of release) of the programs, systems, and libraries developed and maintained by a software organization; and (2) the size of an organization itself as measured by the number of software developers, quality assurance engineers, site reliability engineers, operators, and so on that the organization employs. As with many aspects of software development, what works for small programs simply doesn’t scale to larger development efforts.

What’s more, powerful new language features that are handled perfectly well by a few expert programmers working together in the archetypal garage on a prototype for their new start-up don't always fare as well when they are wantonly exercised by numerous members of a large software development organization. Hence, when we consider the relative safety of a feature, as defined in the next section, we do so with mindfulness that any given feature might be used, and occasionally misused, in very large programs and by a very large number of programmers having a wide range of knowledge, skill, and ability.

\section[{What Do We Mean by {\sffamily\itshape Safely}?}]{What Do We Mean by {\sfbsectionitalRomeo Safely}?}

The ISO C++ Standards Committee, of which we are members, would be remiss --- and downright negligent --- if it allowed any feature of the C++ language to be standardized if that feature were not reliably safe when used as intended. Still, we have chosen the word ``safely" as the moniker for the signature aspect of our book, by which we indicate a comparatively favorable risk-to-reward ratio for using a given feature in a large-scale development environment. By contextualizing the meaning of the term ``safe," we get to apply it to a real-world economy in which everything has a cost in multiple dimensions: risk of misuse, added maintenance burden borne by using a new feature in an older code base, and training needs for developers who might not be familiar with that feature.

Several aspects conspire to offset the value added by the adoption and widespread use of any new language feature, thereby reducing its intrinsic safety. By categorizing features in terms of safety, we strive to capture an appropriately weighted combination of the following factors:
\begin{enumerate}
\item{Number and severity of known deficiencies}
\item{Difficulty in teaching consistent proper use}
\item{Experience level required for consistent proper use}
\item{Risks associated with widespread misuse}
\end{enumerate}
Bottom line: In this book, the degree of safety of a given feature is the relative likelihood that widespread use of that feature will have positive impact and no adverse effect on a large software company’s codebase.

\section[A {\sffamily\itshape Safe} Feature]{A {\sfbsectionitalRomeo Safe} Feature}

Some of the new features of modern C++ add considerable value, are easy to use, and are decidedly hard to misuse unintentionally; hence, ubiquitous adoption of such features is productive, relatively unlikely to become a problem in the context of a large-scale development organization, and to be generally encouraged --- even without training. We identify such staunchly helpful, unflappable C++ features as \textit{safe}.

For example, we categorize the \texttt{override} contextual keyword as a safe feature because it prevents bugs, serves as documentation, cannot easily be misused, and has no serious deficiencies. If someone has heard of this feature and tried to use it and the software compiles, the code base is likely better for it. Using \texttt{override} wherever applicable is always a sound engineering decision.

\section[A {\sffamily\itshape Conditionally Safe} Feature]{A {\sfbsectionitalRomeo Conditionally Safe} Feature}

The preponderance of new features available in modern C++ has important, frequently occurring, and valuable uses, yet how these features are used appropriately, let alone optimally, might not be obvious. What’s more, some of these features are fraught with inherent dangers and deficiencies, requiring explicit training and extra care to circumnavigate their pitfalls.

For example, we deem default member initializers a \textit{conditionally safe} feature because, although they are easy to use per se, the perhaps less-than-obvious unintended consequences of doing so (e.g., tight compile-time coupling) might be prohibitively costly in certain circumstances (e.g., might prevent relink-only patching in production).

\section[An {\sffamily\itshape Unsafe} Feature]{An {\sfbsectionitalRomeo Unsafe} Feature}

When an expert programmer uses any C++ feature appropriately, the feature typically does no direct harm. Yet other developers --- seeing the feature’s use in the code base but failing to appreciate the highly specialized or nuanced reasoning justifying it --- might attempt to use it in what they perceive to be a similar way, yet with profoundly less desirable results. Similarly, maintainers may change the use of a fragile feature altering its semantics in subtle but damaging ways.

Features that are classified as unsafe are those that might have valid, and even very important, use cases, yet our experience indicates that routine or widespread use thereof would be counterproductive in a typical large-scale software-development enterprise.

For example, we deem the final contextual keyword an unsafe feature because the situations in which it would be misused overwhelmingly outnumber those vanishingly few isolated cases in which it is appropriate, let alone valuable. Furthermore, its widespread use would inhibit fine-grained (e.g., hierarchical) reuse, which is critically important to the success of a large organization.

\section{Modern C++ Feature Catalog}

As an essential aspect of its design, this first edition of \textit{Embracing Modern C++ Safely} aims to serve as a comprehensive catalog of C++11 and C++14 language features, presenting vital information for each of them in a clear, concise, consistent, and predictable format to which experienced engineers can readily refer during development or technical discourse.

\subsection{Organization}

This book is divided into five chapters, the middle three of which form the catalog characterizing modern C++ language features grouped by their respective safety classifications:
\begin{itemize}
\item{ Chapter 0: Introduction}
\item{Chapter 1: Safe Features}
\item{Chapter 2: Conditionally Safe Features}
\item{Chapter 3: Unsafe Features}
\item{Chapter 4: Parting Thoughts}
\end{itemize}
For this first edition, the language-feature chapters (1, 2, and 3) each consist of two sections containing, respectively, C++11 and C++14 features having the safety level (\textit{safe}, \textit{conditionally safe}, or \textit{unsafe}) corresponding to that chapter. Recall, however, that Standard Library features are outside the scope of this book.

Each feature resides in its own subsection, rendered in a canonical format:
\begin{itemize}
\item{Description}
\item{Use Cases}
\item{Potential Pitfalls}
\item{Annoyances}
\item{See Also}
\item{Further Reading}
\end{itemize}
By constraining our treatment of each individual feature to this canonized format, we avoid gratuitous variations in rendering, thereby facilitating rapid discovery of whatever particular aspects of a given language feature you are searching for.

\section{How To Use This Book}

Depending on your needs, \textit{Embracing Modern C++ Safely} can be handy in a variety of ways.
\begin{enumerate}
\item{\textbf{Read the entire book from front to back.} If you are conversant with classic C++, consuming this book in its entirety all at once will provide a complete and nuanced practical understanding of each of the language features introduced by C++11 and C++14.}
\item{\textbf{Read the chapters in order but slowly over time.} An incremental, priority-driven approach is also possible and recommended, especially if you're feeling less sure-footed. Understanding and applying first the safe features of Chapter 1 gets you the low-hanging fruit. In time, the conditionally safe features of Chapter 2 will allow you to ease into the breadth of useful modern C++ language features, prioritizing those that are least likely to prove problematic.}
\item{\textbf{Read the first sections of each of the three catalog chapters first.} If you are a developer whose organization uses C++11 but not yet C++14, you can focus on learning everything that can be applied now and then circle back and learn the rest later when it becomes relevant to your evolving organization.}
\item{\textbf{Use the book as a quick-reference guide if and as needed.} Random access is great, too, especially now that you've made it through Chapter 0. If you prefer not to read the book in its entirety (or simply want to refer to it periodically as a refresher), reading any arbitrary individual feature subsection in any order will provide timely access to all relevant details of whichever feature is of immediate interest.}
\end{enumerate}

We wish you would derive value in several ways from the knowledge imbued into \textit{Embracing Modern C++ Safely}, irrespective of how you read it. In addition to helping you become a more knowledgeable and therefore safer developer, this book aims to clarify (whether you are a developer, a lead, or a manager) which features demand more training, attention to detail, experience, peer review, and such. The factual, objective presentation style also makes for excellent input into the preparation of coding standards and style guides that suit the particular needs of a company, project, team, or even just a single discriminating developer (which, of course, we all aim at being). Finally, any C++ software development organization that adopts this book will be taking the first steps toward leveraging modern C++ in a way that maximizes reward while minimizing risks, i.e., by embracing modern C++ \textit{safely}. We are very much looking forward to getting feedback and suggestions for future editions of \textit{Embracing Modern C++ Safely} at www.TODOTODOTODO.com. Happy coding!

  % Chapter 0
\cleardoublepage

%%%%%% RHs redefined for Chs 1-3
%\fancyhead[RO]{\sectionmark}
\fancyhead[LO]{\sfelevenRomeo Section~\thechapter.\thecppxx\hspace*{1em}\cppxx} 
%\fancyhead[RE]{\sectionmark}
%\fancyhead[LE]{\chaptermark}
%%%%%%%%%%
%\addtocontents{toc}{\protect{\protect\setlength{\cftsecindent}{4em}}}

\chapter[Safe Features]{Safe Features}\label{ch-safe}

Intro text should be here.


%%%%%%%%%%%% C++11
\stepcounter{cppxx}
\cftaddnumtitleline{toc}{section}{\thechapter.\thecppxx}{C++11}{\thepage}
% to add to the TOC at the section level
\renewcommand{\cppxx}{C++11}
\addtocontents{toc}{\protect\renewcommand*\protect\cftsecindent{4.5em}}
%\addtocontents{toc}{\protect\renewcommand*\protect\cftsubsecindent{5em}}

%\newpage
%%\cftaddtitleline{toc}{section}{\protect\hspace*{-5em}\framebox{\sffamily\itshape\cppxx}\vspace*{-3ex}}{}
%%\section[Attributes]{Attributes\checkoddpage\Cbox}\label{attributes}
%\section[Attribute Syntax]{Generalized Attribute Syntax}\label{attributes}
%\setcounter{table}{0}
%\setcounter{footnote}{0}
%\setcounter{lstlisting}{0}
%\input{attribute}


\newpage
%\cftaddtitleline{toc}{section}{\protect\hspace*{-5em}\framebox{\sffamily\itshape\cppxx}\vspace*{-3ex}}{}
%\section[Binary Literals]{Binary Literals\checkoddpage\Cbox}\label{binary-literals}
\section[Binary Literals]{Binary Literals}\label{binary-literals}\sectionmark{Binary Literals}
\setcounter{table}{0}
\setcounter{footnote}{0}
\setcounter{lstlisting}{0}
\input{binaryliteral}

\newpage
\section[Consecutive {\tt >}s]{Consecutive Right-Angle Brackets\sectionmark{Consecutive {\RHCode >}s}}\label{consecutive-right-angle-brackets}\sectionmark{Consecutive {\RHCode >}s}
\setcounter{table}{0}
\setcounter{footnote}{0}
\setcounter{lstlisting}{0}
\input{consecutiveangle}

\newpage
\section[{\tt decltype}]{Operator for Extracting Expression Types\sectionmark{\RHCode decltype}}\label{decltype}\sectionmark{\RHCode decltype}
\setcounter{table}{0}
\setcounter{footnote}{0}
\setcounter{lstlisting}{0}
%%%% copyedits in and proofed

The keyword \lstinline!decltype! enables the compile-time inspection of the \romeogloss{declared type} of an \romeogloss{entity} \emph{or} the type and
\romeogloss{value category} of an expression. Note that the special construct \lstinline!decltype(auto)! has a separate meaning; see \featureref{\locationf}{decltypeauto}.
%``\titleref{decltypeauto}" on page~\pageref{decltypeauto}.}

\subsection[Description]{Description}\label{description}

What results from the use of \lstinline!decltype! depends on the nature of
its operand.

\subsubsection[Use with entities]{Use with entities}\label{use-with-(typically-named)-entities}

%If an unparenthesized operand is either an \textbf{id-expression} that names an entity (or a non-type template parameter) or a \textbf{class member access expression}\glossary{class member access expression} (that identifies a class member), \texttt{decltype} yields the \emph{declared type} (the type of the \emph{entity} indicated by the operand):
If the operand is an unparenthesized \romeogloss{id-expression} or unparenthesized member access, \lstinline!decltype! yields the \emph{declared type}, meaning the type of the \emph{entity} indicated by the operand:


\begin{emcppshiddenlisting}[emcppsbatch=e1]
#include <string>  // (ù{\codeincomments{std::string}}ù)
\end{emcppshiddenlisting}
\begin{emcppslisting}[emcppsbatch=e1]
int i;                // (ù{\codeincomments{decltype(i)}}ù)   -> (ù{\codeincomments{int}}ù)
std::string s;        // (ù{\codeincomments{decltype(s)}}ù)   -> (ù{\codeincomments{std::string}}ù)
int* p;               // (ù{\codeincomments{decltype(p)}}ù)   -> (ù{\codeincomments{int}}ù)*
const int& r = *p;    // (ù{\codeincomments{decltype(r)}}ù)   -> (ù{\codeincomments{const int\&}}ù)
struct { char c; } x; // (ù{\codeincomments{decltype(x.c)}}ù) -> (ù{\codeincomments{char}}ù)
double f();           // (ù{\codeincomments{decltype(f)}}ù)   -> (ù{\codeincomments{double()}}ù)
double g(int);        // (ù{\codeincomments{decltype(g)}}ù)   -> (ù{\codeincomments{double(int)}}ù)
\end{emcppslisting}


\subsubsection[Use with expressions]{Use with expressions}\label{use-with-(unnamed)-expressions}

When \lstinline!decltype! is used with any other expression \lstinline!E! of type \lstinline!T!, including parenthesized \romeogloss{id-expression} or parenthesized member access, the result incorporates both the expression's type and
its \romeogloss{value category} (see \featureref{\locationc}{Rvalue-References}):
%\mclhead{Value category of \texttt{E}} & \mclhead{Result of
%\texttt{decltype(E)}} \\
%\emph{prvalue} & \texttt{T} \\
%\emph{lvalue} & \texttt{T\&} \\
%\emph{xvalue} & \texttt{T\&\&} \\
%%%%%%%%%%% changing this per Greg & design discussion
\begin{center}
{\small \begin{tabular}{c|c}\thickhline
\rowcolor[gray]{.9}   {\sffamily\bfseries Value category of {\ttfamily\bfseries E}}
& {\sffamily\bfseries Result of {\ttfamily\bfseries decltype(E)}} \\ \hline
{\romeovalueinside prvalue} &\lstinline!T! \\ \hline
{\romeovalueinside lvalue} & \lstinline!T\&! \\ \hline
{\romeovalueinside xvalue} & \lstinline!T\&\&! \\ \thickhline
\end{tabular}
}
\end{center}
%\noindent The three integer expressions below illustrate the various value
%categories:
%
%begin{emcppslisting}
%decltype(0)   // -> (ù{\codeincomments{int}}ù)   (*prvalue* category)
%
%int i;
%decltype((i)) // -> (ù{\codeincomments{int\&}}ù)  (*lvalue* category)
%
%int&& g();
%decltype(g()) // -> (ù{\codeincomments{int\&\&}}ù) (*xvalue* category)
%\end{emcppslisting}
%begin{emcppslisting}
%int g();
%decltype (g) g1; // id-expression,  (ù{\codeincomments{int ()();}}ù)
%decltype (g()) g2; // function call, (ù{\codeincomments{int}}ù)
%decltype ((g)) g3 = g; // id expression in parentheses, treat as lvalue,  (ù{\codeincomments{int (&)()}}ù)
%decltype ((g())) g4; // function call, parentheses have no effect, prvalue, (ù{\codeincomments{int}}ù)
%\end{emcppslisting}
%%%% new change from Josh
\noindent In general, \romeovalue{prvalues} can be passed to \lstinline!decltype! in a number of ways, including numeric literals, function calls that return by value, and explicitly created temporaries:

\begin{emcppslisting}
decltype(0)   i; // -> (ù{\codeincomments{int}}ù)
int f();
decltype(f()) j; // -> (ù{\codeincomments{int}}ù)
struct S{};
decltype(S()) k; // -> (ù{\codeincomments{S}}ù)
\end{emcppslisting}

\noindent An entity name passed to \lstinline!decltype!, as mentioned above, produces the type of the entity. If an entity name is enclosed in an additional set of parentheses, however, \lstinline!decltype! interprets its argument as an expression and its result incorporates the value category:

\begin{emcppslisting}[emcppsbatch=e2]
int i;
decltype(i)   l = i; // -> (ù{\codeincomments{int}}ù)
decltype((i)) m = i; // -> (ù{\codeincomments{int\&}}ù)
\end{emcppslisting}

\noindent Similarly, for all other \romeovalue{lvalue} expressions, the result of \lstinline!decltype! will be an \romeovalue{lvalue} reference:

\begin{emcppslisting}[emcppsbatch=e2]
int* pi = &i;
decltype(*pi) j = *pi; // -> (ù{\codeincomments{int\&}}ù)
decltype(++i) k = ++i; // -> (ù{\codeincomments{int\&}}ù)
\end{emcppslisting}

\noindent Finally, the value category of the expression will be an \romeovalue{xvalue} if it is a cast to or a function returning an \romeovalue{rvalue} reference:

\begin{emcppslisting}
int i;
decltype(static_cast<int&&>(i)) j = static_cast<int&&>(i); // -> (ù{\codeincomments{int\&\&}}ù)
int&& g();
decltype(g()) k = g();                                     // -> (ù{\codeincomments{int\&\&}}ù)
\end{emcppslisting}


\noindent Much like the \lstinline!sizeof!
operator (which is also resolved at compile time), the expression operand of \lstinline!decltype! is not evaluated:

%begin{emcppslisting}
%int i = 0;
%decltype(i++) j;  // equivalent to (ù{\codeincomments{int j;}}ù)
%assert(i == 0);   // The function (ù{\codeincomments{next}}ù) is never invoked.
%\end{emcppslisting}
\begin{emcppshiddenlisting}[emcppsbatch=e3]
#include <cassert>  // standard C (ù{\codeincomments{assert}}ù) macro
\end{emcppshiddenlisting}
\begin{emcppslisting}[emcppsbatch=e3]
void test1()
{
    int i = 0;
    decltype(i++) j;  // equivalent to (ù{\codeincomments{int j;}}ù)
    assert(i == 0);   // The expression (ù{\codeincomments{i++}}ù) was not evaluated.
}
\end{emcppslisting}
Note that the choice of using the postfix increment is significant; the prefix increment yields a different type:
\begin{emcppslisting}[emcppsbatch=e3]
void test2()
{
    int i = 0;
    int m = 1;
    decltype(++i) k = m; // equivalent to (ù{\codeincomments{int\& k = m;}}ù)
    assert(i == 0); // The expression (ù{\codeincomments{++1}}ù) is not evaluated.
}
\end{emcppslisting}

\subsection[Use Cases]{Use Cases}\label{use-cases-decltype}

\subsubsection[Avoiding unnecessary use of explicit typenames]{Avoiding unnecessary use of explicit typenames}\label{avoiding-unnecessary-use-of-explicit-typenames}

Consider two logically equivalent ways of declaring a vector of
iterators into a list of \lstinline!Widgets!:

\begin{emcppshiddenlisting}[emcppsbatch=e4]
#include <list>    // (ù{\codeincomments{std::list}}ù)   
#include <vector>  // (ù{\codeincomments{std::vector}}ù)
struct Widget {};
\end{emcppshiddenlisting}
\begin{emcppslisting}[emcppsbatch=e4]
std::list<Widget> widgets;
std::vector<std::list<Widget>::iterator> widgetIterators;
    // (1) The full type of (ù{\codeincomments{widgets}}ù) needs to be restated, and (ù{\codeincomments{iterator}}ù)
    // needs to be explicitly named.
\end{emcppslisting}
\begin{emcppshiddenlisting}[emcppsbatch=e4-1]
#include <list>    // (ù{\codeincomments{std::list}}ù)   
#include <vector>  // (ù{\codeincomments{std::vector}}ù)
struct Widget {};
\end{emcppshiddenlisting}
\begin{emcppslisting}[emcppsbatch=e4-1]
std::list<Widget> widgets;
std::vector<decltype(widgets.begin())> widgetIterators;
    // (2) Neither (ù{\codeincomments{std::list}}ù) nor (ù{\codeincomments{Widget}}ù) nor (ù{\codeincomments{iterator}}ù) need be named
    // explicitly.
\end{emcppslisting}

\noindent Notice that, when using \lstinline!decltype!, if the C++ type representing
the widget changes (e.g., from \lstinline!Widget! to, say,
\lstinline!ManagedWidget!) or the container used changes (e.g., from
\lstinline!std::list! to \lstinline!std::vector!), the declaration of
\lstinline!widgetIterators! does not necessarily need to change.

\subsubsection[Expressing type-consistency explicitly]{Expressing type-consistency explicitly}\label{expressing-type-consistency-explicitly}

In some situations, repetition of explicit type names might
inadvertently result in latent defects caused by mismatched types during
maintenance. For example, consider a \lstinline!Packet! class exposing a
\lstinline!const! member function that returns a value of type \lstinline!std::uint8_t!
representing the length of the packet's checksum:

\begin{emcppshiddenlisting}[emcppsbatch=e6]
#include <cstdint>  // (ù{\codeincomments{std::uint8\_t}}ù)
\end{emcppshiddenlisting}
\begin{emcppslisting}[emcppsbatch=e6]
class Packet
{
    // ...
    
public:
    std::uint8_t checksumLength() const;
};
\end{emcppslisting}

\noindent This unsigned 8-bit type was selected to minimize bandwidth usage
as the checksum length is sent over the network. Next, picture a loop
that computes the checksum of a \lstinline!Packet!, using the same type for its iteration variable to match the
type returned by \lstinline!Packet::checksumLength!:

\begin{emcppshiddenlisting}[emcppsbatch=e7]
#include <cstdint>  // (ù{\codeincomments{std::uint8\_t}}ù)
struct Packet
{
public:
    std::uint8_t checksumLength() const;
    std::uint8_t nthByte(std::uint8_t i) const;
};
struct Checksum {
    void appendByte(std::uint8_t b);
};
\end{emcppshiddenlisting}
\begin{emcppslisting}[emcppsbatch=e7]
void f()
{
    Checksum sum;
    Packet data;

    for (std::uint8_t i = 0; i < data.checksumLength(); ++i)  // brittle
    {
        sum.appendByte(data.nthByte(i));
    }
}
\end{emcppslisting}


Now suppose that, over time, the data transmitted by the \lstinline!Packet!
type grows to the point where the range of an \lstinline!std::uint8_t!
value might not be enough to ensure a sufficiently reliable checksum. If
the type returned by \lstinline!checksumLength()! is changed to, say,
\lstinline!std::uint16_t! without updating the type of the iteration
variable \lstinline!i! in lockstep, the loop might
silently{\cprotect\footnote{As of this writing, neither
GCC~9.3 nor Clang~10.0.0 provide
a warning (using \lstinline!-Wall!, \lstinline!-Wextra!, and
\lstinline!-Wpedantic!) for the comparison between \lstinline!std::uint8_t!
and \lstinline!std::uint16_t! --- even if (1) the value returned by
\lstinline!checksumLength! does not fit in a 8-bit integer, and (2) the
body of the function is visible to the compiler. Decorating
\lstinline!checksumLength! with \lstinline!constexpr! causes
\lstinline!clang++! to issue a warning, but this is clearly not a general
solution.}} become infinite.{\cprotect\footnote{The loop
variable is promoted to an \lstinline!unsigned! \lstinline!int! for
comparison purposes but wraps to 0 whenever its value prior to
  being incremented is 255.}}

Had \lstinline!decltype(packet.checksumLength())! been used to express the
type of \lstinline!i!, the types would have remained consistent, and the
ensuing defect would naturally have been avoided:

\begin{emcppshiddenlisting}[emcppsbatch=e7]
void g(const Packet &data)
{
\end{emcppshiddenlisting}
\begin{emcppslisting}[emcppsbatch=e7]
// ...
for (decltype(data.checksumLength()) i = 0; i < data.checksumLength(); ++i)
// ...
\end{emcppslisting}
\begin{emcppshiddenlisting}[emcppsbatch=e7]
{}
} // end g
\end{emcppshiddenlisting}    

\subsubsection[Creating an auxiliary variable of generic type]{Creating an auxiliary variable of generic type}\label{creating-an-auxiliary-variable-of-generic-type}

Consider the task of implementing a generic \lstinline!loggedSum! function
template that returns the sum of two arbitrary objects \lstinline!a! and
\lstinline!b! after logging both the operands and the result value (e.g.,
for debugging or monitoring purposes). To avoid computing the possibly
expensive sum twice, we decide to create an auxiliary function-scope
variable, \lstinline!result!. Since the type of the sum depends on both
\lstinline!a! and \lstinline!b!, we can use
\lstinline!decltype(a!~\lstinline!+!~\lstinline!b)! to infer the type for both
the trailing return type of the
function (see %Section~\ref{trailing-function-return-types},
%``\titleref{trailing-function-return-types}" on page~\pageref{trailing-function-return-types})
\featureref{\locationa}{trailing-function-return-types}) and the auxiliary variable:

\begin{emcppshiddenlisting}[emcppsbatch=e8]
struct LogTrace {};  
constexpr LogTrace LOG_TRACE; // Quick and dirty syntactic simulation of logger
template <typename T>
const LogTrace& operator<<(const LogTrace& lhs, const T&) { return lhs; }
\end{emcppshiddenlisting}
\begin{emcppslisting}[emcppsbatch=e8]
template <typename A, typename B>
auto loggedSum(const A& a, const B& b)
    -> decltype(a + b)                 // (1) exploiting trailing return types
{
    decltype(a + b) result = a + b;    // (2) auxiliary generic variable
    LOG_TRACE << a << " + " << b << " = " << result;
    return result;
}
\end{emcppslisting}
Using \lstinline!decltype(a!~\lstinline!+!~\lstinline!b)! as a return type is significantly different from relying on automatic \romeogloss{\lstinline!return!-type
deduction}; see \featureref{\locationc}{auto-feature}. Note that this particular use involves significant repetition of the expression \lstinline!a+b!.  See
%\textit{\titleref{annoyances-decltype}: \titleref{decltype-mechanical}} on page~\pageref{decltype-mechanical}
\intraref{annoyances-decltype}{decltype-mechanical}
for a discussion of ways in which this might be avoided.

\subsubsection[Determining the validity of a generic expression]{Determining the validity of a generic expression}\label{determining-the-validity-of-a-generic-expression}

In the context of generic-library development, \lstinline!decltype! can be
used in conjunction with \romeogloss{SFINAE} (``Substitution Failure Is Not An Error") to validate an expression involving a
template parameter.

For example, consider the task of writing a generic \lstinline!sortRange!
function template that, given a \romeogloss{range}, either invokes the
\lstinline!sort! member function of the argument (the one specifically
optimized for that type) if available or falls back to the more
general \lstinline!std::sort!:

\begin{emcppslisting}
template <typename Range>
void sortRange(Range& range)
{
    sortRangeImpl(range, 0);
}
\end{emcppslisting}

\noindent The client-facing \lstinline!sortRange! function (above) delegates its
behavior to an overloaded \lstinline!sortRangeImpl! function (below),
invoking the latter with the \lstinline!range! and a \romeogloss{disambiguator}
of type \lstinline!int!. The type of this additional parameter, whose value
is arbitrary, is used to give priority to the \lstinline!sort! member
function at compile time by exploiting overload resolution rules in
the presence of an implicit (\emph{standard}) conversion from
\lstinline!int! to \lstinline!long!:

\begin{emcppshiddenlisting}[emcppsbatch=e9]
#include <algorithm>  // (ù{\codeincomments{std::sort}}ù)
\end{emcppshiddenlisting}
\begin{emcppslisting}[emcppsbatch=e9]
template <typename Range>
void sortRangeImpl(Range& range,
                   long)                  // low priority: standard conversion
{
    // fallback implementation
    std::sort(std::begin(range), std::end(range));
}
\end{emcppslisting}

\noindent The fallback overload of \lstinline!sortRangeImpl! (above) will accept a
\lstinline!long! \romeogloss{disambiguator}, requiring a standard conversion from
\lstinline!int!, and will simply invoke \lstinline!std::sort!. The more
specialized overload of \lstinline!sortRangeImpl! (below) will accept an
\lstinline!int! \romeogloss{disambiguator} requiring no conversions and thus
will be a better match, provided a range-specific sort is available:

\begin{emcppslisting}
template <typename Range>
void sortRangeImpl(Range& range,
                   int,                          // high priority: exact match
                   decltype(range.sort())* = 0)  // check expression validity
{
    // optimized implementation
    range.sort();
}
\end{emcppslisting}

\noindent Note that, by exposing \lstinline!decltype(range.sort())! as part of
\lstinline!sortRangeImpl!'s declaration, the more specialized overload will
be discarded during template substitution if \lstinline!range.sort()! is
not a valid expression for the deduced \lstinline!Range!
type.{\cprotect\footnote{The technique of exposing a possibly unused
unevaluated expression --- e.g., using \lstinline!decltype! --- in a function's
declaration for the purpose of expression-validity detection prior to
template instantiation is commonly known as \romeogloss{expression SFINAE},
which is a restricted form of the more general (classical) SFINAE, and acts exclusively on
expressions visible in a function's signature rather than on frequently obscure
  template-based type computations.}}

The relative position of
\lstinline!decltype(range.sort())! in the signature of
\lstinline!sortRangeImpl! is not significant, as long as it is visible to
the compiler
during template substitution. The example shown in the main text uses a function parameter that is defaulted to
\lstinline!nullptr!. Alternatives involving a trailing return type or a
default template argument are also viable:

%\begin{emcppslisting}[basicstyle={\ttfamily\footnotesize}]
%template <typename Range>
%auto sortRangeImpl(Range& range, int) -> decltype(range.sort(), void());
%    // The comma operator is used to force the return type to (ù{\codeincomments{void}}ù),
%    // regardless of the return type of (ù{\codeincomments{range.sort()}}ù).
%
%template <typename Range, typename = decltype(std::declval<Range&>().sort()>
%auto sortRangeImpl(Range& range, int);
%    // (ù{\codeincomments{std::declval}}ù) is used to generate a reference to (ù{\codeincomments{Range}}ù) that can
%    // be used in an unevaluated expression.
%\end{emcppslisting}
\begin{emcppslisting}
#include <utility>  // (ù{\codeincomments{declval}}ù)
template <typename Range>
auto sortRangeImpl(Range& range, int) -> decltype(range.sort(), void());
    // The comma operator is used to force the return type to (ù{\codeincomments{void}}ù),
    // regardless of the return type of (ù{\codeincomments{range.sort()}}ù).

template <typename Range, typename = decltype(std::declval<Range&>().sort())>
auto sortRangeImpl(Range& range, int) -> void;
    // (ù{\codeincomments{std::declval}}ù) is used to generate a reference to (ù{\codeincomments{Range}}ù) that can
    // be used in an unevaluated expression.
\end{emcppslisting}

\noindent Putting it all together, we see that exactly two possible
outcomes exist for the original client-facing \lstinline!sortRange! function
invoked with a range argument of type \lstinline!R!:

\begin{itemize}
\item{If \lstinline!R! does have a \lstinline!sort! member function, the more specialized overload of\linebreak[4] \lstinline!sortRangeImpl! will be viable as \lstinline!range.sort()! is a well-formed expression and preferred because the \romeogloss{disambiguator} \lstinline!0! (of type \lstinline!int!) requires no conversion.}
\item{Otherwise, the more specialized overload will be discarded during template substitution as \lstinline!range.sort()! is not a well-formed expression, and the only remaining more general \lstinline!sortRangeImpl! overload will be chosen instead.}
\end{itemize}

\subsection[Potential Pitfalls]{Potential Pitfalls}\label{potential-pitfalls}

Perhaps surprisingly, \lstinline!decltype(x)! and \lstinline!decltype((x))!
will sometimes yield different results for the same expression
\lstinline!x!:

\begin{emcppslisting}[emcppsbatch=e10]
int i = 0; // (ù{\codeincomments{decltype(i)}}ù) yields (ù{\codeincomments{int}}ù).
           // (ù{\codeincomments{decltype((i))}}ù) yields (ù{\codeincomments{int\&}}ù).
\end{emcppslisting}

\noindent In the case where the unparenthesized operand is an entity having a
declared type \lstinline!T! and the parenthesized operand is an expression
whose value category is represented (by \lstinline!decltype!) as the same
type \lstinline!T!, the results will coincidentally be the same:

\begin{emcppslisting}[emcppsbatch=e10]
int& ref = i;  // (ù{\codeincomments{decltype(ref)}}ù) yields (ù{\codeincomments{int\&}}ù).
               // (ù{\codeincomments{decltype((ref))}}ù) yields (ù{\codeincomments{int\&}}ù).
\end{emcppslisting}

\noindent Wrapping its operand with parentheses ensures \lstinline!decltype! yields
the \romeogloss{value category} of a given expression. This technique can be
useful in the context of metaprogramming --- particularly in the case of
\romeogloss{value category} propagation.

\subsection[Annoyances]{Annoyances}\label{annoyances-decltype}

\subsubsection{Mechanical repetition of expressions might be required}\label{decltype-mechanical}

As mentioned in
%\textit{\titleref{use-cases-decltype}: \titleref{creating-an-auxiliary-variable-of-generic-type}} on page~\pageref{creating-an-auxiliary-variable-of-generic-type},
\intraref{use-cases-decltype}{creating-an-auxiliary-variable-of-generic-type},
using \lstinline!decltype! to capture a value of an expression that is about to be used, or for the return value of an expression, can often lead to repeating the same expression in multiple places (three distinct ones in the earlier example).

An alternate solution to this problem is to capture the result of the \lstinline!decltype! expression in a \lstinline!typedef!, \lstinline!using! type alias, or as a defaulted template parameter --- but that runs into the problem that it can be used only once the expression is valid.  A defaulted \lstinline!template! parameter cannot reference parameter names as it is written before them, and a type alias cannot be introduced prior to the return type being needed.  A solution to this problem lies in using standard library function \lstinline!std::declval! to create expressions of the appropriate type without needing to reference the actual function parameters by name:

\begin{emcppshiddenlisting}[emcppsbatch=e11]
#include <utility>  // (ù{\codeincomments{std::declval}}ù)
struct LogTrace {};  
constexpr LogTrace LOG_TRACE; // Quick and dirty syntactic simulation of logger
template <typename T>
const LogTrace& operator<<(const LogTrace& lhs, const T&) { return lhs; }
\end{emcppshiddenlisting}
\begin{emcppslisting}[emcppsbatch=e11]
template <typename A, typename B,
          typename Result = decltype(std::declval<const A&>() +
                                     std::declval<const B&>())>
Result loggedSum(const A& a, const B& b)
{
    Result result = a + b;  // no duplication of the (ù{\codeincomments{decltype}}ù) expression
    LOG_TRACE << a << " + " << b << "=" << result;
    return result;
}
\end{emcppslisting}

Here, \lstinline!std::declval!, a function that cannot be executed at runtime and is only appropriate for use in \romeogloss{unevaluated contexts}, produces an expression of the specified type.  When mixed with \lstinline!decltype!, this lets us determine the result types for expressions without needing to construct (or even being able to construct) objects of the needed types.

\subsection[See Also]{See Also}\label{see-also}

\begin{itemize}
\item{%``\titleref{Rvalue-References}" on page~\pageref{Rvalue-References} ---
\seealsoref{Rvalue-References}{\locationc}The \lstinline!decltype! operator yields precise information on whether an expression is an \romeovalue{lvalue} or \romeovalue{rvalue}.}
\item{%``\titleref{alias-declarations-and-alias-templates}" on page~\pageref{alias-declarations-and-alias-templates} ---
\seealsoref{alias-declarations-and-alias-templates}{\locationa}Oftentimes, it is useful to give a name to the type yielded by \lstinline!decltype!, which is done with a \lstinline!using! alias.}
\item{%``\titleref{auto-feature}" on page~\pageref{auto-feature} ---
\seealsoref{auto-feature}{\locationc}The type computed by \lstinline!decltype! is similar to, but distinct from, the type deduction used by \lstinline!auto!.}
\item{%``\titleref{decltypeauto}" on page~\pageref{decltypeauto} ---
\seealsoref{decltypeauto}{\locationf}\lstinline!decltype! type computation rules can be useful in conjunction with an \lstinline!auto! variable.}
\end{itemize}

\subsection[Further Reading]{Further Reading}\label{further-reading}






\newpage
\section[Deleted Functions]{Using {\SecCode =delete} for Arbitrary Functions}\label{deleted-functions}
\setcounter{table}{0}
\setcounter{footnote}{0}
\setcounter{lstlisting}{0}


Use of \texttt{=}~\texttt{delete} in a function's (first) declaration
forces a compilation error upon any attempt to use or access it.

\subsection[Description]{Description}\label{description}

Declaring a particular function (or function overload) to result in a
fatal diagnostic upon invocation can be useful --- e.g., to suppress the
generation of a \emph{special function} or to limit the types of
arguments a particular function is able to accept. In such cases,
\texttt{=}~\texttt{delete;} can be used in place of the body of any
function (on first declaration only) to force a compile-time error if
any attempt is made to invoke it or take its address.

\begin{lstlisting}[language=C++]
void g(double) { }
void g(int) = delete;

void f()
{
    g(3.14);  // OK, (ù{\codeincomments{f(double)}}ù) is invoked.
    g(0);     // Error: (ù{\codeincomments{f(int)}}ù) is deleted.
}
\end{lstlisting}
    
\noindent Notice that deleted functions participate in \emph{overload resolution}
and produce a compile-time error when selected as the best candidate.

\subsection[Use Cases]{Use Cases}\label{use-cases}

\subsubsection[Suppressing special member function generation]{Suppressing special member function generation}\label{suppressing-special-member-function-generation}

When instantiating an object of user-defined type, \textbf{special
member functions}\glossary{special
member functions} that have not been declared explicitly are
often{\cprotect\footnote{The generation of individual special member
functions can be affected by the existence of other user-defined special
member functions or by limitations imposed by the specific types of
any data members or base types. For more information, see Section~\ref{Defaulted-Special-Member-Functions}, ``\titleref{Defaulted-Special-Member-Functions}."}} generated automatically by the compiler.
For certain kinds of types, the notion of \textbf{copy semantics}\glossary{copy semantics}
(including \textbf{move semantics}\glossary{move semantics}{\cprotect\footnote{The two
\textbf{special member functions}\glossary{special member functions} controlling \emph{move} operations
(introduced in C++11) are sometimes implemented as effective
optimizations of copy operations and (rarely) with copy operations
  explicitly deleted; see Section~\ref{Rvalue-References}, ``\titleref{Rvalue-References}."}}) is not meaningful and
hence permitting the generation of copy operations is contraindicated.

Consider a class, \texttt{FileHandle}, that uses the \textbf{RAII}\glossary{RAII} idiom
to safely acquire and release an I/O stream. As \emph{copy semantics}
are typically not meaningful for such resources, we will want to
suppress generation of both the \emph{copy constructor} and \emph{copy
assignment operator}. Prior to C++11, there was no direct way to express
suppression of \emph{special functions} in C++. The commonly
recommended workaround was to declare the two methods \texttt{private}
and leave them unimplemented, typically resulting in a compile-time (or
link-time) error when accessed{\cprotect\footnote{Leaving unimplemented a special
member function that is declared to be private ensures
that there will be at least a link-time error in case that function is
inadvertently accessed from within the implementation of the class
  itself.}}:

\begin{lstlisting}[language=C++]
class FileHandle
{
private:
    // ...

    FileHandle(const FileHandle&);             // not implemented
    FileHandle& operator=(const FileHandle&);  // not implemented

public:
    explicit FileHandle(FILE* filePtr);
    ~FileHandle();

    // ...
};
\end{lstlisting}
    
\noindent With the \texttt{=}~\texttt{delete} syntax, we are able to (1)
explicitly express our intention to make these special member
functions unavailable, (2) do so directly in the \texttt{public} region
of the class, and (3) enable more precise compiler diagnostics:

\begin{lstlisting}[language=C++]
class FileHandle
{
private:
    // ...

public:
    explicit FileHandle(FILE* filePtr);
    ~FileHandle();

    FileHandle(const FileHandle&) = delete;             // make unavailable
    FileHandle& operator=(const FileHandle&) = delete;  // make unavailable

    // ...
};
\end{lstlisting}
    

\subsubsection[Preventing a particular implicit conversion]{Preventing a particular implicit conversion}\label{preventing-a-particular-implicit-conversion}

Certain functions --- especially those that take a \texttt{char} as an
argument --- are prone to inadvertent misuse. As a truly classic
example, consider the C library function \texttt{memset}, which may be used
to write the character \texttt{*} five times in a row, starting at a
specified memory address, \texttt{buf}:

\begin{lstlisting}[language=C++]
#include <cstring>
#include <cstdio>

void f()
{
    char buf[] = "Hello World!";
    memset(buf, 5, '*');  // undefined behavior
    puts(buf);            // expected output: "***** World!"
}
\end{lstlisting}
    
\noindent Sadly, inadvertently reversing the final two arguments is a commonly
recurring error, and the C language provides no help. In C++, we
can target such observed misuse using an extra deleted overload:

\begin{lstlisting}[language=C++]
#include <cstring>  // memset()
void* memset(void* str, int ch, size_t n);      // standard library function
void* memset(void* str, int n, char) = delete;  // defensive against misuse
\end{lstlisting}
    
\noindent Pernicious user errors can now be reported during compilation:

\begin{lstlisting}[language=C++]
// ...
memset(buf, 5, '*');  // Error: (ù{\codeincomments{memset(void\*, int, char)}}ù) is deleted.
// ...
\end{lstlisting}
    

\subsubsection[Preventing all implicit conversions]{Preventing all implicit conversions}\label{preventing-all-implicit-conversions}

The \texttt{ByteStream::send} member function below is designed to work
with 8-bit unsigned integers only. Providing a deleted overload
accepting an \texttt{int} forces a caller to ensure that the argument is
always of the appropriate type:

\begin{lstlisting}[language=C++]
class ByteStream
{
public:
    void send(unsigned char byte) { /* ... */ }
    void send(int) = delete;

    // ...
};

void f()
{
    ByteStream stream;
    stream.send(0);   // Error: (ù{\codeincomments{send(int)}}ù) is deleted.     (1)
    stream.send('a'); // Error: (ù{\codeincomments{send(int)}}ù) is deleted.     (2)
    stream.send(0L);  // Error: ambiguous                 (3)
    stream.send(0U);  // Error: ambiguous                 (4)
    stream.send(0.0); // Error: ambiguous                 (5)
    stream.send(
        static_cast<unsigned char>(100));  // OK          (6)
}
\end{lstlisting}
    
\noindent Invoking \texttt{send} with an \texttt{int} (noted with (1) in the code above) or any integral type
(other than \texttt{unsigned}~\texttt{char}{\cprotect\footnote{Note that
implicitly converting from \texttt{unsigned}~\texttt{char} to either a
\texttt{long} or \texttt{unsigned} integer involves a \textbf{standard
conversion}\glossary{standard
conversion} (not just an \textbf{integral promotion}\glossary{integral promotion}), the same as
  converting to a \texttt{double}.}}) that promotes to \texttt{int} (2)
will map exclusively to the deleted \texttt{send(int)} overload; all
other integral (3 \& 4) and floating-point types (5) are convertible to
both (via a \textbf{standard conversion}\glossary{standard conversion}) and hence will be ambiguous.
An explicit cast to \texttt{unsigned}~\texttt{char} (6) can always be
pressed into service if needed.

\subsubsection[Hiding a structural (nonpolymorphic) base class's member function]{Hiding a structural (nonpolymorphic) base class’s member function}\label{hiding-a-structural-(nonpolymorphic)-base-class's-member-function}

Best practices notwithstanding,{\cprotect\footnote{By publicly deriving
from a concrete class, we do not hide the underlying capabilities,
which can easily be accessed (perhaps accidentally) via assignment to
a pointer or reference to a base class (no casting required). What's
more, inadvertently passing such a class to a function taking the base
class by value will result in \emph{slicing}, which can be especially
problematic when the derived class holds data. Finally, if the derived
class purports to maintain \emph{class invariants} that the base class
does not preserve, this design technique is beyond dubious; a more
robust approach would be to use layering or at least private
inheritance. For more on improving compositional designs at scale, see
  \textbf{{lakos20}}, sections 3.5.10.5 and 3.7.3, pp.~687--703 and
  726--727, respectively.}} it can be cost-effective (in the short
term) to provide an elided ``view'' on a concrete class for (trusted)
clients. Imagine a class \texttt{AudioStream} designed to play sounds
and music that --- in addition to providing basic ``play'' and
``rewind'' operations --- sports a large, robust interface:

\begin{lstlisting}[language=C++]
struct AudioStream
{
    void play();
    void rewind();
    // ...
    // ... (large, robust interface)
    // ...
};
\end{lstlisting}
    

Now suppose that, on short notice, we need to whip up a very similar
class,\linebreak[4] \texttt{ForwardAudioStream}, to use with audio samples that
cannot be rewound (e.g., coming directly from a live feed). Realizing
that we can readily reuse most of \texttt{AudioStream}'s interface, we
pragmatically decide to prototype the new class simply by exploiting
public \textbf{structural inheritance}\glossary{structural inheritance} and then deleting just the lone
unwanted \texttt{rewind} member function:

\begin{lstlisting}[language=C++]
struct ForwardAudioStream : AudioStream
{
    void rewind() = delete; // make just this one function unavailable
};

void f()
{
    ForwardAudioStream stream = FMRadio::getStream();
    stream.play();   // fine
    stream.rewind(); // Error: (ù{\codeincomments{rewind()}}ù) is deleted.
}
\end{lstlisting}
    
\noindent If the need for a \texttt{ForwardAudioStream} type persists, we can always
consider reimplementing it more carefully later.{\cprotect\footnote{\textbf{{lakos20}}, sections 3.5.10.5 and 3.7.3, pp.~687--703 and 726--727}}

\subsection[Potential Pitfalls]{Potential Pitfalls}\label{potential-pitfalls}

None so far

\subsection[Annoyances]{Annoyances}\label{annoyances}

None so far

\subsection[See Also]{See Also}\label{see-also}

\begin{itemize}
\item{Section~\ref{Defaulted-Special-Member-Functions}, ``\titleref{Defaulted-Special-Member-Functions}" — Companion safe C++11 feature that enables \emph{defaulting} (as opposed to \emph{deleting}) special member functions}
\item{Section~\ref{Rvalue-References}, ``\titleref{Rvalue-References}" — Conditionally safe C++11 feature that introduces the two \emph{move} variants to \emph{copy} special member functions}
\end{itemize}

\subsection[Further Reading]{Further Reading}\label{further-reading}

None so far




\newpage
\section[{\tt noreturn}]{The {\SecCode [[noreturn]]} Attribute\sectionmark{\RHCode noreturn}}\label{the-noreturn-attribute}\sectionmark{\RHCode noreturn}
\setcounter{table}{0}
\setcounter{footnote}{0}
\setcounter{lstlisting}{0}



The \texttt{[[noreturn]]} attribute promises that the function to which
it pertains never \nobreak{returns}.

\subsection[Description]{Description}\label{description}

The presence of the standard \texttt{[[noreturn]]} attribute as part of
a function declaration informs both the compiler and human readers that
such a function never returns control flow to the caller:

\begin{lstlisting}[language=C++]
[[noreturn]] void f()
{
    throw;
}
\end{lstlisting}
    
\noindent The \texttt{[[noreturn]]} attribute is not part of a function's type
and is also, therefore, not part of the type of a function pointer. Applying \texttt{[[noreturn]]} to a function
pointer is not an error, though doing so has no actual effect in standard C++; see
{\it\titleref{noreturn-potential-pitfalls}: \titleref{misuse-of-[[noreturn]]-on-function-pointers}} on page~\pageref{misuse-of-[[noreturn]]-on-function-pointers}. Use on a pointer might have
benefits for external tooling, code expressiveness, and future language
evolution:

\begin{lstlisting}[language=C++]
void (*fp [[noreturn]])() = f;
\end{lstlisting}
    

\subsection[Use Cases]{Use Cases}\label{use-cases}

\subsubsection[Better compiler diagnostics]{Better compiler diagnostics}\label{better-compiler-diagnostics}

Consider the task of creating an assertion handler that, when invoked,
always aborts execution of the program after printing some useful
information about the source of the assertion. Since this specific
handler will never return, it is a viable candidate for
\texttt{[[noreturn]]}:

\begin{lstlisting}[language=C++]
[[noreturn]] void abortingAssertionHandler(const char* filename, int line)
{
    LOG_ERROR << "Assertion fired at " << filename << ':' << line;
    std::abort();
}
\end{lstlisting}
    

\noindent The additional information provided by the attribute will allow a
compiler to warn if it determines that a code path in the
function would allow it to return:

\begin{lstlisting}[language=C++]
[[noreturn]] void abortingAssertionHandler(const char* filename, int line)
{
    if (filename)  // just being safe, but see "Further Reading," below
    {
        LOG_ERROR << "Assertion fired at " << filename << ':' << line;
        std::abort();
    }
}  // compile-time warning made possible
\end{lstlisting}
    

\noindent This information can also be used to warn in case unreachable code is
present after\linebreak[4] \texttt{abortingAssertionHandler} is invoked:

\begin{lstlisting}[language=C++]
int main()
{
    // ...
    abortingAssertionHandler("main.cpp", __LINE__);
    std::cout << "We got here.\n";  // compile-time warning made possible
    // ...
}
\end{lstlisting}
    

\noindent Note that this warning is made possible by decorating just the
declaration of the handler function --- i.e., even if the definition of
the function is not visible in the current translation unit.

\subsubsection[Improved runtime performance]{Improved runtime performance}\label{improved-runtime-performance}

If the compiler knows that it is going to invoke a function
that is guaranteed not to return, the compiler is within its rights to optimize
that function by removing what it can now determine to be dead code. As
an example, consider a utility component, \texttt{util}, that defines a
function, \texttt{throwBadAlloc}, that is used to \textbf{insulate} the
throwing of an \texttt{std::bad\_alloc} exception in what would
otherwise be template code fully exposed to clients:

\begin{lstlisting}[language=C++]
// util.h:
[[noreturn]] void throwBadAlloc();

// util.cpp:
#include <util.h>  // (ù{\codeincomments{[[noreturn]] void throwBadAlloc()}}ù)

void throwBadAlloc()  // This redeclaration is also (ù{\codeincomments{[[noreturn]]}}ù).
{
    throw std::bad_alloc();
}
\end{lstlisting}
    

\noindent Irrespective of whether the client compiler warns, the compiler is within its
rights to elide code that is rendered unreachable by the call to the
\texttt{throwBadAlloc} function due to the function being decorated with the
\texttt{[[noreturn]]} attribute on its declaration:

\begin{lstlisting}[language=C++]
#include <util.h>  // [[noreturn]] void throwBadAlloc()

void client()
{
    // ...
    throwBadAlloc();
    // ... (Everything below here can be optimized away.)
}
\end{lstlisting}
    

\noindent Notice that even though \texttt{[[noreturn]]} appeared only on the first
declaration (in the \texttt{util.h} header), the \texttt{[[noreturn]]}
attribute carries over to the redeclaration used in the\linebreak[4] \texttt{throwBadAlloc}
function's definition because the header was included in the
corresponding \texttt{.cpp} file.

\subsection[Potential Pitfalls]{Potential Pitfalls}\label{noreturn-potential-pitfalls}

\subsubsection[\tt{[[noreturn]]} can inadvertently break an otherwise working program]{{\ParaCode [[noreturn]]} can inadvertently break an otherwise working program}\label{[[noreturn]]-can-inadvertently-break-an-otherwise-working-program}

Unlike many attributes, use of \texttt{[[noreturn]]} \emph{can} alter
the semantics of a well-formed program, potentially introducing a
runtime defect and/or making the program ill-formed. If a function that
can potentially return is decorated with \texttt{[[noreturn]]} and then,
in the course of executing a program, it ever does return, that behavior
is \textbf{undefined}.

Consider a \texttt{printAndExit} function whose role is to print a fatal
error message before aborting the program:

\begin{lstlisting}[language=C++]
[[noreturn]] void printAndExit()
{
    std::cout << "Fatal error. Exiting the program.\n";
    assert(false);
}
\end{lstlisting}
    
\noindent The programmer chose to (sloppily) implement termination by using an
assertion, which would not be incorporated into a program compiled with
the preprocessor definition \texttt{NDEBUG} active, and thus
\texttt{printAndExit} would (in that build mode) return normally. If the
compiler of the client is informed that function will not return, the compiler is
free to optimize accordingly. If the function then does return, any
number of hard-to-diagnose defects (e.g., due to incorrectly elided
code) might materialize as a consequence of the ensuing
\textbf{undefined behavior}. Furthermore, within a program, if a function is declared \texttt{[[noreturn]]} in some
translation units but not in others, that program is (inherently)
\textbf{ill-formed, no diagnostic required}.

\subsubsection[Misuse of {\tt[[noreturn]]} on function pointers]{Misuse of {\ParaCode [[noreturn]]} on function pointers}\label{misuse-of-[[noreturn]]-on-function-pointers}

Although the \texttt{[[noreturn]]} attribute is permitted to appertain
to a function pointer (syntactically) for the benefit of external tools,
it has no effect in standard C++; fortunately, most compilers will warn:

\begin{lstlisting}[language=C++]
void (*fp [[noreturn]])();  // not supported by standard C++ (will likely warn)
\end{lstlisting}
    
\noindent What's more, assigning (the address of) a function
that is not decorated with \texttt{[[noreturn]]} to an otherwise
suitable function pointer that is so decorated is perfectly fine:

\begin{lstlisting}[language=C++]
void f() { return; };  // function that always returns

void g()
{
    fp = f;  // "OK" -- that (ù{\codeincomments{fp}}ù) is (ù{\codeincomments{[[noreturn]]}}ù) is silently ignored
}
\end{lstlisting}
    
\noindent Any reliance on \texttt{[[noreturn]]} to have any effect in standard C++
when applied to other than a function's declaration is misguided.

\subsection[Annoyances]{Annoyances}\label{annoyances}

None so far

\subsection[See Also]{See Also}\label{see-also}

\begin{itemize}
\item{Section~\ref{attributes}, ``\titleref{attributes}" — To learn more about allowed attribute placement in general}
\end{itemize}

\subsection[Further Reading]{Further Reading}\label{further-reading}

None so far





\newpage
\section[{\tt nullptr}]{The Null-Pointer-Literal Keyword\sectionmark{\RHCode nullptr}}\label{null-pointer-literal-(nullptr)}\sectionmark{\RHCode nullptr}
\setcounter{table}{0}
\setcounter{footnote}{0}
\setcounter{lstlisting}{0}


The keyword \texttt{nullptr} (unambiguously) denotes the
null-pointer-value literal.

\subsection[Description]{Description}\label{description}

The \texttt{nullptr} keyword is a \textbf{prvalue} (pure rvalue) of type
\texttt{std::nullptr\_t} representing the (implementation-defined)
bit pattern corresponding to a \textbf{null address} on the host platform;
\texttt{nullptr} and other values of type \texttt{std::nullptr\_t}
(along with the integer literal \texttt{0} and the macro \texttt{NULL})
can be converted implicitly to any pointer (or pointer-to-member) type:

\begin{lstlisting}[language=C++]
int data;  // non-member data

int *pi0 = &data;    // Initialize with non-null address.
int *pi1 = nullptr;  // Initialize with null address.
int *pi2 = NULL;     //  "          "    "    "
int *pi3 = 0;        //  "          "    "    "

double f(int x);  // non-member function

double (*pf0)(int) = &f;       // Initialize with non-null address.
double (*pf1)(int) = nullptr;  // Initialize with null address.

struct S
{
    short d_data;    // member data
    float g(int y);  // member function
};

short S::*pmd0 = &S::d_data;  // Initialize with non-null address.
short S::*pmd1 = nullptr;     // Initialize with null address.

float (S::*pmf0)(int) = &S::g;    // Initialize with non-null address.
float (S::*pmf1)(int) = nullptr;  // Initialize with null address.
\end{lstlisting}
    

\noindent Because \texttt{std::nullptr\_t} is its own distinct type, overloading on
it is possible:

\begin{lstlisting}[language=C++]
#include <cstddef>  // (ù{\codeincomments{std::nullptr\_t}}ù)

void g(void*);           // (1)
void g(int);             // (2)
void g(std::nullptr_t);  // (3)

void f()
{
    g("hello");  // OK --- (1) (ù{\codeincomments{void g(void}}ù)*(ù{\codeincomments{)}}ù)
    g(0);        // OK --- (2) (ù{\codeincomments{void g(int)}}ù)
    g(nullptr);  // OK --- (3) (ù{\codeincomments{void g(std::nullptr\_t)}}ù)
    g(NULL);     // Error: ambiguous --- (1), (2), or (3)
}
\end{lstlisting}
    

\subsection[Use Cases]{Use Cases}\label{use-cases}

\subsubsection[Improve type safety]{Improve type safety}\label{improve-type-safety}

In pre-C++11 code bases, use of the \texttt{NULL}
macro{\cprotect\footnote{In the C Standard, the macro \texttt{NULL} is
defined as an \textbf{implementation-defined} integral or \texttt{void*}
constant. Unlike C, C++ forbids conversions from \texttt{void*} to
arbitrary pointer types and instead, prior to C++11, defined
\texttt{NULL} as an \emph{``integral constant expression rvalue of
integer type that evaluates to zero''}; any integer literal (e.g.,
\texttt{0}, \texttt{0L}, \texttt{0U}, \texttt{0LLU}) satisfies this
  criterion.}} was a common way of indicating (mostly to the human
reader) that the literal value it conveys is intended specifically to
represent a \emph{null address} rather than the literal \texttt{int}
value \texttt{0}. From a type-safety perspective, its
implementation-defined (typically integral{\cprotect\footnote{As of
C++11, the definition of \texttt{NULL} has been expanded to --- in
theory --- permit \texttt{nullptr} as a conforming definition; as of
this writing, however, no major compiler vendors do so. Both GCC and
Clang default to \texttt{0L} (\texttt{long}~\texttt{int}) while MSVC
defaults to \texttt{0} (\texttt{int}). Such definitions are unlikely
to change since existing code could cease to compile or (possibly
  silently) present altered run-time behavior.}}) definition, however,
makes the use of \texttt{NULL} only marginally better suited than a raw
literal \texttt{0} to represent a null pointer.

As just one specific illustration of the added type safety provided by
\texttt{nullptr}, imagine that you work for a large software company
that has historically required, as one of its coding standards, that
values returned via output parameters (as opposed to a \texttt{return}
statement) are always returned via pointer to a modifiable
object.{\cprotect\footnote{Functions that return via argument typically
do so to reserve the function's return value to communicate status.
  See \textbf{{lakos96}}, section~9.1.11, pp.~621--628, specifically the
  \emph{Guideline} at the bottom of p.~621: ``Be consistent about
  returning values through arguments (e.g., avoid declaring
  non-\texttt{const} reference parameters)."}} In the illustrative function below, the output parameter's local pointer
variable is ``zeroed'' (shown here in three different ways) to indicate
(and ensure) that nothing more is to be written:

\begin{lstlisting}[language=C++]
int illustrativeFunction(int* x)   // pointer to modifiable integer
{
    // ...
    if (/*...*/)
    {
        x = 0;       // (1) Set pointer (ù{\codeincomments{x}}ù) to null address.
        x = NULL;    // (2) Set pointer (ù{\codeincomments{x}}ù) to null address.
        x = nullptr; // (3) Set pointer (ù{\codeincomments{x}}ù) to null address.
    }
    // ...
    return 0;    // success
}
\end{lstlisting}
    

Now suppose that the function signature is changed (e.g., due to a
change in coding standards in the organization) to accept a reference
instead of a pointer:

\begin{lstlisting}[language=C++]
int illustrativeFunction(int& x)  // reference to modifiable integer
{
    // ...
    if (/*...*/)
    {
        x = 0;       // (1) always compiles; makes what (ù{\codeincomments{x}}ù) refers to 0
        x = NULL;    // (2) implementation-defined (might warn)
        x = nullptr; // (3) always a compile-time error
    }
    // ...
    return 0;    // SUCCESS
}
\end{lstlisting}
    

As the example above demonstrates, how we represent the notion of a
\textbf{null address} matters:
\begin{enumerate}
\item{\texttt{0} — Portable across all implementations but minimal type safety.}
\item{\texttt{NULL} — Implemented as a macro; added type safety (if any) is platform specific.}
\item{\texttt{nullptr} — Portable across all implementations and fully type-safe.}
\end{enumerate}
Use of \texttt{nullptr} instead of \texttt{0} or \texttt{NULL} to denote
a null address maximizes type safety and readability, while avoiding
both macros and implementation-defined behavior.

\subsubsection[Disambiguation of \tt{(int) 0} versus \tt{(T*) 0} during overload resolution]{Disambiguation of {\ParaCode (int) 0} versus {\ParaCode (T*) 0} during overload resolution}\label{disambiguation-of-(int)-0-versus-(t*)-0-during-overload-resolution}

The platform-dependent nature of \texttt{NULL} presents additional
challenges when used to call a function whose overloads differ only in
accepting a pointer or an integral type as the same positional argument
(which might be the case, e.g., in a poorly designed third-party
library):

\begin{lstlisting}[language=C++]
void uglyLibraryFunction(int* p);  // (1)
void uglyLibraryFunction(int  i);  // (2)
\end{lstlisting}
    

\noindent Calling this function with the literal \texttt{0} will always invoke
overload (2), but that might not always be what casual clients expect:

\begin{lstlisting}[language=C++]
void f()
{
    uglyLibraryFunction(0);         // unambiguously invokes (2)
    uglyLibraryFunction((int*) 0);  // unambiguously invokes (1)
    uglyLibraryFunction(nullptr);   // unambiguously invokes (1)
    uglyLibraryFunction(NULL);      // anything! (platform-defined)
    uglyLibraryFunction(0L);        // always ambiguous
    uglyLibraryFunction(0U);        // always ambiguous
}
\end{lstlisting}
    

\noindent\texttt{nullptr} is especially useful when such problematic overloads
are unavoidable because it obviates explicit
casts.{\cprotect\footnote{N.B., Explicitly casting \texttt{0} to an
appropriately typed pointer (other than \texttt{void*}) was at one
  time considered by some to be a best practice, especially in C.}}

\subsubsection[Overloading for a literal null pointer]{Overloading for a literal null pointer}\label{overloading-for-a-literal-null-pointer}

Being a distinct type, \texttt{std::nullptr\_t} can itself participate
in an overload set:

\begin{lstlisting}[language=C++]
void f(int* v);          // (1)
void f(std::nullptr_t);  // (2)

void g()
{
    int* ptr = nullptr;
    f(ptr);      // unambiguously invokes (1)
    f(nullptr);  // unambiguously invokes (2)
}
\end{lstlisting}
    

\noindent Given the relative ease with which a \texttt{nullptr} (above) can be
converted to a typed pointer having the same null-address value, such
overloads are, however, dubious when used to control essential behavior.
Nonetheless, one can envision such use to, say, aid in compile-time
diagnostics when passing a \textbf{null address} would otherwise result in
a runtime error\footnote{see also Section~\ref{deleted-functions}, ``\titleref{deleted-functions}"}:

\begin{lstlisting}[language=C++]
std::size_t strlen(const char* s);
    // The behavior is undefined unless (ù{\codeincomments{s}}ù) is null-terminated.

std::size_t strlen(std::nullptr_t);
    // declared but not defined 
\end{lstlisting}
    

\noindent Another arguably safe use of such an overload for a \texttt{nullptr} is
to avoid a null-pointer check. However, for cases where the client knows the address is null at compile time, better ways typically exist for avoiding the (often
insignificant) overhead to test for a null pointer at runtime.

\subsection[Potential Pitfalls]{Potential Pitfalls}\label{potential-pitfalls}

None so far

\subsection[Annoyances]{Annoyances}\label{annoyances}

None so far

\subsection[See Also]{See Also}\label{see-also}

None so far

\subsection[Further Reading]{Further Reading}\label{further-reading}

None so far




\newpage
\section[{\tt override}]{The {\SecCode override} Member-Function Specifier\sectionmark{\RHCode override}}\label{override}\sectionmark{\RHCode override}
\setcounter{table}{0}
\setcounter{footnote}{0}
\setcounter{lstlisting}{0}
\input{override}

\newpage
\section[{\tt static\_assert}]{Compile-Time Assertions\sectionmark{\RHCode static\_assert}}\label{compile-time-assertions-(static_assert)}\sectionmark{\RHCode static\_assert}
\setcounter{table}{0}
\setcounter{footnote}{0}
\setcounter{lstlisting}{0}


The \texttt{static\_assert} keyword allows programmers to intentionally terminate compilation whenever a given compile-time predicate
 evaluates to \texttt{false}.

\subsection[Description]{Description}\label{description}

Assumptions are
inherent in every program, whether we explicitly document them or not. A common way of validating certain
assumptions at runtime is to use the classic \texttt{assert} macro found
in \texttt{<cassert>}. Such runtime assertions are not always ideal
because (1) the program must already be built and running for them to
even have a chance of being triggered and (2) executing a
\textbf{redundant check}\glossary{redundant check} at runtime typically{\cprotect\footnote{It is
not unheard of for a program having assertions to run faster with them
enabled than disabled --- e.g., when asserting that a pointer is not
null, thereby enabling the optimizer to elide all code branches that
  can be reached only if that pointer were null.}} results in a slower
program. Being able to validate an assertion at compile time avoids
several drawbacks:

\begin{enumerate}
\item{Validation occurs at compile time within a single translation unit, and therefore doesn’t need to wait until a complete program is linked and executed.}
\item{Compile-time assertions can exist in many more places than runtime assertions and are unrelated to program control flow.}
\item{No runtime code will be generated due to a \texttt{static\_assert}, so program performance will not be impacted.}
\end{enumerate}

\subsubsection[Syntax and semantics]{Syntax and semantics}\label{syntax-and-semantics}

We can use \emph{static assertion declarations} to conditionally trigger
controlled compilation failures depending on the truthiness of a
\textbf{constant expression}\glossary{constant expression}. Such declarations are introduced by the
\texttt{static\_assert} keyword, followed by a parenthesized list
consisting of (1) a constant Boolean expression and (2) a mandatory (see
{\it\titleref{static-annoyances}: \titleref{mandatory-string-literal}} on page~\pageref{mandatory-string-literal}) \textbf{string literal}\glossary{string literal}, which will be
part of the compiler diagnostics if the compiler determines that the
assertion fails to hold:

\begin{lstlisting}[language=C++]
static_assert(true, "Never fires.");
static_assert(false, "Always fires.");
\end{lstlisting}
    
\noindent Static assertions can be placed anywhere in the scope of a namespace,
block, or class:

\begin{lstlisting}[language=C++]
static_assert(1 + 1 == 2, "Never fires.");  // (global) namespace scope

template <typename T>
struct S
{
    void f0()
    {
        static_assert(1 + 1 == 3, "Always fires.");  // block scope
    }

    static_assert(!Predicate<T>::value, "Might fire.");  // class scope
};
\end{lstlisting}
    
\noindent Providing a non-constant expression to a \texttt{static\_assert} is
itself a compile-time error:

\begin{lstlisting}[language=C++]
extern bool x;
static_assert(x, "Nice try.");  // Error: (ù{\codeincomments{x}}ù) is not a compile-time constant.
\end{lstlisting}
    

\subsubsection[Evaluation of static assertions in templates]{Evaluation of static assertions in templates}\label{evaluation-of-static-assertions-in-templates}

The C++ Standard does not explicitly specify at precisely what point (during
the compilation process) static assertion declarations are
evaluated.{\cprotect\footnote{By ``evaluated'' here, we mean that the
  asserted expression is processed and its semantic truth\linebreak[4] determined.}}
In particular, when used within the body of a template, a
\texttt{static\_assert} declaration might not be evaluated until
\textbf{template instantiation time}\glossary{template instantiation time}. In practice, however, a
\texttt{static\_assert} that does not depend on any template parameters
is essentially always{\cprotect\footnote{E.g.,
GCC~10.1, Clang~10.0, and MSVC~19.24}} evaluated immediately --- i.e., as
soon as it is parsed and irrespective of whether any subsequent template
instantiations occur:

\begin{lstlisting}[language=C++]
void f1()
{
    static_assert(false, "Impossible!");  // always evaluated immediately...
}                                         // even if (ù{\codeincomments{f1()}}ù) is never invoked

template <typename T>
void f2()
{
    static_assert(false, "Impossible!");  // always evaluated immediately...
}                                         // even if (ù{\codeincomments{f2()}}ù) is never instantiated
\end{lstlisting}
    
\noindent The evaluation of a static assertion that is (1) located within the
body of a class or function template and (2) depends on at least one
template parameter is almost always bypassed during its initial parse since
the value --- true or false --- of the assertion will (in general) depend on the nature
of the template argument:

\begin{lstlisting}[language=C++]
template <typename T>
void f3()
{
    static_assert(sizeof(T) >= 8, "Size < 8.");  // depends on (ù{\codeincomments{T}}ù)
}
\end{lstlisting}
    
\noindent (However, see {\it\titleref{static-potential-pitfalls}: \titleref{static-assertions-in-templates-can-trigger-unintended-compilation-failures}} on page~\pageref{static-potential-pitfalls}.) In the example above, the compiler has no choice but to wait until each
time \texttt{f3} is instantiated because the truth of the predicate will
vary depending on the type provided:

\begin{lstlisting}[language=C++]
void g()
{
    f3<double>();               // OK
    f3<long double>();          // OK
    f3<std::complex<float>>();  // OK
    f3<char>();                 // Error: static assertion failed: Size < 8.
}
\end{lstlisting}
    
\noindent The standard does, however, specify that a program containing any
template definition for which no valid specialization exists is
\textbf{ill formed}\glossary{ill formed} (no diagnostic required), which was the case for
\texttt{f2} but not \texttt{f3}, above. Contrast each of the
\texttt{h*n*} definitions (below) with its correspondingly numbered
\texttt{f*n*} definition (above):

\begin{lstlisting}[language=C++]
void h1()
{
    int a[!sizeof(int) - 1];  // same as (ù{\codeincomments{int a[-1];}}ù) and is ill formed
}

template <typename T>
void h2()
{
    int a[!sizeof(int) - 1];  // always reported as a compile-time error
}

template <typename T>
void h3()
{
    int a[!sizeof(T) - 1];    // typically reported only if instantiated
}
\end{lstlisting}
    
\noindent Both \texttt{f1} and \texttt{h1} are ill-formed, non-template functions,
and both will always be reported at compile time, albeit typically with
decidedly different error messages as demonstrated by GCC 10.x's output:

\begin{lstlisting}[language=C++]
f1: error: static assertion failed: Impossible!
h1: error: size (ù{\codeincomments{-1}}ù) of array (ù{\codeincomments{a}}ù) is negative
\end{lstlisting}
    
\noindent Both \texttt{f2} and \texttt{h2} are ill-formed template functions; the cause of their being ill formed has nothing to do with the
template type and hence will always be reported as a compile-time error
in practice. Finally, \texttt{f3} can be only contextually ill formed
whereas \texttt{h3} is always necessarily ill formed, and yet neither is
reported by typical compilers as such unless and until it has been
instantiated. Reliance on a compiler not to notice that a program is
ill formed is dubious; see {\it\titleref{static-potential-pitfalls}: \titleref{static-assertions-in-templates-can-trigger-unintended-compilation-failures}} on page~\pageref{static-assertions-in-templates-can-trigger-unintended-compilation-failures}.

\subsection[Use Cases]{Use Cases}\label{use-cases}

\subsubsection[Verifying assumptions about the target platform]{Verifying assumptions about the target platform}\label{verifying-assumptions-about-the-target-platform}

Some programs rely on specific properties of the native types provided
by their target platform. Static assertions can help ensure portability
and prevent such programs from being compiled (into a malfunctioning
binary) on, say, an unsupported platform. As an example, consider a
program that relies on the size of an \texttt{int} to be exactly 32
bits (e.g., due to the use of inline \texttt{asm} blocks). Placing a
\texttt{static\_assert} in namespace scope in any of the program's
translation units will (1) ensure that the assumption regarding the size
of \texttt{int} is valid and (2) serve as documentation for readers:

\begin{lstlisting}[language=C++]
#include <ctype>  // (ù{\codeincomments{CHAR\_BIT}}ù)

static_assert(sizeof(int) * CHAR_BIT == 32,
    "An (ù{\codeincomments{int}}ù) must have exactly 32 bits for this program to work correctly.");
\end{lstlisting}
    
\noindent More typically, statically asserting the \emph{size} of an \texttt{int}
avoids having to write code to handle an \texttt{int} type's having
greater or fewer bytes when no such platforms are likely ever to
materialize:

\begin{lstlisting}[language=C++]
static_assert(sizeof(int) == 4, "An (ù{\codeincomments{int}}ù) must have exactly 4 bytes.");
\end{lstlisting}
    

\subsubsection[Preventing misuse of class and function templates]{Preventing misuse of class and function templates}\label{preventing-misuse-of-class-and-function-templates}

Static assertions are often used in practice to constrain class or
function templates to prevent their being instantiated with unsupported
types by either (1) substantially improving compile-time
diagnostics{\cprotect\footnote{Syntactically incompatible types often
lead to absurdly long and notoriously hard-to-read diagnostic messages, especially when deeply nested template expressions are involved.}} or, more critically, (2) actively avoiding
erroneous runtime behavior.
%%%  VR WANTS TO COME BACK TO THIS FOOTNOTE AND ADD A SEE ALSO. 

As an example, consider the \texttt{SmallObjectBuffer<N>} class
templates, which provide storage for arbitrary objects whose size does
not exceed \texttt{N}{\cprotect\footnote{A \texttt{SmallObjectBuffer} is
  similar to C++17's \texttt{std::any} (\textbf{{cppref\_stdany}}) in
  that it can store any object of any type. Instead of performing
  dynamic allocation to support arbitrarily sized objects, however,
  \texttt{SmallObjectBuffer} uses an internal fixed-size buffer, which
  can lead to better performance and cache locality provided (the
  maximum size of) all of the types involved is known.}}:

\begin{lstlisting}[language=C++]
template <std::size_t N>
class SmallObjectBuffer
{
private:
    char d_buffer[N];

public:
    template <typename T>
    void set(const T& object);

    // ...
};
\end{lstlisting}
    
\noindent To prevent buffer overruns, it is important that \texttt{set} accepts
only those objects that will fit in \texttt{d\_buffer}. The use of a
static assertion in the \texttt{set} member function template catches
--- at compile time --- any such misuse:

\begin{lstlisting}[language=C++]
template <std::size_t N>
template <typename T>
void SmallObjectBuffer<N>::set(const T& object)
{
    static_assert(sizeof(T) <= N, "(ù{\codeincomments{object}}ù) does not fit in the small buffer.");
    new (&d_buffer) T(object);
}
\end{lstlisting}
    

The principle of constraining inputs can be applied to most class and
function templates. \texttt{static\_assert} is particularly useful in
conjunction with standard \textbf{type traits}\glossary{type traits} provided in
\texttt{<type\_traits>}. In the \texttt{rotateLeft} function template
(below), we have used two static assertions to ensure that only unsigned
integral types will be accepted:

\begin{lstlisting}[language=C++]
#include <ctype>  // (ù{\codeincomments{CHAR\_BIT}}ù)

template <typename T>
T rotateLeft(T x)
{
    static_assert(std::is_integral<T>::value, "(ù{\codeincomments{T}}ù) must be an integral type.");
    static_assert(std::is_unsigned<T>::value, "(ù{\codeincomments{T}}ù) must be an unsigned type.");

    return (x << 1) | (x >> (sizeof(T) * CHAR_BIT - 1));
}
\end{lstlisting}
    

\subsection[Potential Pitfalls]{Potential Pitfalls}\label{static-potential-pitfalls}

\subsubsection[Static assertions in templates can trigger unintended compilation failures]{Static assertions in templates can trigger unintended compilation failures}\label{static-assertions-in-templates-can-trigger-unintended-compilation-failures}

As mentioned in the description, any program containing a template for
which no valid specialization can be generated is (by definition)
\textbf{ill formed}\glossary{ill formed} (no diagnostic required). Attempting to prevent the
use of, say, a particular function template overload by using a static
assertion that never holds produces such a program:

\begin{lstlisting}[language=C++]
template <bool>
struct SerializableTag { };

template <typename T>
void serialize(char* buffer, const T& object, SerializableTag<true>);  // (1)

template <typename T>
void serialize(char* buffer, const T& object, SerializableTag<false>)  // (2)
{
    static_assert(false, "T must be serializable.");  // independent of (ù{\codeincomments{T}}ù)
        // too obviously ill formed: always a compile-time error
}
\end{lstlisting}
    
\noindent In the example above, the second overload (2) of \texttt{serialize} is
provided with the intent of eliciting a meaningful compile-time error
message in the event that an attempt is made to serialize a
nonserializable type. The program, however, is technically
\emph{ill-formed} and, in this simple case, will likely result in a
compilation failure --- irrespective of whether either overload of
\texttt{serialize} is ever instantiated. 

A commonly attempted workaround
is to make the predicate of the assertion somehow dependent on a
template parameter, ostensibly forcing the compiler to withhold
evaluation of the \texttt{static\_assert} unless and until the template
is actually instantiated (a.k.a. \textbf{instantiation time}\glossary{instantiation time}):

\begin{lstlisting}[language=C++]
template <typename>  // N.B., we make no use of the (nameless) type parameter:
struct AlwaysFalse   // This class exists only to "outwit" the compiler.
{
    enum { value = false };
};

template <typename T>
void serialize(char* buffer, const T& object, SerializableTag<false>)  // (2)
{
    static_assert(AlwaysFalse<T>::value, "T must be serializable.");  // OK
        // less obviously ill formed: compile-time error when instantiated
}
\end{lstlisting}
    
\noindent To implement this version of the second overload, we have provided an
intermediary class template \texttt{AlwaysFalse} that, when instantiated
on any type, contains an enumerator named \texttt{value}, whose value is
\texttt{false}. Although this second implementation is more likely to
produce the desired result (i.e., a controlled compilation failure
only when \texttt{serialize} is invoked with unsuitable arguments),
sufficiently ``smart'' compilers looking at just the current translation
unit would still be able to know that no valid instantiation of
\texttt{serialize} exists and would therefore be well within their
rights to refuse to compile this still technically \emph{ill-formed}
program. 

Equivalent workarounds achieving the same result without a
helper class are possible.

\begin{lstlisting}[language=C++]
template <typename T>
void serialize(char* buffer, const T& object, SerializableTag<false>)  // (2)
{
    static_assert(0 == sizeof(T), "T must be serializable.");  // OK
         // not too obviously ill formed: compile-time error when instantiated
}
\end{lstlisting}
    
\noindent Know that use of this sort of obfuscation is not guaranteed to be either
portable or future-proof: \emph{caveat emptor}.

\subsubsection[Misuse of static assertions to restrict overload sets]{Misuse of static assertions to restrict overload sets}\label{misuse-of-static-assertions-to-restrict-overload-sets}

Even if we are careful to \emph{fool} the compiler into thinking that a
specialization is wrong \emph{only} if instantiated, we still cannot use
this approach to remove a candidate from an overload set because translation
will terminate if the static assertion is triggered. Consider this
(flawed) attempt at writing a \texttt{process} function that will behave
differently depending on the size of the given argument:

\begin{lstlisting}[language=C++]
template <typename T>
void process(const T& x)  // (1) first definition of (ù{\codeincomments{process}}ù) function
{
    static_assert(sizeof(T) <= 32, "Overload for small types");  // BAD IDEA
    // ... (process small types)
}

template <typename T>
void process(const T& x)  // (2) compile-time error: redefinition of function
{
    static_assert(sizeof(T) > 32, "Overload for big types");     // BAD IDEA
    // ... (process big types)
}
\end{lstlisting}
    
\noindent While the intention of the developer might have been to statically
dispatch to one of the two mutually exclusive overloads, the ill-fated
implementation above will not compile because the signatures of the two
overloads are identical, leading to a redefinition error. The semantics
of \texttt{static\_assert} are not suitable for the purposes of
\textbf{compile-time dispatch}\glossary{compile-time dispatch}. 

To achieve the goal of removing (up
front) a specialization from consideration, we will need to employ
\textbf{SFINAE}\glossary{SFINAE}. To do that, we must instead find a way to get the
failing compile-time expression to be part of the function's
\textbf{declaration}\glossary{declaration}{\cprotect\footnote{\textbf{Concepts}\glossary{Concepts} --- a
language feature introduced in C++20 --- provides a far less baroque
alternative to SFINAE that allows for overload sets to be governed by
  the syntactic properties of their (compile-time) template arguments.}}:

\begin{lstlisting}[language=C++]
template <bool> struct Check { };
    // helper class template having a (non-type) boolean template parameter
    // representing a compile-time predicate

template <> struct Check<true> { typedef int Ok; };
    // specialization of (ù{\codeincomments{Check}}ù) that makes the type (ù{\codeincomments{Ok}}ù) manifest *only* if
    // the supplied predicate (boolean template argument) evaluates to (ù{\codeincomments{true}}ù)

template <typename T,
          typename Check<(sizeof(T) <= 32)>::Ok = 0>  // SFINAE
void process(const T& x)  // (1)
{
    // ... (process small types)
}

template <typename T,
          typename Check<(sizeof(T) > 32)>::Ok = 0>  // SFINAE
void process(const T& x)  // (2)
{
    // ... (process big types)
}
\end{lstlisting}
    
\noindent The (empty) \texttt{Check} helper class template in conjunction with
just one of its two possible specializations (above) conditionally
exposes the \texttt{Ok} type alias \emph{only} if the provided boolean
template parameter evaluates to \texttt{true}. (Otherwise, by default, it
does not.) 

During the substitution phase of template instantiation,
exactly one of the two overloads of the \texttt{process} function will
attempt to access a nonexisting \texttt{Ok} type alias via the
\texttt{Check<false>} instantiation, which again, by default, is
nonexistent. Although such an error would typically result in a
compilation failure, in the context of template argument substitution it
will instead result in only the offending overload's being discarded,
giving other (valid) overloads a chance to be selected:

\begin{lstlisting}[language=C++]
void client()
{
    process(SmallType());  // discards (2), selects (1)
    process(BigType());    // discards (1), selects (2)
}
\end{lstlisting}
    
\noindent This general technique of pairing template specializations is used widely
in modern C++ programming. For another, often more convenient way of
constraining overloads using \textbf{expression SFINAE}\glossary{expression SFINAE}, see Section~\ref{trailing-function-return-types}, ``\titleref{trailing-function-return-types}."

\subsection[Annoyances]{Annoyances}\label{static-annoyances}

\subsubsection[Mandatory string literal]{Mandatory string literal}\label{mandatory-string-literal}

Many compilation failures caused by static assertions are
self-explanatory since the offending line (which necessarily contains the
predicate code) is displayed as part of the compiler diagnostic. In
those situations, the message required{\cprotect\footnote{As of C++17,
  the message argument of a static assertion is optional.}} as part of
\texttt{static\_assert}'s grammar is redundant:

\begin{lstlisting}[language=C++]
static_assert(std::is_integral<T>::value, "(ù{\codeincomments{T}}ù) must be an integral type.");
\end{lstlisting}
    
\noindent Developers commonly provide an empty string literal in these
cases:

\begin{lstlisting}[language=C++]
static_assert(std::is_integral<T>::value, "");
\end{lstlisting}
    

\subsection[See Also]{See Also}\label{see-also}

\begin{itemize}
\item{Section~\ref{trailing-function-return-types}, ``\titleref{trailing-function-return-types}" — Safe C++11 feature that allows fine-grained control over overload resolution by enabling \textbf{expression SFINAE}\glossary{expression SFINAE} as part of a function’s \textbf{declaration}\glossary{declaration}}
\end{itemize}

\subsection[Further reading]{Further reading}\label{further-reading}

None so far




\newpage
\section[Trailing Return]{Trailing Function Return Types}\label{trailing-function-return-types}
\setcounter{table}{0}
\setcounter{footnote}{0}
\setcounter{lstlisting}{0}
\input{trailingreturn}


\newpage
\section[{\tt alignas}]{The {\SecCode alignas} Decorator\sectionmark{\RHCode alignas}}\label{alignas}\sectionmark{\RHCode alignas}
\setcounter{table}{0}
\setcounter{footnote}{0}
\setcounter{lstlisting}{0}
% 8 Feb 2021 JMB, file comment cleanup

\emcppsFeature{
    short={\lstinline!alignas!},
    tocshort={\TOCCode alignas},
    long={The {\SecCode alignas} Decorator},
    toclong={The \lstinline!alignas! Decorator},
    rhshort={\RHCode alignas},
}{alignas}
%\section[{\tt alignas}]{The {\SecCode alignas} Decorator\sectionmark{\RHCode alignas}}\label{alignas}\sectionmark{\RHCode alignas}
\setcounter{table}{0}
\setcounter{footnote}{0}
\setcounter{lstlisting}{0}

\texttt{alignas}, a keyword that acts like an attribute, is used to widen (make more strict)
the \textbf{alignment} of a \textbf{variable}, \textbf{user-defined
type}, or \textbf{data member}.

\subsection[Description]{Description}\label{description}

The \texttt{alignas} specifier provides a means of further restricting
the granularity at which (1) a particular object of arbitrary type, (2)
a user-defined type (\texttt{class}, \texttt{struct}, \texttt{union}, or
\texttt{enum}), or (3) an individual data member is permitted to reside
within the virtual-memory-address space.

\subsubsection[Restricting the alignment of a particular object]{Restricting the alignment of a particular object}\label{restricting-the-alignment-of-a-particular-object}

In its most basic form, \texttt{alignas} acts like an
attribute that accepts (as an argument) an
\textbf{integral constant expression} representing an explicitly
supplied minimum alignment value:

\begin{lstlisting}[language=C++]
alignas(64) int i;    // OK, (ù{\codeincomments{i}}ù) is aligned on a 64-byte address boundary.
int j alignas(8), k;  // OK, (ù{\codeincomments{j}}ù) is 8-byte aligned; (ù{\codeincomments{k}}ù) remains naturally aligned.
\end{lstlisting}

\noindent If more than one alignment pertains to a given object, the most
restrictive alignment value is applied:

\begin{lstlisting}[language=C++]
alignas(4) alignas(8) alignas(2) char m;  // OK, (ù{\codeincomments{m}}ù) is 8-byte aligned.
alignas(8) int n alignas(16);             // OK, (ù{\codeincomments{n}}ù) is 16-byte aligned.
\end{lstlisting}

\noindent For a program to be \textbf{well formed}, a specified alignment value
must satisfy several\linebreak[4] \mbox{requirements}:

\begin{enumerate}
\item{Be either zero or a non-negative integral power of two of type \texttt{std::size\_t} (0, 1, 2, 4, 8, 16\dots).}
\item{Be at least the minimum alignment\cprotect\footnote{The minimum alignment of an entity is the least restrictive memory-address boundary at which the entity can be placed and have the program continue to work properly. This value is platform dependent and often subject to compiler controls but, by default, is often well approximated by \textbf{natural alignment}; see \textit{\titleref{alignas-appendix}: \titleref{natural-alignment}} on page \pageref{natural-alignment}.} required by the decorated entity.}
\item{Be no more than the largest alignment\cprotect\footnote{The notion of the largest supported alignment is characterized by both \textbf{maximal alignment} and the maximum \textbf{extended alignment}. \textbf{Maximal alignment} is defined as that most restrictive alignment that is valid in \emph{all} contexts on the current platform. All fundamental and pointer types necessarily have a minimal alignment that is less than or equal to \texttt{alignof(std::max\_align\_t)} — typically 8 or 16. Any alignment value greater than \textbf{maximal alignment} is an \textbf{extended alignment} value. Whether any extended alignment is supported (and in which contexts) is implementation defined. On typical platforms, extended alignment will often be as large as $2^{18}$ or $2^{19}$, however implementations may warn when the alignment of a global object exceeds some maximal hardware threshold (such as the size of a physical memory page, e.g., 4096 or 8192). For \textbf{automatic variables} (defined on the program stack), making alignment more restrictive than what would naturally be employed  is seldom desired because at most one thread is able to access proximately located variables there unless explicitly passed in via address to separate threads; see \textit{\titleref{alignas-use-cases}: \titleref{avoiding-false-sharing-among-distinct-objects-in-a-multi-threaded-program}} on page~\pageref{avoiding-false-sharing-among-distinct-objects-in-a-multi-threaded-program}. Note that, in the case of \texttt{i} in the first code snippet on page~\pageref{restricting-the-alignment-of-a-particular-object}, a conforming platform that did not support an extended alignment of 64 would be required to report an error at compile time.} supported on the platform in the context in which the entity appears.}
\end{enumerate}

\noindent Additionally, if the specified alignment value is zero, the
\texttt{alignas} specifier is ignored:

\begin{lstlisting}[language=C++]
// Static variables declared at namespace scope
alignas(32) int i0;  // OK, aligned on a 32-byte boundary (extended alignment)
alignas(16) int i1;  // OK, aligned on a 16-byte boundary (maximum alignment)
alignas(8)  int i2;  // OK, aligned on an 8-byte boundary
alignas(7)  int i3;  // error: not a power of two
alignas(4)  int i4;  // OK, no change to alignment boundary
alignas(2)  int i5;  // error: less than minimum alignment on this platform
alignas(0)  int i6;  // OK, (ù{\codeincomments{alignas}}ù) specifier ignored

alignas(1024 * 16) int i7;
    // OK, might warn: e.g., exceeds (physical) page size on current platform

alignas(1024 * 1024 * 512) int i8;
    // (likely) compile-time error: e.g., exceeds maximum size of object file

alignas(8) char buf[128];  // create 8-byte-aligned, 128-byte character buffer

void f()
{
    // automatic variables declared at function scope
    alignas(4)  double e0;  // error: less than minimum alignment on this platform
    alignas(8)  double e1;  // OK, no-change to (8-byte) alignment boundary
    alignas(16) double e2;  // OK, aligned to maximum (fundamental) alignment value
    alignas(32) double e3;  // OK, maximum alignment value exceeded; might warn
}
\end{lstlisting}

\subsubsection[Restricting the alignment of a user-defined type]{Restricting the alignment of a user-defined type}\label{restricting-the-alignment-of-a-user-defined-type}

The \texttt{alignas} specifier can also be used to specify alignment for
user-defined types (UDTs), such as a \texttt{class}, \texttt{struct},
\texttt{union}, or \texttt{enum}. When specifying the alignment of a UDT,
 the \texttt{alignas} keyword is placed \emph{after} the
type specifier (e.g., \texttt{class}) and just before the name of the
type (e.g., \texttt{C}):

\begin{lstlisting}[language=C++]
class alignas(2)  C { };  // OK, aligned on a 2-byte boundary; size = 2
struct alignas(4) S { };  // OK, aligned on a 4-byte boundary; size = 4
union alignas(8)  U { };  // OK, aligned on an 8-byte boundary; size = 8
enum alignas(16)  E { };  // OK, aligned on a 16-byte boundary; size = 4
\end{lstlisting}

\noindent Notice that, for each of \texttt{class}, \texttt{struct}, and
\texttt{union} above, the \texttt{sizeof} objects of that type increased
to match the alignment; in the case of the \texttt{enum}, however, the
size remains that of the default \textbf{underlying type} (e.g., 4
bytes) on the current platform.{\cprotect\footnote{When \texttt{alignas}
is applied to an enumeration \texttt{E}, the Standard does not
indicate whether padding bits are added to \texttt{E}'s object
representation or not, affecting the result of \texttt{sizeof(E)}. The
implementation variance resulting from this lack of clarity in the
  Standard was captured in \textbf{miller17}. The outcome of the core
  issue was to completely remove permission for \texttt{alignas} to be
  applied to enumerations (see \textbf{iso18a}). Therefore, conforming implementations will
  eventually stop accepting the \texttt{alignas} specifier on
  enumerations in the future.}}

Again, specifying an alignment that is less than what would occur
naturally or else is restricted explicitly is ill formed:

\begin{lstlisting}[language=C++]
struct alignas(2) T0 { int i; };
    // Error: Alignment of (ù{\codeincomments{T0}}ù) (2) is less than that of (ù{\codeincomments{int}}ù) (4).
struct alignas(1) T1 { C c; };
    // Error: Alignment of (ù{\codeincomments{T1}}ù) (1) is less than that of (ù{\codeincomments{C}}ù) (2).
\end{lstlisting}


\subsubsection[Restricting the alignment of individual data members]{Restricting the alignment of individual data members}\label{restricting-the-alignment-of-individual-data-members}

Within a user-defined type (\texttt{class}, \texttt{struct}, or
\texttt{union}), using the attribute-like syntax of the
\texttt{alignas} keyword to specify the alignments of individual data
members is possible:

\begin{lstlisting}[language=C++]
struct T2
{
    alignas(8)  char   x;  // size   1; alignment 8
    alignas(16) int    y;  // size   4; alignment 16
    alignas(64) double y;  // size   8; alignment 64
};  // size 128; alignment 64
\end{lstlisting}

\noindent The effect here is the same as if we had added the padding explicitly
and then set the alignment of the structure overall:

\begin{lstlisting}[language=C++]
struct alignas(64) T3
{
    char   x;      // size   1; alignment 8
    char   a[15];  // padding
    int    y;      // size   4; alignment 16
    char   b[44];  // padding
    double z;      // size   8; alignment 64
    char   c[56];  // padding (optional)
};  // size 128; alignment 64
\end{lstlisting}

\noindent Again, if more than one attribute pertains to a given data member,
the maximum applicable alignment value is applied:

\begin{lstlisting}[language=C++]
struct T4
{
    alignas(2) char
        c1 alignas(1),  // size 1; alignment 2
        c2 alignas(2),  // size 1; alignment 2
        c4 alignas(4);  // size 1; alignment 4
};                      // size 8; alignment 4
\end{lstlisting}


\subsubsection[Matching the alignment of another type]{Matching the alignment of another type}\label{matching-the-alignment-of-another-type}

The \texttt{alignas} specifier also accepts (as an argument) a type
identifier. In its alternate form, \texttt{alignas(T)} is strictly
equivalent to \texttt{alignas(alignof(T))}:

\begin{lstlisting}[language=C++]
alignas(int) char c;  // equivalent to (ù{\codeincomments{alignas(alignof(int)) char c;}}ù)
\end{lstlisting}


\subsection[Use Cases]{Use Cases}\label{alignas-use-cases}

\subsubsection[Creating a sufficiently aligned object buffer]{Creating a sufficiently aligned object buffer}\label{creating-a-sufficiently-aligned-object-buffer}

When writing low-level, system-infrastructure code, constructing an object within a raw buffer is sometimes useful. As a minimal
example, consider a function that uses a local character buffer to
create an object of type \texttt{std::complex<long}~\texttt{double>} on
the program stack using placement \texttt{new}:

\begin{lstlisting}[language=C++]
void f()
{
    // ...
    char objectBuffer[sizeof(std::complex<long double>)];  // BAD IDEA
    // ...
    new(objectBuffer) std::complex<long double>(1.0, 0.0);  // Might dump core!
    // ...
}
\end{lstlisting}

\noindent The essential problem with the code above is that \texttt{objectBuffer},
being an array of characters (each having an alignment of 1), is itself
byte aligned. The compiler is therefore free to place it on any address
boundary. On the other hand, \texttt{std::complex<long}~\texttt{double>} is an aggregate consisting of two \texttt{long}~\texttt{double} objects
and therefore necessarily requires (at least) the same strict alignment
(typically 16) as the two \texttt{long}~\texttt{double} objects it comprises. Previous
solutions to this problem involved creating a \texttt{union} of the
object buffer and some maximally aligned type (e.g.,
\texttt{std::max\_align\_t}):

\begin{lstlisting}[language=C++]
#include <cstddef>  // (ù{\codeincomments{std::max\_align\_t}}ù)

void f()
{
    // ...

    union {                                              // awkward workaround
        std::max_align_t dummy;  // (ù{\codeincomments{typedef}}ù) to maximally aligned type
        char objectBuffer[sizeof(std::complex<long double>)];
    } objectBuffer;

    // ...

    new(&objectBuffer) std::complex<long double>(1.0, 0.0);  // OK

    // ...
}
\end{lstlisting}

\noindent Using the alternate syntax for \texttt{alignas}, we can avoid gratuitous
complexity and just state our intentions explicitly:

\begin{lstlisting}[language=C++]
void f()
{
    // ...

    alignas(std::complex<long double>) char objectBuffer[
                              sizeof(std::complex<long double>)];  // GOOD IDEA

    // ...

    new(objectBuffer) std::complex<long double>(1.0, 0.0);  // OK

    // ...
}
\end{lstlisting}


\subsubsection[Ensuring proper alignment for architecture-specific instructions]{Ensuring proper alignment for architecture-specific instructions}\label{ensuring-proper-alignment-for-architecture-specific-instructions}

Architecture-specific instructions or compiler intrinsics might require
the data they act on to have a specific alignment. One example of such
intrinsics is the \emph{Streaming SIMD Extensions (SSE)}\footnote{\textbf{inteliig}, ``Technologies: SSE"} instruction set available on the x86
architecture. SSE instructions operate on groups of four 32-bit
single-precision floating-point numbers at a time, which are required to
be 16-byte aligned.{\cprotect\footnote{``Data must be
16-byte aligned when loading to and storing from the 128-bit XMM
  registers used by SSE/SSE2/SSE3/SSSE3'': see \textbf{{intel16}},
  section 4.4.4, ``Data Alignment for 128-Bit Data," pp. 4-19--4-20.}} The
\texttt{alignas} specifier can be used to create a type satisfying this requirement:

\begin{lstlisting}[language=C++]
struct SSEVector
{
    alignas(16) float d_data[4];
};
\end{lstlisting}

\noindent Each object of the \texttt{SSEVector} type above is guaranteed always to
be aligned to a 16-byte boundary and can therefore be safely (and
conveniently) used with SSE intrinsics:

\begin{lstlisting}[language=C++]
#include <xmmintrin.h>  // (ù{\codeincomments{\_\_m128}}ù) and (ù{\codeincomments{\_mm\_XXX}}ù) functions

void f()
{
    const SSEVector v0 = {0.0f, 1.0f, 2.0f, 3.0f};
    const SSEVector v1 = {10.0f, 10.0f, 10.0f, 10.0f};

    __m128 sseV0 = _mm_load_ps(v0.d_data);
    __m128 sseV1 = _mm_load_ps(v1.d_data);
        // (ù{\codeincomments{\_mm\_load\_ps}}ù) requires the given (ù{\codeincomments{float}}ù) array to be 16-byte aligned.
        // The data is loaded into a dedicated 128-bit CPU register.

    __m128 sseResult = _mm_add_ps(sseV0, sseV1);
        // sum two 128-bit registers; typically generates an (ù{\codeincomments{addps}}ù) instruction

    SSEVector vResult;
    _mm_store_ps(vResult.d_data, sseResult);
        // Store the result of the sum back into a (ù{\codeincomments{float}}ù) array.

    assert(vResult.d_data[0] == 10.0f);
    assert(vResult.d_data[1] == 11.0f);
    assert(vResult.d_data[2] == 12.0f);
    assert(vResult.d_data[3] == 13.0f);
}
\end{lstlisting}


\subsubsection[Avoiding \textbf{false sharing} among distinct objects in a multi-threaded program]{Avoiding \textbf{false sharing} among distinct objects in a multi-threaded program}\label{avoiding-false-sharing-among-distinct-objects-in-a-multi-threaded-program}

In the context of an application where multithreading has been employed
to improve performance, seeing
a previously single-threaded workflow become even less performant after
a parallelization attempt can be surprising (and disheartening). One possible insidious cause of such
disappointing results comes from \textbf{false sharing} --- a situation
in which multiple threads unwittingly harm each other's performance
while writing to logically independent variables that happen to reside
on the same \textbf{cache line}; see {\it \titleref{alignas-appendix}:} {\it\titleref{cache-lines,-l1,-l2,-and-l3-cache,-pages,-and-virtual-memory}} on page~\pageref{cache-lines,-l1,-l2,-and-l3-cache,-pages,-and-virtual-memory}.

As a simple (purely pedagogical) illustration of the potential
performance degradation resulting from \textbf{false sharing}, consider
a function that spawns separate threads to repeatedly increment
(concurrently) logically distinct variables that happen to reside in
close proximity on the program stack:

\begin{lstlisting}[language=C++]
#include <thread>  // (ù{\codeincomments{std::thread}}ù)

volatile int target = 0;  // updated asynchronously from multiple threads

void incrementJob(int* p);
    // Repeatedly increment (ù{\codeincomments{*p}}ù) a large, fixed number of times;
    // periodically write its current value to (ù{\codeincomments{target}}ù).

void f()
{
    int i0 = 0;  // Here, (ù{\codeincomments{i0}}ù) and (ù{\codeincomments{i1}}ù) likely share the same cache line,
    int i1 = 0;  // i.e., byte-aligned memory block on the program stack.

    std::thread t0(&incrementJob, &i0);
    std::thread t1(&incrementJob, &i1);
        // Spawn two parallel jobs incrementing the respective variables.

    t0.join();
    t1.join();
        // Wait for both jobs to be completed.
}
\end{lstlisting}

\noindent In the simplistic example above, the proximity in memory between
\texttt{i0} and \texttt{i1} can result in their belonging to the same
\textbf{cache line}, thus leading to \textbf{false sharing}. By
prepending \texttt{alignas(64)} to the declaration of both integers, we
ensure that the two variables reside on distinct cache lines:

\begin{lstlisting}[language=C++]
// ...

void f()
{
    alignas(64) int i0 = 0;   // Assuming a cache line on this platform is 64
    alignas(64) int i1 = 0;   // bytes, (ù{\codeincomments{i0}}ù) and (ù{\codeincomments{i1}}ù) will be on separate ones.

    // ...
\end{lstlisting}

\noindent As an empirical demonstration of the effects of \textbf{false sharing},
a benchmark program repeatedly calling \texttt{f} completed its
execution seven times faster on average when compared to the same
program without use of \texttt{alignas}.{\cprotect\footnote{The benchmark
program was compiled using Clang 11.0.0 using \texttt{-Ofast},
\texttt{-march=native}, and \mbox{\texttt{-std=c++11}}. The program was then
executed on a machine running Windows 10 x64, equipped with an Intel
Core i7-9700k CPU (8 cores, 64-byte cache line size). Over the
course of multiple runs, the version of the benchmark without
\texttt{alignas} took 18.5967ms to complete (on average), while the
version with \texttt{alignas} took 2.45333ms to complete (on average).
  See \textbf{{[PRODUCTION: CODE PROVIDED WITH BOOK] alignasbenchmark}} for the source code of the program.}}

\subsubsection[Avoiding \textbf{false sharing} within a single thread-aware object]{Avoiding \textbf{false sharing} within a single thread-aware object}\label{avoiding-false-sharing-within-a-single-thread-aware-object}

A real-world scenario where the need for preventing \textbf{false
sharing} is fundamental occurs in the implementation of high-performance
concurrent data structures. As an example, a thread-safe ring buffer
might make use of \texttt{alignas} to ensure that the indices of the
head and tail of the buffer are aligned at the start of a cache line
(typically 64, 128, or 256 bytes), thereby preventing them from
occupying the same one.

\begin{lstlisting}[language=C++]
class ThreadSafeRingBuffer
{
    alignas(cpuCacheSize) std::atomic<std::size_t> d_head;
    alignas(cpuCacheSize) std::atomic<std::size_t> d_tail;

    // ...
};
\end{lstlisting}

\noindent Not aligning \texttt{d\_head} and \texttt{d\_tail} (above) to the CPU
cache size might result in poor performance of the
\texttt{ThreadSafeRingBuffer} because CPU cores that need to access only one
of the variables will inadvertently load the other one as well,
triggering expensive hardware-level coherency mechanisms between the
cores' caches. On the other hand, specifying such substantially stricter
alignment on consecutive data members necessarily increases the size of
the object; see {\it\titleref{potential-pitfalls}:} {\it\titleref{stricter-alignment-might-reduce-cache-utilization}} on page~\pageref{stricter-alignment-might-reduce-cache-utilization}.

\subsection[Potential Pitfalls]{Potential Pitfalls}\label{potential-pitfalls}

\subsubsection[Underspecifying alignment is not universally reported]{Underspecifying alignment is not universally reported}\label{underspecifying-alignment-is-not-universally-reported}

The Standard is clear when it comes to underspecifying
alignment{\cprotect\footnote{\textbf{{cpp11}}, section 7.6.2, ``Alignment Specifier," paragraph~5, pp. 179}}:
\begin{quote}
The combined effect of all \textit{alignment-specifiers} in a declaration
shall not specify an alignment that is less strict than the alignment
that would be required for the entity being declared if all
\textit{alignment-specifiers} were omitted (including those in other
declarations).
\end{quote}
The compiler is required to honor the
specified value if it is a \textbf{fundamental
alignment},{\cprotect\footnote{``If the constant expression evaluates to
a fundamental alignment, the alignment requirement of the declared
entity shall be the specified fundamental alignment'': \textbf{cpp11}, section~7.6.2, ``Alignment Specifier," paragraph~2, item~2, p. 178.}} so
imagining how this would lead to anything other than an ill-formed
program is difficult:

\begin{lstlisting}[language=C++]
alignas(4) void* p;              // (1) Error: (ù{\codeincomments{alignas(4)}}ù) is below minimum, 8.

struct alignas(2) S { int x; };  // (2) Error: (ù{\codeincomments{alignas(2)}}ù) is below minimum, 4.

struct alignas(2) T { };
struct alignas(1) U { T e; };    // (3) Error: (ù{\codeincomments{alignas(1)}}ù) is below minimum, 2.
\end{lstlisting}

\noindent Each of the three errors above are reported by Clang, but GCC
doesn't issue so much as a warning (let alone the required error) ---
even in the most pedantic warning mode. Thus, one could write a program,
involving statements like those above, that happens to work on one
platform (e.g., GCC) but fails to compile on another (e.g.,
Clang).{\cprotect\footnote{Underspecifying alignment is not reported at
all by GCC 10.1, using the
\texttt{-std=c++11} \texttt{-Wall} \texttt{-Wextra} \mbox{\texttt{-Wpedantic}}
flags. With the same set of options, Clang 10.0 produces a compilation
failure. MSVC v19.24 will produce a warning and ignore any alignment
  less than the minimum one.}}

\subsubsection[Incompatibly specifying alignment is \textbf{IFNDR}]{Incompatibly specifying alignment is \textbf{IFNDR}}\label{incompatibly-specifying-alignment-is-ifndr}

It is permissible to forward declare a user-defined type (UDT)
without an \texttt{alignas} specifier so long as all defining
declarations of the type have either no \texttt{alignas} specifier or
have the same one. Similarly, if any forward declaration of a
user-defined type has an \texttt{alignas} specifier, then all defining
declarations of the type must have the same specifier and that specifier
must be \emph{equivalent to} (not necessarily \emph{the same as}) that
in the forward declaration:

\begin{lstlisting}[language=C++]
struct Foo;                  // OK, does not specify an alignment
struct alignas(double) Foo;  // OK, must be equivalent to every definition
struct alignas(8) Foo;       // OK, all definitions must be identical.
struct alignas(8) Foo { };   // OK, equivalent to each decl. specifying (ù{\codeincomments{alignas}}ù)
struct Foo;                  // OK, has no effect
struct alignas(8) Foo;       // OK, has no effect; might warn after definition
\end{lstlisting}

\noindent Specifying an alignment in a forward declaration without specifying an
equivalent one in the defining declaration is \textbf{ill formed; no diagnostic
is required (IFNDR)} if the two declarations appear in distinct translation
units:

\begin{lstlisting}[language=C++]
struct alignas(4) Bar;      // OK, forward declaration
struct Bar { };             // error: missing (ù{\codeincomments{alignas}}ù) specifier

struct alignas(4) Baz;      // OK, forward declaration
struct alignas(8) Baz { };  // error: non-equivalent (ù{\codeincomments{alignas}}ù) specifier
\end{lstlisting}

\noindent Both of the errors above are flagged by Clang, but neither of them is
reported by GCC. Note that when the inconsistency occurs across
translation units, no mainstream compiler is likely to diagnose the
problem:

\begin{lstlisting}[language=C++]
// file1.cpp:
struct Bam { char ch; } bam, *p = &bam;

// file2.cpp:
struct alignas(int) Bam;  // Error: definition of (ù{\codeincomments{Bam}}ù) lacks alignment specifier.
extern Bam* p;            //        (no diagnostic required)
\end{lstlisting}

\noindent Any program incorporating both translation units above is
\textbf{ill formed, no diagnostic required}.

\subsubsection[Stricter alignment might reduce cache utilization]{Stricter alignment might reduce cache utilization}\label{stricter-alignment-might-reduce-cache-utilization}

User-defined types having artificially stricter alignments than would
naturally occur on the host platform means that fewer of them can fit
within any given level of physical cache within the hardware. Types
having data members whose alignment is artificially widened tend to be
larger and thus suffer the same lost cache utilization. As an
alternative to enforcing stricter alignment to avoid \textbf{false
sharing}, consider organizing a multithreaded program such that tight
clusters of repeatedly accessed objects are always acted upon by only a
single thread at a time, e.g., using local (arena) memory allocators;
see {\it\titleref{alignas-appendix}:} {\it\titleref{cache-lines,-l1,-l2,-and-l3-cache,-pages,-and-virtual-memory}} on page~\pageref{cache-lines,-l1,-l2,-and-l3-cache,-pages,-and-virtual-memory}.

\subsection[See Also]{See Also}\label{see-also}

\begin{itemize}
\item{Section~\ref{alignof}, ``\titleref{alignof}" — Safe C++11 feature that inspects the alignment of a given type}
\item{Section~\ref{attributes}, ``\titleref{attributes}" — Safe C++11 feature that shows how other attributes (following the conventional attribute notation) are used to annotate source code, improve error diagnostics, and implicitly code generation} %%% WORDS MISSING!
\end{itemize}

\subsection[Further Reading]{Further Reading}\label{further-reading}

None so far

\subsection[Appendix]{Appendix}\label{alignas-appendix}

\subsubsection[Natural Alignment]{Natural Alignment}\label{natural-alignment}

By default, fundamental, pointer, and enumerated types typically reside
on an address boundary that divides the size of the object; we refer to
such alignment as \textbf{natural alignment}{\cprotect\footnote{Sizes
and alignment shown here are typical but not specifically required by
the standard. On some platforms, one can request that all
types be \textbf{byte aligned}. While such a representation is more
compact, entities that span memory boundaries can require multiple
fetch operations leading to run times that are typically
significantly (sometimes as much as an order of magnitude) slower when
  run in this ``packed'' mode.}}:

\begin{lstlisting}[language=C++]
char   c;  // size 1; alignment 1; boundaries: 0x00, 0x01, 0x02, 0x03, ...
short  s;  // size 2; alignment 2: boundaries: 0x00, 0x02, 0x04, 0x06, ...
int    i;  // size 4; alignment 4; boundaries: 0x00, 0x04, 0x08, 0x0c, ...
float  f;  // size 4; alignment 4; boundaries: 0x00, 0x04, 0x08, 0x0c, ...
double d;  // size 8; alignment 8; boundaries: 0x00, 0x08, 0x10, 0x18, ...
\end{lstlisting}

\noindent For aggregates (including arrays) or user-defined types, the alignment
is typically that of the most strictly aligned subelement:

\begin{lstlisting}[language=C++]
struct S0
{
    char a;  // size 1; alignment 1
    char b;  // size 1; alignment 1
    int  c;  // size 4; alignment 4
};           // size 8; alignment 4

struct S1
{
    char a;  // size  1; alignment 1
    int  b;  // size  4; alignment 4
    char c;  // size  1; alignment 1
};           // size 12; alignment 4

struct S2
{
    int  a;  // size 4; alignment 4
    char b;  // size 1; alignment 1
    char c;  // size 1; alignment 1
};           // size 8; alignment 4

struct S3
{
    char a;  // size 1; alignment 1
    char b;  // size 1; alignment 1
};           // size 2; alignment 1

struct S4
{
    char a[2];  // size 2; alignment 1
};              // size 2; alignment 1
\end{lstlisting}

\noindent Size and alignment behave similarly with respect to \textbf{structural
inheritance}:

\begin{lstlisting}[language=C++]
struct D0 : S0
{
    double d;  // size  8; alignment 8
};             // size 16; alignment 8

struct D1 : S1
{
    double d;  // size  8; alignment 8
};             // size 24; alignment 8

struct D2 : S2
{
    int d;  // size  4; alignment 4
};          // size 12; alignment 4

struct D3 : S3
{
    int d;  // size 4; alignment 4
};          // size 8; alignment 4

struct D4 : S4
{
    char d;  // size 1; alignment 1
};           // size 3; alignment 1
\end{lstlisting}

\noindent Finally, virtual functions invariably introduce an implicit
virtual-table-pointer member having a size and alignment corresponding
to that of a memory address (e.g., 4 or 8) on the host platform:

\begin{lstlisting}[language=C++]
struct S5
{
    virtual ~S5();
};                   // size 8; alignment 8

struct D5 : S5
{
    char d;  // size  1; alignment 1
};           // size 16; alignment 8
\end{lstlisting}


\subsubsection[Cache lines; L1, L2, and L3 cache; pages; and virtual memory]{Cache lines; L1, L2, and L3 cache; pages; and virtual memory}\label{cache-lines,-l1,-l2,-and-l3-cache,-pages,-and-virtual-memory}

Modern computers are highly complex systems, and a detailed
understanding of their intricacies is unnecessary to achieve most of
the performance benefits. Still, certain general themes and
rough thresholds aid in understanding how to squeeze just a bit
more out of the underlying hardware. In this section, we sketch
fundamental concepts that are common to all modern computer hardware;
although the precise details will vary, the general ideas remain
essentially the same.

In its most basic form, a computer consists of central processing unit
(CPU) having internal registers that access main memory (MM).
Registers in the CPU (on the order of hundreds of bytes) are among the
fastest forms of memory, while main memory (typically many gigabytes) is
orders of magnitude slower. An almost universally observed phenomenon is
that of \textbf{locality of reference}, which suggests that data that
resides in close proximity (in the virtual address space) is more likely
to be accessed together in rapid succession than more distant data.

To exploit the phenomenon of \textbf{locality of reference}, computers
introduce the notion of a cache that, while much faster than main
memory, is also much smaller. Programs that attempt to amplify
\textbf{locality of reference} will, in turn, often be rewarded with
faster run times. The organization of a cache and, in fact, the number of
levels of cache (e.g., L1, L2, L3,$\ldots$) will vary, but the basic
design parameters are, again, more or less the same. A given level of
cache will have a certain total size in bytes (invariably an integral
power of two). The cache will be segmented into what are called
\textbf{cache lines} whose size (a smaller power of two) divides that of
the cache itself. When the CPU accesses main memory, it first looks to
see if that memory is in the cache; if it is, the value is returned
quickly (known as a \textbf{cache hit}). Otherwise, the cache line(s)
containing that data is (are) fetched (from the next higher level of cache
or from main memory) and placed into the cache (known as a \textbf{cache
miss}), possibly ejecting other less recently used
ones.{\cprotect\footnote{Conceptually, the cache is often thought of as
being able to hold any arbitrary subset of the most recently accessed
cache lines. This kind of cache is known as \textbf{fully
associative}. Although it provides the best hit rate, a \textbf{fully
associative} cache requires the most power along with significant
additional chip area to perform the fully parallel lookup. \textbf{Direct-mapped} cache associativity is at the
other extreme. In direct mapped, each memory location has exactly one location
available to it in the cache. If another memory location mapping to
that location is needed, the current cache line must be flushed from
the cache. Although this approach has the lowest hit rate, lookup
times, chip area, and power consumption are all minimized (optimally).
Between these two extremes is a continuum that is referred to as
\textbf{set associative}. A \textbf{set associate} cache has more than
  one (typically 2, 4, or 8; see \textbf{solihin15}, section~5.2.1, ``Placement Policy," pp. 136--141, and \textbf{hruska20})
  location in which each memory location in main memory can reside.
  Note that, even with a relatively small $N$, as $N$ increases, an $N$-way
  \textbf{set associative} cache quickly approaches the hit rate of a fully
  associative cache at greatly reduced collateral cost; for most
  software-design purposes, any loss in hit rate due to set
  associativity of a cache can be safely ignored.}}

Data residing in distinct cache lines is physically independent and can
be written concurrently by multiple threads. Logically unrelated data
residing in the same cache line, however, is nonetheless physically
coupled; two threads that write to such logically unrelated data will find
themselves synchronized by the hardware. Such unexpected and typically
undesirable sharing of a cache line by unrelated data acted upon by two
concurrent threads is known as \textbf{false sharing}. One way of
avoiding \textbf{false sharing} is to align such data on a cache-line
boundary, thus rendering accidental collocation of such data on the same
cache line impossible. Another (more broad-based) design approach that
avoids lowering cache utilization is to ensure that data acted upon by a
given thread is kept physically separate --- e.g., through the use of
local (arena) memory allocators.{\cprotect\footnote{\textbf{{lakos17}},
  \textbf{{lakos19}}, \textbf{{lakos22}}}}

Finally, even data that is not currently in cache but resides nearby in
MM can benefit from locality. The virtual address space, synonymous with
the size of a \texttt{void*} (typically 64-bits on modern general-purpose hardware), has historically well exceeded the physical memory
available to the CPU. The operating system must therefore maintain a
mapping (in main memory) from what is resident in physical memory and
what resides in secondary storage (e.g., on disc). In addition,
essentially all modern hardware provides a
\textbf{TLB}{\cprotect\footnote{A translation-lookaside buffer (TLB) is
a kind of address-translation cache that is typically part of a chip's
memory management unit (MMU). A TLB holds a recently accessed subset
of the complete mapping (itself maintained in MM) from virtual memory
address to physical ones. A TLB is used to reduce access time when the
requisite pages are already resident in memory; its size (e.g., 4K) is
capped at the number of bytes of physical memory (e.g., 32Gb) divided
by the number of bytes in each physical page (e.g., 8Kb), but could be
smaller. Because it resides on chip, is typically an order of
magnitude faster (SRAM versus DRAM), and requires only a single lookup
(as opposed to two or more when going out to MM), there is an enormous
  premium on minimizing TLB misses.}} that caches the addresses of the
most recently accessed physical pages, providing yet another advantage
to having the \textbf{working set} (i.e., the current set of frequently
accessed pages) remain small and densely packed with relevant
data.{\cprotect\footnote{Note that memory for handle-body
types (e.g., \texttt{std::vector} or \texttt{std::deque}) and
especially node-based containers (e.g., \texttt{std::map} and
\texttt{std::unordered\_map}), originally allocated within a single
page, can --- through deallocation and reallocation (or even move
operations) --- become scattered across multiple (perhaps many)
 pages, thus causing what was originally a relatively small \textbf{working set}
to no longer fit within physical memory. This phenomenon, known as
\textbf{diffusion} (which is a distinct concept from
\textbf{fragmentation}), is what typically leads to a substantial
runtime performance degradation (due to \textbf{thrashing}) in large,
long-running programs. Such \textbf{diffusion} can be mitigated by
judicious use of local arena memory allocators (and deliberate
avoidance of \textbf{move operations} across disparate localities of
  frequent memory usage).}} What's more, dense working sets, in addition
to facilitating hits for repeat access, increase the likelihood that
data that is coresident on a page (or cache line) will be needed soon
(i.e., in effect acting as a prefetch).{\cprotect\footnote{We sometimes
lightheartedly refer to the beneficial prefetch of unrelated data that
is accidentally needed subsequently (e.g., within a single thread) due
to high locality within a cache line (or a physical page) as
  \textbf{true sharing}.}} Table~\ref{table-alignas-appendix} provides a summary of typical physical parameters found in modern computers today.

%%%%%%TODO: TBD: @Lori: Format this as a figure with a caption. @jslakos: Double-check accuracy of values.

\begin{table}[h!]
\begin{center}
\begin{threeparttable}
\caption{Various sizes and access speeds of typical memory for modern \mbox{computers}}\label{table-alignas-appendix} \vspace{1.5ex}
{\small \begin{tabular}{c|c|c}\thickhline
\rowcolor[gray]{.9}    {\sffamily\bfseries Memory Type} & {\sffamily\bfseries Typical Memory Size (Bytes)} & {\sffamily\bfseries Typical Access Times}\\\hline
CPU Registers & 512 \ldots~2048 & $\sim$250ps\\ \hline
Cache Line  & 64 \ldots~256 & NA\\ \hline
L1 Cache & 16Kb \ldots~64Kb & $\sim$1ns\\ \hline
L2 Cache & 1Mb \ldots~2Mb & $\sim$10ns\\ \hline
L3 Cache & 8Mb \ldots~32Mb & $\sim$80ns--120ns\\ \hline
L4 Cache & 32Mb \ldots~128Mb & $\sim$100ns--200ns\\ \hline
Set Associativity  &  2 \ldots~64 & NA\\ \hline
TL & 4 words \ldots~65536 & 10ns \ldots~50ns\\ \hline
Physical Memory Page & 512 \ldots~8192 & 100ns \ldots~500ns\\ \hline
\rule{0pt}{3.5mm}Virtual Memory & $2^{32}$ bytes \ldots~$2^{64}$ bytes & $\sim$10$\mu$s--50$\mu$s\\ \hline
Solid-State Disc (SSD) & 256Gb \ldots~16Tb & $\sim$25$\mu$s--100$\mu$s\\ \hline
Mechanical Disc & Huge & $\sim$5ms--10ms\\ \hline
Clock Speed & NA & $\sim$4GHz \\ \thickhline
\end{tabular}
}
\end{threeparttable}
    \end{center}
\end{table}
\clearpage



\newpage
\section[{\tt deprecated}]{The {\SecCode [[deprecated]]} Attribute}\label{deprecated}\label{the-standard-[[deprecated]]-attribute}
\setcounter{table}{0}
\setcounter{footnote}{0}
\setcounter{lstlisting}{0}
\input{deprecated}

\newpage
\section[Raw String Literals '14]{Syntax for Unprocessed String Contents}
\setcounter{table}{0}
\setcounter{footnote}{0}
\setcounter{lstlisting}{0}
PLACEHOLDER

\newpage
\section[{\tt using} Aliases]{Type/Template Aliases (Extended {\SecCode typedef})}\label{alias-declarations-and-alias-templates}
PLACEHOLDER

\newpage
\section[Defaulted Functions]{Using {\SecCode =default} for Special Member Functions}\label{Defaulted-Special-Member-Functions}\label{defaulted-special-member-functions}
PLACEHOLDER

\newpage
\section[Attribute Syntax]{Generalized Attribute Support}\label{attributes}
PLACEHOLDER




%\addtocontents{toc}{\newpage}


%%%%%%%%%%%%%%% C++14
%\cftaddtitleline{toc}{chapter}{\sffamily\itshape C++14}{}
% to add to the TOC at the section level
\addtocontents{toc}{\protect\renewcommand*\protect\cftsecindent{1.25em}}
\stepcounter{cppxx}
\cftaddnumtitleline{toc}{section}{\thechapter.\thecppxx}{C++14}{\thepage}
\renewcommand{\cppxx}{C++14}

\newpage
\addtocontents{toc}{\protect\renewcommand*\protect\cftsecindent{4.5em}}
\section[Digit Separators]{The Digit Separator: {\SecCode '}}
\label{digitseparator} % used with placeholder file
\label{digit-separators} % label had changed when actual file arrived
\setcounter{table}{0}
\setcounter{footnote}{0}
\setcounter{lstlisting}{0}
PLACEHOLDER

\newpage
\section[{\tt constexpr} Functions '14]{Relaxed Restrictions on {\SecCode constexpr} Functions}
\setcounter{table}{0}
\setcounter{footnote}{0}
\setcounter{lstlisting}{0}
PLACEHOLDER

\newpage
\section[{\tt code0} in Short]{Longer Wordier Feature Title with {\SecCode Code0}\sectionmark{{\RHCode code0} in Short}}\label{Feature}\label{testmaterial}\sectionmark{{\RHCode code0} in Short}
\setcounter{table}{0}
\setcounter{footnote}{0}
\setcounter{lstlisting}{0}


Lorem ipsum dolor sit amet, consectetur
adipiscing elit. Suspendisse non nisi ex. Ut non sollicitudin est. Donec
laoreet, risus pretium egestas condimentum, ex eros tempor diam, quis
congue nisi metus in tellus. Proin tempor ac lectus nec elementum.
Maecenas augue turpis, pellentesque sed eros sit amet, tincidunt pretium
nulla. Nunc lacus ligula, ullamcorper a porta eget, fringilla sed orci.
Maecenas eget ultricies risus. Donec varius vehicula diam.

\subsection[Description]{Description}\label{description}

Lorem ipsum dolor sit amet, consectetur
adipiscing elit. Suspendisse non nisi ex. Ut non sollicitudin est. Donec
laoreet, risus pretium egestas condimentum, ex eros tempor diam, quis
congue nisi metus in tellus. Proin tempor ac lectus nec elementum.
Maecenas augue turpis, pellentesque sed eros sit amet, tincidunt pretium
nulla. Nunc lacus ligula, ullamcorper a porta eget, fringilla sed orci.
Maecenas eget ultricies risus. Donec varius vehicula diam.

\subsubsection[C++ attribute syntax]{C++ attribute syntax}\label{c++-attribute-syntax}

Lorem ipsum dolor sit amet, consectetur
adipiscing elit.{\cprotect\footnote{\authorsnote We will have
some footnotes that are authors' notes.}} Suspendisse non nisi ex. Ut non sollicitudin est. Donec
laoreet, risus pretium egestas condimentum, ex eros tempor diam, quis
congue nisi metus in tellus. Proin tempor ac lectus nec elementum.
Maecenas augue turpis, pellentesque sed eros sit amet, tincidunt pretium
nulla. Nunc lacus ligula, ullamcorper a porta eget, fringilla sed orci.
Maecenas eget ultricies risus. Donec varius vehicula diam.

Nested unordered list:

\begin{itemize}
\item{One\begin{itemize}
\item{Sub}
\item{Sub}
\end{itemize}
}
\item{Two}
\item{Three}
\end{itemize}

Donec
laoreet, risus pretium egestas condimentum, ex eros tempor diam, quis
congue nisi metus in tellus. Proin tempor ac lectus nec elementum.
Maecenas augue turpis, pellentesque sed eros sit amet, tincidunt pretium
nulla:

Nested ordered list:

\begin{enumerate}\setlength{\leftmarginii}{6mm}
\item{One\begin{enumerate}
\item{Sub}
\item{Sub}
\end{enumerate}
}
\item{Two}
\item{Three}
\end{enumerate}

This feature, when used in conjunction with \emph{explicit instantiation
definitions}, can significantly improve compilation times for a set of
translation units that often instantiate common templates:

\noindent\begin{minipage}{.45\textwidth}
\begin{lstlisting}[language=C++, caption=code 1,frame=tb]{Name}
void code()
{

}
\end{lstlisting}
\end{minipage}\hfill
\begin{minipage}{.45\textwidth}
\begin{lstlisting}[language=C++, caption=code 2,frame=tb]{Name}
void code()
{

}
\end{lstlisting}
\end{minipage}

Nullam nibh tortor, finibus ut lectus eu, convallis vehicula libero. Cras maximus ligula nisl, a eleifend mauris venenatis nec. Suspendisse potenti. Maecenas viverra laoreet mauris ut laoreet. Donec non felis risus. Ut faucibus, justo id luctus lobortis, nisi lacus pharetra est, id pretium arcu ligula ac magna. Morbi a nisl sit amet lacus vehicula tempor nec sit amet nibh. 

\noindent\begin{minipage}{.45\textwidth}
\begin{lstlisting}[language=C++, caption=code 1 with a long and wrapping title,frame=tb]{Name}
void code()
{

}
\end{lstlisting}
\end{minipage}\hfill
\begin{minipage}{.45\textwidth}
\begin{lstlisting}[language=C++, caption=code 2 with a long and wrapping title,frame=tb]{Name}
void code()
{

}
\end{lstlisting}
\end{minipage}

Pellentesque id lorem ac sem ullamcorper semper. Donec imperdiet sapien in nisi faucibus, et tempus lacus hendrerit. Nulla laoreet risus eu tortor ultrices, sed bibendum neque sodales. Ut faucibus ipsum id convallis convallis. Vivamus vehicula rutrum metus in semper. Morbi fringilla ex vel vulputate vehicula. Maecenas ut porttitor massa.

\subsection[Appendix to the Feature]{Appendix to the Feature}
Suspendisse lorem libero, egestas non semper eleifend, finibus ut purus. Quisque leo nulla, lacinia vel dui et, vestibulum facilisis erat. Vestibulum scelerisque auctor diam, ut fringilla diam elementum non. Vivamus sed mauris lobortis, blandit urna quis, iaculis neque. Phasellus sit amet venenatis sapien, vel vestibulum augue. In pellentesque sit amet enim a rutrum. Fusce nec nulla et nisl faucibus efficitur vel non diam. Nulla lobortis feugiat augue, a lobortis nunc imperdiet eget. Sed in neque ultricies, efficitur mi eu, sollicitudin massa. Vivamus sed dui convallis, mattis sem a, scelerisque augue. Sed non ultricies neque. Praesent tincidunt feugiat lorem, a sagittis dui eleifend ac. Nullam molestie nulla quis risus molestie dignissim. Aenean egestas enim ut tellus vulputate finibus. Integer iaculis sodales gravida.

\begin{table}[h!] 
\begin{center}
\begin{threeparttable}
\caption{Testing Table Numbering}\label{test-table1}\vspace{1.5ex} 
{\small \begin{tabular}{c|c|c}\thickhline 
\rowcolor[gray]{.9}   {\sffamily\bfseries Compiler} & {\sffamily\bfseries Compiler-Specific} &
{\sffamily\bfseries Standard-Conforming} \\ \hline 
GCC &\texttt{\_\_attribute\_\_((pure))} & \texttt{[[gnu::pure]]} \\ \hline
Clang & \texttt{\_\_attribute\_\_((no\_sanitize))} &\texttt{[[clang::no\_sanitize]]} \\ \hline 
MSVC & \texttt{declspec(deprecated)} & \texttt{[[deprecated]]} \\ \thickhline
\end{tabular}
}
\end{threeparttable} 
\end{center}
\end{table}

Donec mattis, ex a ornare auctor, tortor felis pulvinar odio, eu mollis tortor orci id ligula. Vestibulum dignissim magna vitae lectus cursus, ac sollicitudin mi pulvinar. Praesent vulputate lorem lectus, eu pulvinar mi finibus eu. Donec sit amet massa massa. Sed eget molestie quam. Praesent iaculis diam ut nunc condimentum auctor. Morbi porttitor varius ante vel venenatis. Integer ut porttitor elit, id lacinia turpis. Integer id iaculis ante, placerat ornare tortor. Sed et consectetur est, sit amet rhoncus turpis.

\begin{table}[h!] 
\begin{center}
\begin{threeparttable}
\caption{Testing Table Numbering}\label{test-table1}\vspace{1.5ex} 
{\small \begin{tabular}{c|c|c}\thickhline 
\rowcolor[gray]{.9}   {\sffamily\bfseries Compiler} & {\sffamily\bfseries Compiler-Specific} &
{\sffamily\bfseries Standard-Conforming} \\ \hline 
GCC &\texttt{\_\_attribute\_\_((pure))} & \texttt{[[gnu::pure]]} \\ \hline
Clang & \texttt{\_\_attribute\_\_((no\_sanitize))} &\texttt{[[clang::no\_sanitize]]} \\ \hline 
MSVC & \texttt{declspec(deprecated)} & \texttt{[[deprecated]]} \\ \thickhline
\end{tabular}
}
\end{threeparttable} 
\end{center}
\end{table}

Phasellus ac tellus ipsum. Donec ornare id urna a blandit. Pellentesque finibus nulla augue, vitae pretium justo imperdiet ac. Aenean condimentum nibh eget turpis viverra gravida. Nullam diam lectus, egestas id ex quis, blandit facilisis libero. Proin laoreet ante dictum tristique finibus. Vestibulum pretium varius ipsum, sed blandit ex venenatis a. Nam pharetra pretium accumsan. Integer accumsan purus elementum tortor aliquam, commodo semper leo commodo. Nunc finibus varius erat, non hendrerit leo dapibus at. Donec mattis porta ex, eu ornare nibh condimentum id. Sed sit amet erat sit amet urna volutpat ullamcorper. Aenean feugiat eget orci ac feugiat. Nulla facilisi. Cras augue lacus, placerat sit amet vestibulum at, mollis sit amet nisl.

Nam id sem sagittis, placerat justo eget, eleifend justo. Sed ultrices rhoncus quam ut auctor. Nam pellentesque risus orci, a rhoncus arcu lobortis in. Vestibulum tristique nisi sed sollicitudin aliquam. Praesent ut odio sapien. Ut pretium sollicitudin nisi sit amet blandit. Cras ut eleifend elit. Aliquam vel tincidunt lacus, non sodales justo. Sed euismod sapien non diam euismod, sit amet vulputate nisi rhoncus. Vestibulum leo diam, dapibus suscipit blandit molestie, suscipit eget ante. Donec tristique rhoncus purus, nec pellentesque lorem volutpat non. Quisque purus dui, egestas ut ultricies sed, pharetra sit amet turpis. Praesent ut ligula porttitor, convallis velit sit amet, rhoncus urna. 






  % Chapter 1
\cleardoublepage

\chapter[Conditionally Safe Features]{Conditionally Safe Features}\label{ch-conditional}

%%%%%%%%%%%%%%% C++11
\stepcounter{cppxx}
\addtocontents{toc}{\protect\renewcommand*\protect\cftsecindent{1.25em}}
\cftaddnumtitleline{toc}{section}{\thechapter.\thecppxx}{C++11}{\thepage}
% to add to the TOC at the section level
\renewcommand{\cppxx}{C++11}

\newpage
\addtocontents{toc}{\protect\renewcommand*\protect\cftsecindent{4.5em}}
\section[{\tt alignas}]{The {\SecCode alignas} Decorator}\label{alignas}
% 8 Feb 2021 JMB, file comment cleanup

\emcppsFeature{
    short={\lstinline!alignas!},
    tocshort={\TOCCode alignas},
    long={The {\SecCode alignas} Decorator},
    toclong={The \lstinline!alignas! Decorator},
    rhshort={\RHCode alignas},
}{alignas}
%\section[{\tt alignas}]{The {\SecCode alignas} Decorator\sectionmark{\RHCode alignas}}\label{alignas}\sectionmark{\RHCode alignas}
\setcounter{table}{0}
\setcounter{footnote}{0}
\setcounter{lstlisting}{0}

\texttt{alignas}, a keyword that acts like an attribute, is used to widen (make more strict)
the \textbf{alignment} of a \textbf{variable}, \textbf{user-defined
type}, or \textbf{data member}.

\subsection[Description]{Description}\label{description}

The \texttt{alignas} specifier provides a means of further restricting
the granularity at which (1) a particular object of arbitrary type, (2)
a user-defined type (\texttt{class}, \texttt{struct}, \texttt{union}, or
\texttt{enum}), or (3) an individual data member is permitted to reside
within the virtual-memory-address space.

\subsubsection[Restricting the alignment of a particular object]{Restricting the alignment of a particular object}\label{restricting-the-alignment-of-a-particular-object}

In its most basic form, \texttt{alignas} acts like an
attribute that accepts (as an argument) an
\textbf{integral constant expression} representing an explicitly
supplied minimum alignment value:

\begin{lstlisting}[language=C++]
alignas(64) int i;    // OK, (ù{\codeincomments{i}}ù) is aligned on a 64-byte address boundary.
int j alignas(8), k;  // OK, (ù{\codeincomments{j}}ù) is 8-byte aligned; (ù{\codeincomments{k}}ù) remains naturally aligned.
\end{lstlisting}

\noindent If more than one alignment pertains to a given object, the most
restrictive alignment value is applied:

\begin{lstlisting}[language=C++]
alignas(4) alignas(8) alignas(2) char m;  // OK, (ù{\codeincomments{m}}ù) is 8-byte aligned.
alignas(8) int n alignas(16);             // OK, (ù{\codeincomments{n}}ù) is 16-byte aligned.
\end{lstlisting}

\noindent For a program to be \textbf{well formed}, a specified alignment value
must satisfy several\linebreak[4] \mbox{requirements}:

\begin{enumerate}
\item{Be either zero or a non-negative integral power of two of type \texttt{std::size\_t} (0, 1, 2, 4, 8, 16\dots).}
\item{Be at least the minimum alignment\cprotect\footnote{The minimum alignment of an entity is the least restrictive memory-address boundary at which the entity can be placed and have the program continue to work properly. This value is platform dependent and often subject to compiler controls but, by default, is often well approximated by \textbf{natural alignment}; see \textit{\titleref{alignas-appendix}: \titleref{natural-alignment}} on page \pageref{natural-alignment}.} required by the decorated entity.}
\item{Be no more than the largest alignment\cprotect\footnote{The notion of the largest supported alignment is characterized by both \textbf{maximal alignment} and the maximum \textbf{extended alignment}. \textbf{Maximal alignment} is defined as that most restrictive alignment that is valid in \emph{all} contexts on the current platform. All fundamental and pointer types necessarily have a minimal alignment that is less than or equal to \texttt{alignof(std::max\_align\_t)} — typically 8 or 16. Any alignment value greater than \textbf{maximal alignment} is an \textbf{extended alignment} value. Whether any extended alignment is supported (and in which contexts) is implementation defined. On typical platforms, extended alignment will often be as large as $2^{18}$ or $2^{19}$, however implementations may warn when the alignment of a global object exceeds some maximal hardware threshold (such as the size of a physical memory page, e.g., 4096 or 8192). For \textbf{automatic variables} (defined on the program stack), making alignment more restrictive than what would naturally be employed  is seldom desired because at most one thread is able to access proximately located variables there unless explicitly passed in via address to separate threads; see \textit{\titleref{alignas-use-cases}: \titleref{avoiding-false-sharing-among-distinct-objects-in-a-multi-threaded-program}} on page~\pageref{avoiding-false-sharing-among-distinct-objects-in-a-multi-threaded-program}. Note that, in the case of \texttt{i} in the first code snippet on page~\pageref{restricting-the-alignment-of-a-particular-object}, a conforming platform that did not support an extended alignment of 64 would be required to report an error at compile time.} supported on the platform in the context in which the entity appears.}
\end{enumerate}

\noindent Additionally, if the specified alignment value is zero, the
\texttt{alignas} specifier is ignored:

\begin{lstlisting}[language=C++]
// Static variables declared at namespace scope
alignas(32) int i0;  // OK, aligned on a 32-byte boundary (extended alignment)
alignas(16) int i1;  // OK, aligned on a 16-byte boundary (maximum alignment)
alignas(8)  int i2;  // OK, aligned on an 8-byte boundary
alignas(7)  int i3;  // error: not a power of two
alignas(4)  int i4;  // OK, no change to alignment boundary
alignas(2)  int i5;  // error: less than minimum alignment on this platform
alignas(0)  int i6;  // OK, (ù{\codeincomments{alignas}}ù) specifier ignored

alignas(1024 * 16) int i7;
    // OK, might warn: e.g., exceeds (physical) page size on current platform

alignas(1024 * 1024 * 512) int i8;
    // (likely) compile-time error: e.g., exceeds maximum size of object file

alignas(8) char buf[128];  // create 8-byte-aligned, 128-byte character buffer

void f()
{
    // automatic variables declared at function scope
    alignas(4)  double e0;  // error: less than minimum alignment on this platform
    alignas(8)  double e1;  // OK, no-change to (8-byte) alignment boundary
    alignas(16) double e2;  // OK, aligned to maximum (fundamental) alignment value
    alignas(32) double e3;  // OK, maximum alignment value exceeded; might warn
}
\end{lstlisting}

\subsubsection[Restricting the alignment of a user-defined type]{Restricting the alignment of a user-defined type}\label{restricting-the-alignment-of-a-user-defined-type}

The \texttt{alignas} specifier can also be used to specify alignment for
user-defined types (UDTs), such as a \texttt{class}, \texttt{struct},
\texttt{union}, or \texttt{enum}. When specifying the alignment of a UDT,
 the \texttt{alignas} keyword is placed \emph{after} the
type specifier (e.g., \texttt{class}) and just before the name of the
type (e.g., \texttt{C}):

\begin{lstlisting}[language=C++]
class alignas(2)  C { };  // OK, aligned on a 2-byte boundary; size = 2
struct alignas(4) S { };  // OK, aligned on a 4-byte boundary; size = 4
union alignas(8)  U { };  // OK, aligned on an 8-byte boundary; size = 8
enum alignas(16)  E { };  // OK, aligned on a 16-byte boundary; size = 4
\end{lstlisting}

\noindent Notice that, for each of \texttt{class}, \texttt{struct}, and
\texttt{union} above, the \texttt{sizeof} objects of that type increased
to match the alignment; in the case of the \texttt{enum}, however, the
size remains that of the default \textbf{underlying type} (e.g., 4
bytes) on the current platform.{\cprotect\footnote{When \texttt{alignas}
is applied to an enumeration \texttt{E}, the Standard does not
indicate whether padding bits are added to \texttt{E}'s object
representation or not, affecting the result of \texttt{sizeof(E)}. The
implementation variance resulting from this lack of clarity in the
  Standard was captured in \textbf{miller17}. The outcome of the core
  issue was to completely remove permission for \texttt{alignas} to be
  applied to enumerations (see \textbf{iso18a}). Therefore, conforming implementations will
  eventually stop accepting the \texttt{alignas} specifier on
  enumerations in the future.}}

Again, specifying an alignment that is less than what would occur
naturally or else is restricted explicitly is ill formed:

\begin{lstlisting}[language=C++]
struct alignas(2) T0 { int i; };
    // Error: Alignment of (ù{\codeincomments{T0}}ù) (2) is less than that of (ù{\codeincomments{int}}ù) (4).
struct alignas(1) T1 { C c; };
    // Error: Alignment of (ù{\codeincomments{T1}}ù) (1) is less than that of (ù{\codeincomments{C}}ù) (2).
\end{lstlisting}


\subsubsection[Restricting the alignment of individual data members]{Restricting the alignment of individual data members}\label{restricting-the-alignment-of-individual-data-members}

Within a user-defined type (\texttt{class}, \texttt{struct}, or
\texttt{union}), using the attribute-like syntax of the
\texttt{alignas} keyword to specify the alignments of individual data
members is possible:

\begin{lstlisting}[language=C++]
struct T2
{
    alignas(8)  char   x;  // size   1; alignment 8
    alignas(16) int    y;  // size   4; alignment 16
    alignas(64) double y;  // size   8; alignment 64
};  // size 128; alignment 64
\end{lstlisting}

\noindent The effect here is the same as if we had added the padding explicitly
and then set the alignment of the structure overall:

\begin{lstlisting}[language=C++]
struct alignas(64) T3
{
    char   x;      // size   1; alignment 8
    char   a[15];  // padding
    int    y;      // size   4; alignment 16
    char   b[44];  // padding
    double z;      // size   8; alignment 64
    char   c[56];  // padding (optional)
};  // size 128; alignment 64
\end{lstlisting}

\noindent Again, if more than one attribute pertains to a given data member,
the maximum applicable alignment value is applied:

\begin{lstlisting}[language=C++]
struct T4
{
    alignas(2) char
        c1 alignas(1),  // size 1; alignment 2
        c2 alignas(2),  // size 1; alignment 2
        c4 alignas(4);  // size 1; alignment 4
};                      // size 8; alignment 4
\end{lstlisting}


\subsubsection[Matching the alignment of another type]{Matching the alignment of another type}\label{matching-the-alignment-of-another-type}

The \texttt{alignas} specifier also accepts (as an argument) a type
identifier. In its alternate form, \texttt{alignas(T)} is strictly
equivalent to \texttt{alignas(alignof(T))}:

\begin{lstlisting}[language=C++]
alignas(int) char c;  // equivalent to (ù{\codeincomments{alignas(alignof(int)) char c;}}ù)
\end{lstlisting}


\subsection[Use Cases]{Use Cases}\label{alignas-use-cases}

\subsubsection[Creating a sufficiently aligned object buffer]{Creating a sufficiently aligned object buffer}\label{creating-a-sufficiently-aligned-object-buffer}

When writing low-level, system-infrastructure code, constructing an object within a raw buffer is sometimes useful. As a minimal
example, consider a function that uses a local character buffer to
create an object of type \texttt{std::complex<long}~\texttt{double>} on
the program stack using placement \texttt{new}:

\begin{lstlisting}[language=C++]
void f()
{
    // ...
    char objectBuffer[sizeof(std::complex<long double>)];  // BAD IDEA
    // ...
    new(objectBuffer) std::complex<long double>(1.0, 0.0);  // Might dump core!
    // ...
}
\end{lstlisting}

\noindent The essential problem with the code above is that \texttt{objectBuffer},
being an array of characters (each having an alignment of 1), is itself
byte aligned. The compiler is therefore free to place it on any address
boundary. On the other hand, \texttt{std::complex<long}~\texttt{double>} is an aggregate consisting of two \texttt{long}~\texttt{double} objects
and therefore necessarily requires (at least) the same strict alignment
(typically 16) as the two \texttt{long}~\texttt{double} objects it comprises. Previous
solutions to this problem involved creating a \texttt{union} of the
object buffer and some maximally aligned type (e.g.,
\texttt{std::max\_align\_t}):

\begin{lstlisting}[language=C++]
#include <cstddef>  // (ù{\codeincomments{std::max\_align\_t}}ù)

void f()
{
    // ...

    union {                                              // awkward workaround
        std::max_align_t dummy;  // (ù{\codeincomments{typedef}}ù) to maximally aligned type
        char objectBuffer[sizeof(std::complex<long double>)];
    } objectBuffer;

    // ...

    new(&objectBuffer) std::complex<long double>(1.0, 0.0);  // OK

    // ...
}
\end{lstlisting}

\noindent Using the alternate syntax for \texttt{alignas}, we can avoid gratuitous
complexity and just state our intentions explicitly:

\begin{lstlisting}[language=C++]
void f()
{
    // ...

    alignas(std::complex<long double>) char objectBuffer[
                              sizeof(std::complex<long double>)];  // GOOD IDEA

    // ...

    new(objectBuffer) std::complex<long double>(1.0, 0.0);  // OK

    // ...
}
\end{lstlisting}


\subsubsection[Ensuring proper alignment for architecture-specific instructions]{Ensuring proper alignment for architecture-specific instructions}\label{ensuring-proper-alignment-for-architecture-specific-instructions}

Architecture-specific instructions or compiler intrinsics might require
the data they act on to have a specific alignment. One example of such
intrinsics is the \emph{Streaming SIMD Extensions (SSE)}\footnote{\textbf{inteliig}, ``Technologies: SSE"} instruction set available on the x86
architecture. SSE instructions operate on groups of four 32-bit
single-precision floating-point numbers at a time, which are required to
be 16-byte aligned.{\cprotect\footnote{``Data must be
16-byte aligned when loading to and storing from the 128-bit XMM
  registers used by SSE/SSE2/SSE3/SSSE3'': see \textbf{{intel16}},
  section 4.4.4, ``Data Alignment for 128-Bit Data," pp. 4-19--4-20.}} The
\texttt{alignas} specifier can be used to create a type satisfying this requirement:

\begin{lstlisting}[language=C++]
struct SSEVector
{
    alignas(16) float d_data[4];
};
\end{lstlisting}

\noindent Each object of the \texttt{SSEVector} type above is guaranteed always to
be aligned to a 16-byte boundary and can therefore be safely (and
conveniently) used with SSE intrinsics:

\begin{lstlisting}[language=C++]
#include <xmmintrin.h>  // (ù{\codeincomments{\_\_m128}}ù) and (ù{\codeincomments{\_mm\_XXX}}ù) functions

void f()
{
    const SSEVector v0 = {0.0f, 1.0f, 2.0f, 3.0f};
    const SSEVector v1 = {10.0f, 10.0f, 10.0f, 10.0f};

    __m128 sseV0 = _mm_load_ps(v0.d_data);
    __m128 sseV1 = _mm_load_ps(v1.d_data);
        // (ù{\codeincomments{\_mm\_load\_ps}}ù) requires the given (ù{\codeincomments{float}}ù) array to be 16-byte aligned.
        // The data is loaded into a dedicated 128-bit CPU register.

    __m128 sseResult = _mm_add_ps(sseV0, sseV1);
        // sum two 128-bit registers; typically generates an (ù{\codeincomments{addps}}ù) instruction

    SSEVector vResult;
    _mm_store_ps(vResult.d_data, sseResult);
        // Store the result of the sum back into a (ù{\codeincomments{float}}ù) array.

    assert(vResult.d_data[0] == 10.0f);
    assert(vResult.d_data[1] == 11.0f);
    assert(vResult.d_data[2] == 12.0f);
    assert(vResult.d_data[3] == 13.0f);
}
\end{lstlisting}


\subsubsection[Avoiding \textbf{false sharing} among distinct objects in a multi-threaded program]{Avoiding \textbf{false sharing} among distinct objects in a multi-threaded program}\label{avoiding-false-sharing-among-distinct-objects-in-a-multi-threaded-program}

In the context of an application where multithreading has been employed
to improve performance, seeing
a previously single-threaded workflow become even less performant after
a parallelization attempt can be surprising (and disheartening). One possible insidious cause of such
disappointing results comes from \textbf{false sharing} --- a situation
in which multiple threads unwittingly harm each other's performance
while writing to logically independent variables that happen to reside
on the same \textbf{cache line}; see {\it \titleref{alignas-appendix}:} {\it\titleref{cache-lines,-l1,-l2,-and-l3-cache,-pages,-and-virtual-memory}} on page~\pageref{cache-lines,-l1,-l2,-and-l3-cache,-pages,-and-virtual-memory}.

As a simple (purely pedagogical) illustration of the potential
performance degradation resulting from \textbf{false sharing}, consider
a function that spawns separate threads to repeatedly increment
(concurrently) logically distinct variables that happen to reside in
close proximity on the program stack:

\begin{lstlisting}[language=C++]
#include <thread>  // (ù{\codeincomments{std::thread}}ù)

volatile int target = 0;  // updated asynchronously from multiple threads

void incrementJob(int* p);
    // Repeatedly increment (ù{\codeincomments{*p}}ù) a large, fixed number of times;
    // periodically write its current value to (ù{\codeincomments{target}}ù).

void f()
{
    int i0 = 0;  // Here, (ù{\codeincomments{i0}}ù) and (ù{\codeincomments{i1}}ù) likely share the same cache line,
    int i1 = 0;  // i.e., byte-aligned memory block on the program stack.

    std::thread t0(&incrementJob, &i0);
    std::thread t1(&incrementJob, &i1);
        // Spawn two parallel jobs incrementing the respective variables.

    t0.join();
    t1.join();
        // Wait for both jobs to be completed.
}
\end{lstlisting}

\noindent In the simplistic example above, the proximity in memory between
\texttt{i0} and \texttt{i1} can result in their belonging to the same
\textbf{cache line}, thus leading to \textbf{false sharing}. By
prepending \texttt{alignas(64)} to the declaration of both integers, we
ensure that the two variables reside on distinct cache lines:

\begin{lstlisting}[language=C++]
// ...

void f()
{
    alignas(64) int i0 = 0;   // Assuming a cache line on this platform is 64
    alignas(64) int i1 = 0;   // bytes, (ù{\codeincomments{i0}}ù) and (ù{\codeincomments{i1}}ù) will be on separate ones.

    // ...
\end{lstlisting}

\noindent As an empirical demonstration of the effects of \textbf{false sharing},
a benchmark program repeatedly calling \texttt{f} completed its
execution seven times faster on average when compared to the same
program without use of \texttt{alignas}.{\cprotect\footnote{The benchmark
program was compiled using Clang 11.0.0 using \texttt{-Ofast},
\texttt{-march=native}, and \mbox{\texttt{-std=c++11}}. The program was then
executed on a machine running Windows 10 x64, equipped with an Intel
Core i7-9700k CPU (8 cores, 64-byte cache line size). Over the
course of multiple runs, the version of the benchmark without
\texttt{alignas} took 18.5967ms to complete (on average), while the
version with \texttt{alignas} took 2.45333ms to complete (on average).
  See \textbf{{[PRODUCTION: CODE PROVIDED WITH BOOK] alignasbenchmark}} for the source code of the program.}}

\subsubsection[Avoiding \textbf{false sharing} within a single thread-aware object]{Avoiding \textbf{false sharing} within a single thread-aware object}\label{avoiding-false-sharing-within-a-single-thread-aware-object}

A real-world scenario where the need for preventing \textbf{false
sharing} is fundamental occurs in the implementation of high-performance
concurrent data structures. As an example, a thread-safe ring buffer
might make use of \texttt{alignas} to ensure that the indices of the
head and tail of the buffer are aligned at the start of a cache line
(typically 64, 128, or 256 bytes), thereby preventing them from
occupying the same one.

\begin{lstlisting}[language=C++]
class ThreadSafeRingBuffer
{
    alignas(cpuCacheSize) std::atomic<std::size_t> d_head;
    alignas(cpuCacheSize) std::atomic<std::size_t> d_tail;

    // ...
};
\end{lstlisting}

\noindent Not aligning \texttt{d\_head} and \texttt{d\_tail} (above) to the CPU
cache size might result in poor performance of the
\texttt{ThreadSafeRingBuffer} because CPU cores that need to access only one
of the variables will inadvertently load the other one as well,
triggering expensive hardware-level coherency mechanisms between the
cores' caches. On the other hand, specifying such substantially stricter
alignment on consecutive data members necessarily increases the size of
the object; see {\it\titleref{potential-pitfalls}:} {\it\titleref{stricter-alignment-might-reduce-cache-utilization}} on page~\pageref{stricter-alignment-might-reduce-cache-utilization}.

\subsection[Potential Pitfalls]{Potential Pitfalls}\label{potential-pitfalls}

\subsubsection[Underspecifying alignment is not universally reported]{Underspecifying alignment is not universally reported}\label{underspecifying-alignment-is-not-universally-reported}

The Standard is clear when it comes to underspecifying
alignment{\cprotect\footnote{\textbf{{cpp11}}, section 7.6.2, ``Alignment Specifier," paragraph~5, pp. 179}}:
\begin{quote}
The combined effect of all \textit{alignment-specifiers} in a declaration
shall not specify an alignment that is less strict than the alignment
that would be required for the entity being declared if all
\textit{alignment-specifiers} were omitted (including those in other
declarations).
\end{quote}
The compiler is required to honor the
specified value if it is a \textbf{fundamental
alignment},{\cprotect\footnote{``If the constant expression evaluates to
a fundamental alignment, the alignment requirement of the declared
entity shall be the specified fundamental alignment'': \textbf{cpp11}, section~7.6.2, ``Alignment Specifier," paragraph~2, item~2, p. 178.}} so
imagining how this would lead to anything other than an ill-formed
program is difficult:

\begin{lstlisting}[language=C++]
alignas(4) void* p;              // (1) Error: (ù{\codeincomments{alignas(4)}}ù) is below minimum, 8.

struct alignas(2) S { int x; };  // (2) Error: (ù{\codeincomments{alignas(2)}}ù) is below minimum, 4.

struct alignas(2) T { };
struct alignas(1) U { T e; };    // (3) Error: (ù{\codeincomments{alignas(1)}}ù) is below minimum, 2.
\end{lstlisting}

\noindent Each of the three errors above are reported by Clang, but GCC
doesn't issue so much as a warning (let alone the required error) ---
even in the most pedantic warning mode. Thus, one could write a program,
involving statements like those above, that happens to work on one
platform (e.g., GCC) but fails to compile on another (e.g.,
Clang).{\cprotect\footnote{Underspecifying alignment is not reported at
all by GCC 10.1, using the
\texttt{-std=c++11} \texttt{-Wall} \texttt{-Wextra} \mbox{\texttt{-Wpedantic}}
flags. With the same set of options, Clang 10.0 produces a compilation
failure. MSVC v19.24 will produce a warning and ignore any alignment
  less than the minimum one.}}

\subsubsection[Incompatibly specifying alignment is \textbf{IFNDR}]{Incompatibly specifying alignment is \textbf{IFNDR}}\label{incompatibly-specifying-alignment-is-ifndr}

It is permissible to forward declare a user-defined type (UDT)
without an \texttt{alignas} specifier so long as all defining
declarations of the type have either no \texttt{alignas} specifier or
have the same one. Similarly, if any forward declaration of a
user-defined type has an \texttt{alignas} specifier, then all defining
declarations of the type must have the same specifier and that specifier
must be \emph{equivalent to} (not necessarily \emph{the same as}) that
in the forward declaration:

\begin{lstlisting}[language=C++]
struct Foo;                  // OK, does not specify an alignment
struct alignas(double) Foo;  // OK, must be equivalent to every definition
struct alignas(8) Foo;       // OK, all definitions must be identical.
struct alignas(8) Foo { };   // OK, equivalent to each decl. specifying (ù{\codeincomments{alignas}}ù)
struct Foo;                  // OK, has no effect
struct alignas(8) Foo;       // OK, has no effect; might warn after definition
\end{lstlisting}

\noindent Specifying an alignment in a forward declaration without specifying an
equivalent one in the defining declaration is \textbf{ill formed; no diagnostic
is required (IFNDR)} if the two declarations appear in distinct translation
units:

\begin{lstlisting}[language=C++]
struct alignas(4) Bar;      // OK, forward declaration
struct Bar { };             // error: missing (ù{\codeincomments{alignas}}ù) specifier

struct alignas(4) Baz;      // OK, forward declaration
struct alignas(8) Baz { };  // error: non-equivalent (ù{\codeincomments{alignas}}ù) specifier
\end{lstlisting}

\noindent Both of the errors above are flagged by Clang, but neither of them is
reported by GCC. Note that when the inconsistency occurs across
translation units, no mainstream compiler is likely to diagnose the
problem:

\begin{lstlisting}[language=C++]
// file1.cpp:
struct Bam { char ch; } bam, *p = &bam;

// file2.cpp:
struct alignas(int) Bam;  // Error: definition of (ù{\codeincomments{Bam}}ù) lacks alignment specifier.
extern Bam* p;            //        (no diagnostic required)
\end{lstlisting}

\noindent Any program incorporating both translation units above is
\textbf{ill formed, no diagnostic required}.

\subsubsection[Stricter alignment might reduce cache utilization]{Stricter alignment might reduce cache utilization}\label{stricter-alignment-might-reduce-cache-utilization}

User-defined types having artificially stricter alignments than would
naturally occur on the host platform means that fewer of them can fit
within any given level of physical cache within the hardware. Types
having data members whose alignment is artificially widened tend to be
larger and thus suffer the same lost cache utilization. As an
alternative to enforcing stricter alignment to avoid \textbf{false
sharing}, consider organizing a multithreaded program such that tight
clusters of repeatedly accessed objects are always acted upon by only a
single thread at a time, e.g., using local (arena) memory allocators;
see {\it\titleref{alignas-appendix}:} {\it\titleref{cache-lines,-l1,-l2,-and-l3-cache,-pages,-and-virtual-memory}} on page~\pageref{cache-lines,-l1,-l2,-and-l3-cache,-pages,-and-virtual-memory}.

\subsection[See Also]{See Also}\label{see-also}

\begin{itemize}
\item{Section~\ref{alignof}, ``\titleref{alignof}" — Safe C++11 feature that inspects the alignment of a given type}
\item{Section~\ref{attributes}, ``\titleref{attributes}" — Safe C++11 feature that shows how other attributes (following the conventional attribute notation) are used to annotate source code, improve error diagnostics, and implicitly code generation} %%% WORDS MISSING!
\end{itemize}

\subsection[Further Reading]{Further Reading}\label{further-reading}

None so far

\subsection[Appendix]{Appendix}\label{alignas-appendix}

\subsubsection[Natural Alignment]{Natural Alignment}\label{natural-alignment}

By default, fundamental, pointer, and enumerated types typically reside
on an address boundary that divides the size of the object; we refer to
such alignment as \textbf{natural alignment}{\cprotect\footnote{Sizes
and alignment shown here are typical but not specifically required by
the standard. On some platforms, one can request that all
types be \textbf{byte aligned}. While such a representation is more
compact, entities that span memory boundaries can require multiple
fetch operations leading to run times that are typically
significantly (sometimes as much as an order of magnitude) slower when
  run in this ``packed'' mode.}}:

\begin{lstlisting}[language=C++]
char   c;  // size 1; alignment 1; boundaries: 0x00, 0x01, 0x02, 0x03, ...
short  s;  // size 2; alignment 2: boundaries: 0x00, 0x02, 0x04, 0x06, ...
int    i;  // size 4; alignment 4; boundaries: 0x00, 0x04, 0x08, 0x0c, ...
float  f;  // size 4; alignment 4; boundaries: 0x00, 0x04, 0x08, 0x0c, ...
double d;  // size 8; alignment 8; boundaries: 0x00, 0x08, 0x10, 0x18, ...
\end{lstlisting}

\noindent For aggregates (including arrays) or user-defined types, the alignment
is typically that of the most strictly aligned subelement:

\begin{lstlisting}[language=C++]
struct S0
{
    char a;  // size 1; alignment 1
    char b;  // size 1; alignment 1
    int  c;  // size 4; alignment 4
};           // size 8; alignment 4

struct S1
{
    char a;  // size  1; alignment 1
    int  b;  // size  4; alignment 4
    char c;  // size  1; alignment 1
};           // size 12; alignment 4

struct S2
{
    int  a;  // size 4; alignment 4
    char b;  // size 1; alignment 1
    char c;  // size 1; alignment 1
};           // size 8; alignment 4

struct S3
{
    char a;  // size 1; alignment 1
    char b;  // size 1; alignment 1
};           // size 2; alignment 1

struct S4
{
    char a[2];  // size 2; alignment 1
};              // size 2; alignment 1
\end{lstlisting}

\noindent Size and alignment behave similarly with respect to \textbf{structural
inheritance}:

\begin{lstlisting}[language=C++]
struct D0 : S0
{
    double d;  // size  8; alignment 8
};             // size 16; alignment 8

struct D1 : S1
{
    double d;  // size  8; alignment 8
};             // size 24; alignment 8

struct D2 : S2
{
    int d;  // size  4; alignment 4
};          // size 12; alignment 4

struct D3 : S3
{
    int d;  // size 4; alignment 4
};          // size 8; alignment 4

struct D4 : S4
{
    char d;  // size 1; alignment 1
};           // size 3; alignment 1
\end{lstlisting}

\noindent Finally, virtual functions invariably introduce an implicit
virtual-table-pointer member having a size and alignment corresponding
to that of a memory address (e.g., 4 or 8) on the host platform:

\begin{lstlisting}[language=C++]
struct S5
{
    virtual ~S5();
};                   // size 8; alignment 8

struct D5 : S5
{
    char d;  // size  1; alignment 1
};           // size 16; alignment 8
\end{lstlisting}


\subsubsection[Cache lines; L1, L2, and L3 cache; pages; and virtual memory]{Cache lines; L1, L2, and L3 cache; pages; and virtual memory}\label{cache-lines,-l1,-l2,-and-l3-cache,-pages,-and-virtual-memory}

Modern computers are highly complex systems, and a detailed
understanding of their intricacies is unnecessary to achieve most of
the performance benefits. Still, certain general themes and
rough thresholds aid in understanding how to squeeze just a bit
more out of the underlying hardware. In this section, we sketch
fundamental concepts that are common to all modern computer hardware;
although the precise details will vary, the general ideas remain
essentially the same.

In its most basic form, a computer consists of central processing unit
(CPU) having internal registers that access main memory (MM).
Registers in the CPU (on the order of hundreds of bytes) are among the
fastest forms of memory, while main memory (typically many gigabytes) is
orders of magnitude slower. An almost universally observed phenomenon is
that of \textbf{locality of reference}, which suggests that data that
resides in close proximity (in the virtual address space) is more likely
to be accessed together in rapid succession than more distant data.

To exploit the phenomenon of \textbf{locality of reference}, computers
introduce the notion of a cache that, while much faster than main
memory, is also much smaller. Programs that attempt to amplify
\textbf{locality of reference} will, in turn, often be rewarded with
faster run times. The organization of a cache and, in fact, the number of
levels of cache (e.g., L1, L2, L3,$\ldots$) will vary, but the basic
design parameters are, again, more or less the same. A given level of
cache will have a certain total size in bytes (invariably an integral
power of two). The cache will be segmented into what are called
\textbf{cache lines} whose size (a smaller power of two) divides that of
the cache itself. When the CPU accesses main memory, it first looks to
see if that memory is in the cache; if it is, the value is returned
quickly (known as a \textbf{cache hit}). Otherwise, the cache line(s)
containing that data is (are) fetched (from the next higher level of cache
or from main memory) and placed into the cache (known as a \textbf{cache
miss}), possibly ejecting other less recently used
ones.{\cprotect\footnote{Conceptually, the cache is often thought of as
being able to hold any arbitrary subset of the most recently accessed
cache lines. This kind of cache is known as \textbf{fully
associative}. Although it provides the best hit rate, a \textbf{fully
associative} cache requires the most power along with significant
additional chip area to perform the fully parallel lookup. \textbf{Direct-mapped} cache associativity is at the
other extreme. In direct mapped, each memory location has exactly one location
available to it in the cache. If another memory location mapping to
that location is needed, the current cache line must be flushed from
the cache. Although this approach has the lowest hit rate, lookup
times, chip area, and power consumption are all minimized (optimally).
Between these two extremes is a continuum that is referred to as
\textbf{set associative}. A \textbf{set associate} cache has more than
  one (typically 2, 4, or 8; see \textbf{solihin15}, section~5.2.1, ``Placement Policy," pp. 136--141, and \textbf{hruska20})
  location in which each memory location in main memory can reside.
  Note that, even with a relatively small $N$, as $N$ increases, an $N$-way
  \textbf{set associative} cache quickly approaches the hit rate of a fully
  associative cache at greatly reduced collateral cost; for most
  software-design purposes, any loss in hit rate due to set
  associativity of a cache can be safely ignored.}}

Data residing in distinct cache lines is physically independent and can
be written concurrently by multiple threads. Logically unrelated data
residing in the same cache line, however, is nonetheless physically
coupled; two threads that write to such logically unrelated data will find
themselves synchronized by the hardware. Such unexpected and typically
undesirable sharing of a cache line by unrelated data acted upon by two
concurrent threads is known as \textbf{false sharing}. One way of
avoiding \textbf{false sharing} is to align such data on a cache-line
boundary, thus rendering accidental collocation of such data on the same
cache line impossible. Another (more broad-based) design approach that
avoids lowering cache utilization is to ensure that data acted upon by a
given thread is kept physically separate --- e.g., through the use of
local (arena) memory allocators.{\cprotect\footnote{\textbf{{lakos17}},
  \textbf{{lakos19}}, \textbf{{lakos22}}}}

Finally, even data that is not currently in cache but resides nearby in
MM can benefit from locality. The virtual address space, synonymous with
the size of a \texttt{void*} (typically 64-bits on modern general-purpose hardware), has historically well exceeded the physical memory
available to the CPU. The operating system must therefore maintain a
mapping (in main memory) from what is resident in physical memory and
what resides in secondary storage (e.g., on disc). In addition,
essentially all modern hardware provides a
\textbf{TLB}{\cprotect\footnote{A translation-lookaside buffer (TLB) is
a kind of address-translation cache that is typically part of a chip's
memory management unit (MMU). A TLB holds a recently accessed subset
of the complete mapping (itself maintained in MM) from virtual memory
address to physical ones. A TLB is used to reduce access time when the
requisite pages are already resident in memory; its size (e.g., 4K) is
capped at the number of bytes of physical memory (e.g., 32Gb) divided
by the number of bytes in each physical page (e.g., 8Kb), but could be
smaller. Because it resides on chip, is typically an order of
magnitude faster (SRAM versus DRAM), and requires only a single lookup
(as opposed to two or more when going out to MM), there is an enormous
  premium on minimizing TLB misses.}} that caches the addresses of the
most recently accessed physical pages, providing yet another advantage
to having the \textbf{working set} (i.e., the current set of frequently
accessed pages) remain small and densely packed with relevant
data.{\cprotect\footnote{Note that memory for handle-body
types (e.g., \texttt{std::vector} or \texttt{std::deque}) and
especially node-based containers (e.g., \texttt{std::map} and
\texttt{std::unordered\_map}), originally allocated within a single
page, can --- through deallocation and reallocation (or even move
operations) --- become scattered across multiple (perhaps many)
 pages, thus causing what was originally a relatively small \textbf{working set}
to no longer fit within physical memory. This phenomenon, known as
\textbf{diffusion} (which is a distinct concept from
\textbf{fragmentation}), is what typically leads to a substantial
runtime performance degradation (due to \textbf{thrashing}) in large,
long-running programs. Such \textbf{diffusion} can be mitigated by
judicious use of local arena memory allocators (and deliberate
avoidance of \textbf{move operations} across disparate localities of
  frequent memory usage).}} What's more, dense working sets, in addition
to facilitating hits for repeat access, increase the likelihood that
data that is coresident on a page (or cache line) will be needed soon
(i.e., in effect acting as a prefetch).{\cprotect\footnote{We sometimes
lightheartedly refer to the beneficial prefetch of unrelated data that
is accidentally needed subsequently (e.g., within a single thread) due
to high locality within a cache line (or a physical page) as
  \textbf{true sharing}.}} Table~\ref{table-alignas-appendix} provides a summary of typical physical parameters found in modern computers today.

%%%%%%TODO: TBD: @Lori: Format this as a figure with a caption. @jslakos: Double-check accuracy of values.

\begin{table}[h!]
\begin{center}
\begin{threeparttable}
\caption{Various sizes and access speeds of typical memory for modern \mbox{computers}}\label{table-alignas-appendix} \vspace{1.5ex}
{\small \begin{tabular}{c|c|c}\thickhline
\rowcolor[gray]{.9}    {\sffamily\bfseries Memory Type} & {\sffamily\bfseries Typical Memory Size (Bytes)} & {\sffamily\bfseries Typical Access Times}\\\hline
CPU Registers & 512 \ldots~2048 & $\sim$250ps\\ \hline
Cache Line  & 64 \ldots~256 & NA\\ \hline
L1 Cache & 16Kb \ldots~64Kb & $\sim$1ns\\ \hline
L2 Cache & 1Mb \ldots~2Mb & $\sim$10ns\\ \hline
L3 Cache & 8Mb \ldots~32Mb & $\sim$80ns--120ns\\ \hline
L4 Cache & 32Mb \ldots~128Mb & $\sim$100ns--200ns\\ \hline
Set Associativity  &  2 \ldots~64 & NA\\ \hline
TL & 4 words \ldots~65536 & 10ns \ldots~50ns\\ \hline
Physical Memory Page & 512 \ldots~8192 & 100ns \ldots~500ns\\ \hline
\rule{0pt}{3.5mm}Virtual Memory & $2^{32}$ bytes \ldots~$2^{64}$ bytes & $\sim$10$\mu$s--50$\mu$s\\ \hline
Solid-State Disc (SSD) & 256Gb \ldots~16Tb & $\sim$25$\mu$s--100$\mu$s\\ \hline
Mechanical Disc & Huge & $\sim$5ms--10ms\\ \hline
Clock Speed & NA & $\sim$4GHz \\ \thickhline
\end{tabular}
}
\end{threeparttable}
    \end{center}
\end{table}
\clearpage



\newpage
%\cftaddtitleline{toc}{section}{\protect\hspace*{-5em}\framebox{\sffamily\itshape\cppxx}\vspace*{-3ex}}{}
% \section[{\tt auto}]{{\SecCode auto}\checkoddpage\Cbox}\label{auto-feature}
%\input{auto-placeholder}
%\cftaddtitleline{toc}{section}{\protect\hspace*{-5em}\framebox{\sffamily\itshape\cppxx}\vspace*{-3ex}}{}
\section[{\tt auto} Variables]{Variables of Automatically Deduced Type}\label{auto}\label{auto-feature}
\input{auto-placeholder}

\newpage
 \section[Braced Init]{Brace-Initialization Syntax: {\SecCode \{\}}}\label{bracedinit}
\input{bracedinit-placeholder}

\newpage
\section[{\tt constexpr} Functions]{Compile-Time Evaluatable Functions}\label{constexprfunc}


\emcppsFeature{
    short={\lstinline!constexpr! Functions},
    tocshort={{\TOCCode constexpr} Functions},
    long={Compile-Time Evaluatable Functions},
    rhshort={{\RHCode constexpr} Functions},
}{constexprfunc}
\setcounter{table}{0}
\setcounter{footnote}{0}
\setcounter{lstlisting}{0}
%\section[{\tt constexpr} Functions]{Compile-Time Evaluatable Functions}\label{constexprfunc}

placeholder text.........
 
 

\newpage
\section[{\tt constexpr} Variables]{Compile-Time Accessible Variables}\label{constexprvar}


\emcppsFeature{
    short={\lstinline!constexpr! Variables},
    tocshort={{\TOCCode constexpr} Variables},
    long={Compile-Time Accessible Variables},
    rhshort={{\RHCode constexpr} Variables},
}{constexprvar}
\setcounter{table}{0}
\setcounter{footnote}{0}
\setcounter{lstlisting}{0}
%\section[{\tt constexpr} Variables]{Compile-Time Accessible Variables}\label{constexprvar}

placeholder text...........




\newpage
\section[Default Member Init]{Default {\SecCode class}/{\SecCode union} Member Initializers}\label{Default-Member-Initializers}
\input{defmemberinit-placeholder}

\newpage
\section[{\tt enum}~{\tt class}]{Strongly Typed Scoped Enumerations}\label{enumclass}
% 01 Feb, 2021 JMB - packet 4 compilation fixes

\texttt{enum}~\texttt{class} is an alternative to the classic
\texttt{enum} construct that simultaneously provides both stronger
typing and an enclosing scope for its enumerated values.

\subsection[Description]{Description}\label{description-enumclass}

Classic, C-style enumerations are useful and continue to fulfill
important engineering needs:

\begin{emcppslisting}
enum EnumName { e_Enumerator0 /*= value0 */, e_EnumeratorN /* = valueN */ };
    // classic, C-style (ù{\codeincomments{enum}}ù): enumerators are neither type-safe nor scoped
\end{emcppslisting}

\noindent For more examples where the classic \texttt{enum} shines, see \textit{\titleref{potential-pitfalls-enumclass}: \titleref{strong-typing-of-an-enum-class-can-be-counterproductive}} on page~\pageref{strong-typing-of-an-enum-class-can-be-counterproductive} and \textit{\titleref{annoyances-enumclass}: \titleref{scoped-enumerations-do-not-necessarily-add-value}} on page~\pageref{scoped-enumerations-do-not-necessarily-add-value}. Still,
innumerable practical situations occur in which enumerators that are
both scoped and more type-safe would be preferred; see \textit{\titleref{introducing-the-c++11-enum-class}} on page~\pageref{introducing-the-c++11-enum-class}.

\subsubsection[Drawbacks and workarounds relating to unscoped C++03 enumerations]{Drawbacks and workarounds relating to unscoped C++03 enumerations}\label{drawbacks-and-workarounds-relating-to-unscoped-c++03-enumerations}

Since the enumerators of a classic \texttt{enum} leak out into the
enclosing scope, if two unrelated enumerations that happen to use the
same enumerator name appear in the same scope, an ambiguity could ensue:

\begin{emcppslisting}
enum Color { e_RED, e_ORANGE, e_YELLOW };    // OK
enum Fruit { e_APPLE, e_ORANGE, e_BANANA };  // Error, (ù{\codeincomments{e\_ORANGE}}ù) is redefined.
\end{emcppslisting}

\noindent Note that
we use a lowercase, single-letter prefix, such as \texttt{e\_}, to
ensure that the uppercase enumerator name is less likely to collide
  with a legacy macro, which is especially useful in header files. The problems associated with the use of unscoped enumerations is
exacerbated when those enumerations are placed in their own respective
header files in the global or some other large namespace scope, such as
\texttt{std}, for general reuse. In such cases, latent defects will
typically not manifest unless and until the two enumerations are
included in the same translation unit.

If the only issue were the leakage of the enumerators into the enclosing
scope, then the long-established workaround of enclosing the enumeration
within a \texttt{struct} would suffice:

\begin{emcppslisting}[emcppsbatch=e1]
struct Color { enum Enum { e_RED, e_ORANGE, e_YELLOW };  };   // OK
struct Fruit { enum Enum { e_APPLE, e_ORANGE, e_BANANA }; };  // OK (scoped)
\end{emcppslisting}

\noindent Employing the C++03 workaround in the above code snippet implies that,
when passing such an explicitly scoped, classical \texttt{enum} into a
function, the distinguishing name of the \texttt{enum} is subsumed by
its enclosing \texttt{struct} and the \texttt{enum} name itself, such as
\texttt{Enum}, becomes \textbf{boilerplate code}:

\begin{emcppslisting}[emcppsbatch=e1]
int enumeratorValue1 = Color::e_ORANGE;  // OK
int enumeratorValue2 = Fruit::e_ORANGE;  // OK

void colorFunc(Color::Enum color);  // enumerated (scoped) (ù{\codeincomments{Color}}ù) parameter
void fruitFunc(Fruit::Enum fruit);  // enumerated (scoped) (ù{\codeincomments{Fruit}}ù) parameter
\end{emcppslisting}

\noindent Hence, adding \emph{just} scope to a classic, C++03 \texttt{enum} is
easily doable and might be exactly what is indicated; see
\textit{\titleref{potential-pitfalls-enumclass}: \titleref{strong-typing-of-an-enum-class-can-be-counterproductive}} on page~\pageref{strong-typing-of-an-enum-class-can-be-counterproductive}.

\subsubsection[Drawbacks relating to weakly typed, C++03 enumerators]{Drawbacks relating to weakly typed, C++03 enumerators}\label{drawbacks-relating-to-weakly-typed,-c++03-enumerators}

Historically, C++03 enumerations have been employed to represent at
least two distinct concepts:

\begin{enumerate}
\item{A collection of related, but not necessarily unique, named integral values}
\item{A pure, perhaps ordered, set of named entities in which cardinal value has no relevance}
\end{enumerate}

\noindent It will turn out that the modern \texttt{enum}~\texttt{class} feature,
which we will discuss in \textit{\titleref{description-enumclass}: \titleref{introducing-the-c++11-enum-class}}, is more closely aligned with this second
concept.

A classic enumeration, by default, has an implementation-defined
\textbf{underlying type} (see ``\titleref{explicit-enumeration-underlying-type}" on page~\pageref{explicit-enumeration-underlying-type}), which it uses to represent
variables of that enumerated type as well as the values of its
enumerators. While implicit conversion \emph{to} an enumerated type is
never permitted, when implicitly converting \emph{from} a classical
\texttt{enum} type to some arithmetic type, the \texttt{enum} promotes
to integral types in a way similar to how its underlying type would
promote using the rules of \textbf{integral promotion} and
\textbf{standard conversion}:

\begin{emcppslisting}
void f()
{
    enum A { e_A0, e_A1, e_A2 };  // classic, C-style C++03 (ù{\codeincomments{enum}}ù)
    enum B { e_B0, e_B1, e_B2 };  //    "       "       "     "

    A a;  // Declare object (ù{\codeincomments{a}}ù) to be of type (ù{\codeincomments{A}}ù).
    B b;  //   "       "    (ù{\codeincomments{b}}ù)  " "  "       (ù{\codeincomments{B}}ù).

    a = e_B2;  // Error, cannot convert (ù{\codeincomments{e\_B2}}ù) to (ù{\codeincomments{enum}}ù) type (ù{\codeincomments{A}}ù)
    b = e_B2;  // OK, assign the value (ù{\codeincomments{e\_B2}}ù) (numerically 2) to (ù{\codeincomments{b}}ù).
    a = b;     // Error, cannot convert enumerator (ù{\codeincomments{b}}ù) to (ù{\codeincomments{enum}}ù) type (ù{\codeincomments{A}}ù)
    b = b;     // OK, self-assignment
    a = 1;     // Error, invalid conversion from (ù{\codeincomments{int 1}}ù) to (ù{\codeincomments{enum}}ù) type (ù{\codeincomments{A}}ù)
    a = 0;     // Error, invalid conversion from (ù{\codeincomments{int 0}}ù) to (ù{\codeincomments{enum}}ù) type (ù{\codeincomments{A}}ù)

    bool     v = a;     // OK
    char     w = e_A0;  // OK
    unsigned y = e_B1;  // OK
    float    x = b;     // OK
    double   z = e_A2;  // OK
    char*    p = e_B0;  // Error, unable to convert (ù{\codeincomments{e\_B0}}ù) to (ù{\codeincomments{char*}}ù)
    char*    q = +e_B0; // Error, invalid conversion of (ù{\codeincomments{int}}ù) to (ù{\codeincomments{char*}}ù)
}
\end{emcppslisting}

\noindent Notice that, in this example, the final two diagnostics for the
attempted initializations of \texttt{p} and \texttt{q}, respectively,
differ slightly. In the first, we are trying to initialize a pointer,
\texttt{p}, with an enumerated type, \texttt{B}. In the second, we have
creatively used the built-in unary-plus operator to explicitly promote
the enumerator to an integral type before attempting to assign it to a
pointer, \texttt{q}. Even though the numerical value of the enumerator
is \texttt{0} and such is known at compile time, implicit
conversion to a pointer type from anything but the literal integer
constant \texttt{0} is not
permitted. Excluding esoteric user-defined
types, only a literal \texttt{0} or, as of C++11, a value of type
\texttt{std::nullptr\_t} is implicitly convertible to an arbitrary
  pointer type; see ``\titleref{null-pointer-literal-(nullptr)}" on page~\pageref{null-pointer-literal-(nullptr)}.

C++ fully supports comparing values of \emph{classic} \texttt{enum}
types with values of arbitrary \textbf{arithmetic type} as well as those
of the same enumerated type; the operands of a comparator will be
promoted to a sufficiently large integer type and the comparison will be
done with those values. Comparing values having distinct enumerated
types, however, is deprecated and will typically elicit a
warning.{\cprotect\footnote{As of C++20, attempting to compare two values of distinct classically enumerated
types is a compile-time error. Note that explicitly converting at least one of them to an
integral type --- for example, using built-in unary plus --- both
makes our intentions clear and avoids warnings.

\begin{emcppshiddenlisting}[emcppsbatch=e2]
    enum A { e_A0, e_A1, e_A2 };
    enum B { e_B0, e_B1, e_B2 };
\end{emcppshiddenlisting}
\begin{emcppslisting}[emcppsbatch=e2,basicstyle={\ttfamily\footnotesize}]
void test()
{
  if (e_A0 < 0)       { /* ... */ }  // OK, comparison with integral type
  if (1.0 != e_B1)    { /* ... */ }  // OK, comparison with arithmetic type
  if (A() <= e_A2)    { /* ... */ }  // OK, comparison with same enumerated type
  if (e_A0 == e_B0)   { /* ... */ }  // warning, deprecated (error as of C++20)
  if ( e_A0 == +e_B0) { /* ... */ }  // OK, unary (ù{\codeincomments{+}}ù) converts to integral type
  if (+e_A0 ==  e_B0) { /* ... */ }  // OK,   "        "     "     "      "
  if (+e_A0 == +e_B0) { /* ... */ }  // OK,   "        "     "     "      "
}
\end{emcppslisting}
      }}

\subsubsection[Introducing the C++11 {\tt enum class}]{Introducing the C++11 {\SubsubsecCode enum class}}\label{introducing-the-c++11-enum-class}

With the advent of modern C++, we now have a new, alternative
enumeration construct, \texttt{enum}~\texttt{class}, that simultaneously
addresses strong type safety and lexical scoping, two distinct and often
desirable properties:

\begin{emcppslisting}
enum class Name { e_Enumerator0 /* = value0 */, e_EnumeratorN /* = valueN */ };
    // (ù{\codeincomments{enum class}}ù) enumerators are both type-safe and scoped
\end{emcppslisting}

\noindent Another major distinction is that the default
\textbf{underlying type} for a C-style \texttt{enum} is
\textbf{implementation defined}, whereas, for an
\texttt{enum}~\texttt{class}, it is always an \texttt{int}. See
\textit{\titleref{description-enumclass}: \titleref{enum-class-and-underlying-type}} on page~\pageref{enum-class-and-underlying-type} and
\textit{\titleref{potential-pitfalls-enumclass}: \titleref{external-use-of-opaque-enumerators-enumclass}} on page~\pageref{external-use-of-opaque-enumerators-enumclass}.

The enumerators within an \texttt{enum}~\texttt{class} are all scoped by
its name, while classic enumerations leak the enumerators into the
enclosing scope:

\begin{emcppslisting}
enum Vehicle  { e_CAR,   e_TRAIN, e_PLANE };
enum Geometry { e_POINT, e_LINE,  e_PLANE };  // Error, (ù{\codeincomments{e\_PLANE}}ù) is redefined.
\end{emcppslisting}

\noindent Unlike unscoped enumerations, \texttt{enum}~\texttt{class} does not leak
its enumerators into the enclosing scope and can therefore help avoid
collisions with other enumerations having like-named enumerators defined
in the same scope:

\begin{emcppslisting}[emcppsbatch=e3]
enum       VehicleUnscoped  { e_CAR, e_TRAIN, e_PLANE };
struct     VehicleScopedExplicitly { enum Enum { e_CAR, e_TRAIN, e_PLANE }; };
enum class VehicleScopedImplicitly { e_CAR, e_BOAT,  e_PLANE };
\end{emcppslisting}

\noindent Just like an unscoped \texttt{enum} type, an object of type
\texttt{enum}~\texttt{class} is passed as a parameter to a function
using the enumerator name itself:

\begin{emcppslisting}[emcppsbatch=e3]
void f1(VehicleUnscoped value);          // classic enumeration passed by value
void f2(VehicleScopedImplicitly value);  // modern enumeration passed by value
\end{emcppslisting}

If we use the approach for
adding scope to enumerators that is described in \textit{\titleref{description-enumclass}: \titleref{drawbacks-relating-to-weakly-typed,-c++03-enumerators}} on page~\pageref{drawbacks-relating-to-weakly-typed,-c++03-enumerators}, the name of the enclosing \texttt{struct} together with a
consistent name for the enumeration, such as \texttt{Enum}, has to be used to indicate an
enumerated type:

\begin{emcppslisting}[emcppsbatch=e3,basicstyle={\ttfamily\footnotesize}]
void f3(VehicleScopedExplicitly::Enum value);
    // classically scoped (ù{\codeincomments{enum}}ù) passed by value
\end{emcppslisting}

\noindent Qualifying the enumerators of a scoped enumeration is the same,
irrespective of whether the scoping is explicit or implicit:

\begin{emcppslisting}[emcppsbatch=e3]
void g()
{
    f1(VehicleUnscoped::e_PLANE);
        // call (ù{\codeincomments{f1}}ù) with an explicitly scoped enumerator

    f2(VehicleScopedImplicitly::e_PLANE);
        // call (ù{\codeincomments{f2}}ù) with an implicitly scoped enumerator
}
\end{emcppslisting}

\noindent Apart from implicit scoping, the modern, C++11
\texttt{enum}~\texttt{class} deliberately does \emph{not} support
implicit conversion, in any context, to its \textbf{underlying type}:

\begin{emcppslisting}[emcppsbatch=e3]
void h()
{
    int i1 = VehicleScopedExplicitly::e_PLANE;
       // OK, scoped C++03 (ù{\codeincomments{enum}}ù) (implicit conversion)
 
    int i2 = VehicleScopedImplicitly::e_PLANE;
        // Error, no implicit conversion to underlying type

    if (VehicleScopedExplicitly::e_PLANE > 3) {} // OK
    if (VehicleScopedImplicitly::e_PLANE > 3) {} // Error, implicit conversion
}
\end{emcppslisting}

\noindent Enumerators of an \texttt{enum}~\texttt{class} do, however, admit
equality and ordinal comparisons within their own type:

\begin{emcppslisting}[emcppsbatch=e4]
enum class E { e_A, e_B, e_C };  // By default, enumerators increase from 0.

static_assert(E::e_A < E::e_C, "");  // OK, comparison between same-type values
static_assert(0 == E::e_A, "");      // Error, no implicit conversion from (ù{\codeincomments{E}}ù)
static_assert(0 == static_cast<int>(E::e_A), "");  // OK, explicit conversion

void f(E v)
{
    if (v > E::e_A) { /* ... */ }  // OK, comparing values of same type, (ù{\codeincomments{E}}ù)
}
\end{emcppslisting}

\noindent Note that incrementing an enumerator variable from one strongly typed
enumerator's value to the next requires an explicit cast; see \textit{\titleref{potential-pitfalls-enumclass}: \titleref{strong-typing-of-an-enum-class-can-be-counterproductive}} on page~\pageref{strong-typing-of-an-enum-class-can-be-counterproductive}.

\subsubsection[{\tt enum} {\tt class} and underlying type]{{\SubsubsecCode enum} {\SubsubsecCode class} and underlying type}\label{enum-class-and-underlying-type}

Since C++11, both scoped and unscoped enumerations permit explicit
specification of their integral \textbf{underlying type}:

\begin{emcppslisting}
enum Ec : char { e_X, e_Y, e_Z };
    // underlying type is (ù{\codeincomments{char}}ù)

static_assert(1 == sizeof(Ec),     "");
static_assert(1 == sizeof Ec::e_X, "");

enum class Es : short { e_X, e_Y, e_Z };
    // underlying type is (ù{\codeincomments{short int}}ù)

static_assert(sizeof(short) == sizeof(Es),     "");
static_assert(sizeof(short) == sizeof Es::e_X, "");
\end{emcppslisting}

\noindent Unlike a classic \texttt{enum}, which has an
\textbf{implementation-defined} default \textbf{underlying type}, the
default \textbf{underlying type} for an \texttt{enum}~\texttt{class} is
always \texttt{int}:

\begin{emcppslisting}
enum class Ei { e_X, e_Y, e_Z };
    // When not specified the underlying type of an (ù{\codeincomments{enum class}}ù) is (ù{\codeincomments{int}}ù).

static_assert(sizeof(int) == sizeof(Ei),     "");
static_assert(sizeof(int) == sizeof Ei::e_X, "");
\end{emcppslisting}

\noindent Note that, because the default \textbf{underlying type} of an
\texttt{enum}~\texttt{class} is specified by the Standard, eliding the
enumerators of an \texttt{enum}~\texttt{class} in
a local redeclaration is \emph{always} possible; see
\textit{\titleref{potential-pitfalls-enumclass}: \titleref{external-use-of-opaque-enumerators-enumclass}} on page~\pageref{external-use-of-opaque-enumerators-enumclass} and ``\titleref{enumopaque}" on page~\pageref{enumopaque}.

\subsection[Use Cases]{Use Cases}\label{use-cases-enumclass}

\subsubsection[Avoiding unintended implicit conversions to arithmetic types]{Avoiding unintended implicit conversions to arithmetic types}\label{avoiding-unintended-implicit-conversions-to-arithmetic-types}

Suppose that we want to represent the result of selecting one of a fixed
number of alternatives from a drop-down menu as a simple unordered set
of uniquely valued named integers. For example, this might be the case
when configuring a product, such as a vehicle, for purchase:

\begin{emcppslisting}
struct Trans
{
    enum Enum { e_MANUAL, e_AUTOMATIC };  // classic, C++03 scoped (ù{\codeincomments{enum}}ù)
};
\end{emcppslisting}

\noindent Although automatic promotion of a classic enumerator to \texttt{int}
works well when typical use of the enumerator involves knowing its
cardinal value, such promotions are less than ideal when cardinal values
have no role in intended usage:

\begin{emcppshiddenlisting}[emcppsbatch=e5]
void attachDoor(int i);
\end{emcppshiddenlisting}
\begin{emcppslisting}[emcppsbatch=e5]
class Car { /* ... */ };

struct Trans
{                                         // explicitly scoped
    enum Enum { e_MANUAL, e_AUTOMATIC };  // classic (ù{\codeincomments{enum}}ù)
};                                        // (BAD IDEA)

int buildCar(Car* result, int numDoors, Trans::Enum trans)
{
    int status = Trans::e_MANUAL;    // Bug, accidental misuse!

    for (int i = 0; i < trans; ++i)  // Bug, accidental misuse!
    {
         attachDoor(i);
    }

    return status;
}
\end{emcppslisting}

\noindent As shown in the example above, it is never correct for a value of type
\texttt{Trans::Enum} to be assigned to, compared with, or otherwise
modified like an integer; hence, \emph{any} such use would necessarily
be considered a mistake and, ideally, flagged by the compiler as an
error. The stronger typing provided by \texttt{enum}~\texttt{class}
achieves this goal:

\begin{emcppslisting}[emcppsignore={Fails to build with all Declarations marked Error}]
class Car { /* ... */ };

enum class Trans { e_MANUAL, e_AUTOMATIC };  // modern (ù{\codeincomments{enum class}}ù) (GOOD IDEA)

int buildCar(Car* result, int numDoors, Trans trans)
{
    int status = Trans::e_MANUAL;    // Error, incompatible types

    for (int i = 0; i < trans; ++i)  // Error, incompatible types
    {
         attachDoor(i);
    }

    return status;
}
\end{emcppslisting}

\noindent By deliberately choosing the \texttt{enum}~\texttt{class} over the
\emph{classic} \texttt{enum} above, we automate the detection of many common
kinds of accidental misuse. Secondarily, we slightly simplify the interface of the function signature by removing the extra \texttt{::Enum} boilerplate qualifications required of an explicitly
scoped, less-type-safe, classic \texttt{enum}, but see \textit{\titleref{potential-pitfalls-enumclass}: \titleref{strong-typing-of-an-enum-class-can-be-counterproductive}} on page~\pageref{strong-typing-of-an-enum-class-can-be-counterproductive}.

In an unlikely event that the numeric value of a strongly typed
enumerator is needed (e.g., for serialization), it can be extracted
explicitly via a \texttt{static\_cast}:

\begin{emcppshiddenlisting}[emcppsbatch=e6]
class Car { /* ... */ };

enum class Trans { e_MANUAL, e_AUTOMATIC };  // modern (ù{\codeincomments{enum class}}ù) (GOOD IDEA)
\end{emcppshiddenlisting}
\begin{emcppslisting}[emcppsbatch=e6]
const int manualIntegralValue    = static_cast<int>(Trans::e_MANUAL);
const int automaticIntegralValue = static_cast<int>(Trans::e_AUTOMATIC);
static_assert(0 == manualIntegralValue,    "");
static_assert(1 == automaticIntegralValue, "");
\end{emcppslisting}

\subsubsection[Avoiding namespace pollution]{Avoiding namespace pollution}\label{avoiding-namespace-pollution}

Classic, C-style enumerations do not provide scope for their
enumerators, leading to unintended latent name collisions:

\begin{emcppslisting}
// vehicle.h:
// ...
enum Vehicle  { e_CAR, e_TRAIN, e_PLANE };  // classic, C-style (ù{\codeincomments{enum}}ù)
// ...

// geometry.h:
// ...
enum Geometry { e_POINT, e_LINE,  e_PLANE };  // classic, C-style (ù{\codeincomments{enum}}ù)
// ...

// client.cpp:
#include <vehicle.h>  // OK
#include <geometry.h> // Error, (ù{\codeincomments{e\_PLANE}}ù) redefined
// ...
\end{emcppslisting}

\noindent The common workaround is to wrap the \texttt{enum} in a \texttt{struct}
or \texttt{namespace}:

\begin{emcppslisting}
// vehicle.h:
// ...
struct Vehicle {                            // explicitly scoped
    enum Enum { e_CAR, e_TRAIN, e_PLANE };  // classic, C-style (ù{\codeincomments{enum}}ù)
};
// ...

// geometry.h:
// ...
struct Geometry {                            // explicitly scoped
    enum Enum { e_POINT, e_LINE, e_PLANE };  // classic, C-style (ù{\codeincomments{enum}}ù)
};
// ...

// client.cpp:
#include <vehicle.h>    // OK
#include <geometry.h>   // OK, enumerators are scoped explicitly.
// ...
\end{emcppslisting}

\noindent If implicit conversions of enumerators to integral types are not
required, we can achieve the same scoping effect with much more type
safety and slightly less boilerplate --- i.e., without the
\texttt{::Enum} when declaring a variable --- by employing
\texttt{enum}~\texttt{class} instead:

\begin{emcppslisting}
// vehicle.h:
// ...
enum class Vehicle { e_CAR, e_TRAIN, e_PLANE };
// ...

// geometry.h:
// ...
enum class Geometry { e_POINT, e_LINE, e_PLANE };
// ...

// client.cpp:
#include <vehicle.h>  // OK
#include <geometry.h>   // OK, enumerators are scoped implicitly.
// ...
\end{emcppslisting}

\subsubsection[Improving overloading disambiguation]{Improving overloading disambiguation}\label{improving-overloading-disambiguation}

Overloaded functions are notorious for providing opportunities for
misuse. Maintenance difficulties are exacerbated when arguments for
these overloads are convertible to more than a single parameter in the
function. As an illustration of the compounding of such maintenance
difficulties, suppose that we have a widely used, named type,
\texttt{Color}, and the numeric values of its enumerators are small,
unique, and irrelevant. Imagine we have chosen to represent
\texttt{Color} as a \emph{classic} \texttt{enum}:

\begin{emcppslisting}[emcppsbatch=e7]
struct Color
{                                           // explicitly scoped
    enum Enum { e_RED, e_BLUE /*, ...*/ };  // classic, C-style (ù{\codeincomments{enum}}ù)
};                                          // (BAD IDEA)
\end{emcppslisting}

\noindent Suppose further that we have provided two overloaded functions, each
having two parameters, with one signature's parameters including the
enumeration \texttt{Color}:

\begin{emcppslisting}[emcppsbatch=e7]
void clearScreen(int pattern, int orientation);          // (0)
void clearScreen(Color::Enum background, double alpha);  // (1)
\end{emcppslisting}

\noindent Depending on the types of the arguments supplied, one or the other
functions will be selected or else the call will be ambiguous and the
program will fail to compile{\cprotect\footnote{GCC version 7.4.0
incorrectly diagnoses both ambiguity errors as warnings, although it
states in the warning that it is an error:

\begin{lstlisting}[language=bash,basicstyle={\ttfamily\footnotesize}]
warning: ISO C++ says that these are ambiguous, even though the worst conversion for the
         first is better than the worst conversion for the second:

note: candidate 1: void clearScreen(int, int)
void clearScreen(int pattern, int orientation);
     ^~~~~~~~~~~
note: candidate 2: void clearScreen(Color::Enum, double)
void clearScreen(Color::Enum background, double alpha;
     ^~~~~~~~~~~
\end{lstlisting}
      }}:

\begin{emcppslisting}[emcppsbatch=e7]
void f0()
{
    clearScreen(1           , 1           );  // calls (0) above
    clearScreen(1           , 1.0         );  // calls (0) above
    clearScreen(1           , Color::e_RED);  // calls (0) above

    clearScreen(1.0         , 1           );  // calls (0) above
    clearScreen(1.0         , 1.0         );  // calls (0) above
    clearScreen(1.0         , Color::e_RED);  // calls (0) above

    clearScreen(Color::e_RED, 1           );  // Error, ambiguous call
    clearScreen(Color::e_RED, 1.0         );  // calls (1) above
    clearScreen(Color::e_RED, Color::e_RED);  // Error, ambiguous call
}
\end{emcppslisting}

\noindent Now suppose that we had instead defined our \texttt{Color} enumeration
as a modern \texttt{enum}~\texttt{class}:

\begin{emcppslisting}[emcppsbatch=e8]
enum class Color { e_RED, e_BLUE /*, ...*/ };

void clearScreen(int pattern, int orientation);    // (2)
void clearScreen(Color background, double alpha);  // (3)
\end{emcppslisting}

\noindent The function that will be called from a given set of arguments becomes
clear:

\begin{emcppslisting}[emcppsbatch=e8]
void f1()
{
    clearScreen(1           , 1           );  // calls (2) above
    clearScreen(1           , 1.0         );  // calls (2) above
    clearScreen(1           , Color::e_RED);  // Error, no matching function

    clearScreen(1.0         , 1           );  // calls (2) above
    clearScreen(1.0         , 1.0         );  // calls (2) above
    clearScreen(1.0         , Color::e_RED);  // Error, no matching function

    clearScreen(Color::e_RED, 1           );  // calls (3) above
    clearScreen(Color::e_RED, 1.0         );  // calls (3) above
    clearScreen(Color::e_RED, Color::e_RED);  // Error, no matching function
}
\end{emcppslisting}

\noindent Returning to our original, classic-\texttt{enum} design, suppose that we
find we need to add a third parameter, \texttt{bool}~\texttt{z}, to the
second overload:

\begin{emcppshiddenlisting}[emcppsbatch=e9]
struct Color
{
    enum Enum { e_RED, e_BLUE /*, ...*/ };
};
\end{emcppshiddenlisting}
\begin{emcppslisting}[emcppsbatch=e9]
void clearScreen(int pattern, int orientation);                  // (0)
void clearScreen(Color::Enum background, double alpha, bool z);  // (4) classic
\end{emcppslisting}

\noindent If our plan is that any existing client calls involving
\texttt{Color::Enum} will now be flagged as errors, we are going to be
very disappointed:

\begin{emcppslisting}[emcppsbatch=e9]
void f2()
{
    clearScreen(Color::e_RED, 1.0);  // calls (0) above
}
\end{emcppslisting}

\noindent In fact, every combination of arguments above --- all nine of them ---
will call function (0) above with no warnings at all:

\begin{emcppslisting}[emcppsbatch=e9]
void f3()
{
    clearScreen(1           , 1           );  // calls (0) above
    clearScreen(1           , 1.0         );  // calls (0) above
    clearScreen(1           , Color::e_RED);  // calls (0) above

    clearScreen(1.0         , 1           );  // calls (0) above
    clearScreen(1.0         , 1.0         );  // calls (0) above
    clearScreen(1.0         , Color::e_RED);  // calls (0) above

    clearScreen(Color::e_RED, 1           );  // calls (0) above
    clearScreen(Color::e_RED, 1.0         );  // calls (0) above
    clearScreen(Color::e_RED, Color::e_RED);  // calls (0) above
}
\end{emcppslisting}

\noindent Finally, let's suppose again that we have used
\texttt{enum}~\texttt{class} to implement our \texttt{Color}\linebreak[4]
\mbox{enumeration}:

\begin{emcppshiddenlisting}[emcppsbatch=e10]
enum class Color { e_RED, e_BLUE /*, ...*/ };
\end{emcppshiddenlisting}
\begin{emcppslisting}[emcppsbatch=e10]
void clearScreen(int pattern, int orientation);            // (2)
void clearScreen(Color background, double alpha, bool z);  // (5) modern

void f4()
{
    clearScreen(Color::e_RED, 1.0);  // Error, no matching function
}
\end{emcppslisting}

\noindent And in fact, the \emph{only} calls that succeed unmodified are precisely
those that do not involve the enumeration \texttt{Color}, as desired:

\begin{emcppslisting}[emcppsbatch=e10]
void f5()
{
    clearScreen(1           , 1           );  // calls (2) above
    clearScreen(1           , 1.0         );  // calls (2) above
    clearScreen(1           , Color::e_RED);  // Error, no matching function

    clearScreen(1.0         , 1           );  // calls (2) above
    clearScreen(1.0         , 1.0         );  // calls (2) above
    clearScreen(1.0         , Color::e_RED);  // Error, no matching function

    clearScreen(Color::e_RED, 1           );  // Error, no matching function
    clearScreen(Color::e_RED, 1.0         );  // Error, no matching function
    clearScreen(Color::e_RED, Color::e_RED);  // Error, no matching function
}
\end{emcppslisting}

\noindent Bottom line: Having a \emph{pure} enumeration --- such as
\texttt{Color}, used widely in function signatures --- be strongly typed
can only help to expose accidental misuse but, again, see
\textit{\titleref{potential-pitfalls-enumclass}: \titleref{strong-typing-of-an-enum-class-can-be-counterproductive}} on page~\pageref{strong-typing-of-an-enum-class-can-be-counterproductive}.

Note that strongly typed enumerations help to avoid accidental misuse by
requiring an explicit \emph{cast} should conversion to an arithmetic
type be desired:

\begin{emcppslisting}[emcppsbatch=e10]
void f6()
{
    clearScreen(Color::e_RED, 1.0);                   // Error, no match
    clearScreen(static_cast<int>(Color::e_RED), 1.0); // OK, calls (2) above
    clearScreen(Color::e_RED, 1.0, false);            // OK, calls (5) above
}
\end{emcppslisting}


\subsubsection[Encapsulating implementation details within the enumerators themselves]{Encapsulating implementation details within the enumerators themselves}\label{encapsulating-implementation-details-within-the-enumerators-themselves}

In rare cases, providing a pure, ordered enumeration having unique (but
not necessarily contiguous) numerical values that exploit lower-order
bits to categorize and make readily available important
individual properties might offer an advantage, such as in performance. Note that in order to preserve the ordinality of the enumerators
overall, the higher-level bits must encode their relative order. The
lower-level bits are then available for arbitrary use in the
  implementation..

For example, suppose that we have a \texttt{MonthOfYear} enumeration
that encodes the months that have 31 days in their least-significant bit
and an accompanying \texttt{inline} function to quickly determine
whether a given enumerator represents such a month:

\begin{emcppslisting}
#include <type_traits>  // (ù{\codeincomments{std::underlying\_type}}ù)

enum class MonthOfYear : unsigned char  // optimized to flag long months
{
    e_JAN = ( 1 << 4) + 0x1,
    e_FEB = ( 2 << 4) + 0x0,
    e_MAR = ( 3 << 4) + 0x1,
    e_APR = ( 4 << 4) + 0x0,
    e_MAY = ( 5 << 4) + 0x1,
    e_JUN = ( 6 << 4) + 0x0,
    e_JUL = ( 7 << 4) + 0x1,
    e_AUG = ( 8 << 4) + 0x1,
    e_SEP = ( 9 << 4) + 0x0,
    e_OCT = (10 << 4) + 0x1,
    e_NOV = (11 << 4) + 0x0,
    e_DEC = (12 << 4) + 0x1
};

bool hasThirtyOneDays(MonthOfYear month)
{
    return static_cast<std::underlying_type<MonthOfYear>::type>(month) & 0x1;
}
\end{emcppslisting}

\noindent In the example above, we are using a new cross-cutting
feature of all enumerated types that allows the client defining the
type to specify its underlying type precisely. In this case, we have
chosen an \texttt{unsigned}~\texttt{char} to maximize the number of
flag bits while keeping the overall size to a single byte. Three bits
remain available. Had we needed more flag bits, we could have just as
easily used a larger underlying type, such as
  \texttt{unsigned}~\texttt{short}; see ``\titleref{explicit-enumeration-underlying-type}" on page~\pageref{explicit-enumeration-underlying-type}.

\noindent In case \texttt{enum}s are used for encoding purposes, the public clients are not intended to make use of the
cardinal values; hence clients are well advised to treat them as
implementation details, potentially subject to change without notice.
Representing this enumeration using the modern
\texttt{enum}~\texttt{class}, instead of an explicitly scoped classic
\texttt{enum}, deters clients from making any use (apart from same-type
comparisons) of the cardinal values assigned to the enumerators. Notice
that implementors of the \texttt{hasThirtyOneDays} function will require
a verbose but efficient \texttt{static\_cast} to resolve the
cardinal value of the enumerator and thus make the requested
determination as efficiently as possible.

\subsection[Potential Pitfalls]{Potential Pitfalls}\label{potential-pitfalls-enumclass}

\subsubsection[Strong typing of an {\tt enum} {\tt class} can be counterproductive]{Strong typing of an {\SubsubsecCode enum} {\SubsubsecCode class} can be counterproductive}\label{strong-typing-of-an-enum-class-can-be-counterproductive}

The additive value in using a modern \texttt{enum}~\texttt{class} is
governed \emph{solely} by whether its stronger typing, \emph{not} its
implicit scoping, of its enumerators would be beneficial in its
anticipated typical usage. If the expectation is that the client will
never need to know the specific values of the enumerators, then use
of the modern \texttt{enum}~\texttt{class} is often just what's needed.
But if the cardinal values themselves are ever needed during typical
use, extracting them will require the client to perform an explicit
cast. Beyond mere inconvenience, encouraging clients to use casts
invites defects.

Suppose, for example, we have a function, \texttt{setPort}, from an
external library that takes an integer port number:

\begin{emcppslisting}[emcppsbatch=e11]
int setPort(int portNumber);
    // Set the current port; return (ù{\codeincomments{0}}ù) on success and a nonzero value otherwise.
\end{emcppslisting}

\noindent Suppose further that we have used the modern
\texttt{enum}~\texttt{class} feature to implement an enumeration,
\texttt{SysPort}, that identifies well-known ports on our system:

\begin{emcppslisting}[emcppsbatch=e11]
enum class SysPort { e_INPUT = 27, e_OUTPUT = 29, e_ERROR = 32, e_CTRL = 6 };
    // enumerated port values used to configure our systems
\end{emcppslisting}

\noindent Now suppose we want to call the function \texttt{f} using one of these
enumerated values:

\begin{emcppslisting}[emcppsbatch=e11]
void setCurrentPortToCtrl()
{
    setPort(SysPort::e_CTRL);  // Error, cannot convert (ù{\codeincomments{SetPort}}ù) to (ù{\codeincomments{int}}ù)
}
\end{emcppslisting}

\noindent Unlike the situation for a \emph{classic} \texttt{enum}, no implicit
conversion occurs from an \texttt{enum}~\texttt{class} to its underlying
integral type, so anyone using this enumeration will be forced to
somehow explicitly \textbf{cast} the enumerator to some arithmetic type.
There are, however, multiple choices for performing this cast:

\begin{emcppslisting}[emcppsbatch=e11]
#include <type_traits>  // (ù{\codeincomments{std::underlying\_type}}ù)

void test()
{
    setPort(int(SysPort::e_CTRL));                                       // (1)
    setPort((int)SysPort::e_CTRL);                                       // (2)
    setPort(static_cast<int>(SysPort::e_CTRL));                          // (3)
    setPort(static_cast<std::underlying_type<SysPort>::type>(            // (4)
                                                         SysPort::e_CTRL));
    setPort(static_cast<int>(                                            // (5)
             static_cast<std::underlying_type<SysPort>::type>(SysPort::e_CTRL)));
}
\end{emcppslisting}

\noindent Any of the above casts would work in this case, but consider a future
where a platform changed \texttt{setPort} to take a \texttt{long} and
the control port was changed to a value that cannot be represented as an
\texttt{int}:

\begin{emcppslisting}
int setPort(long portNumber);
enum class SysPort : unsigned { e_INPUT = 27, e_OUTPUT = 29, e_ERROR = 32,
                                e_CTRL = 0x80000000 };
    // enumerated port values used to configure our systems
\end{emcppslisting}

\noindent Only casting method (4) above will pass the correct value for
\texttt{e\_CTRL} to this new \texttt{setPort} implementation. The other
variations will all pass a negative number for the port, which would
certainly not be the intention of the user writing this code. A classic
C-style \texttt{enum} would have avoided any manually written cast
entirely and the proper value would propagate into \texttt{setPort} even
as the range of values used for ports changes:

\begin{emcppshiddenlisting}[emcppsbatch=e12]
#include <type_traits>  // (ù{\codeincomments{std::is\_same}}ù)
int setPort(long portNumber);
\end{emcppshiddenlisting}
\begin{emcppslisting}[emcppsbatch=e12]
struct SysPort  // explicit scoping for a classic, C-style (ù{\codeincomments{enum}}ù)
{
    enum Enum { e_INPUT = 27, e_OUTPUT = 29, e_ERROR = 32,
                e_CTRL = 0x80000000 };

    // Note that the underlying type of (ù{\codeincomments{Enum}}ù) is implicit and will be
    // large enough to represent all of these values.
    static_assert(
        std::is_same<std::underlying_type<Enum>::type,unsigned>::value, "");
};

void setCurrentPortToCtrl()
{
    setPort(SysPort::e_CTRL);  // OK, (ù{\codeincomments{SysPort::Enum}}ù) promotes to (ù{\codeincomments{long}}ù).
}
\end{emcppslisting}

\noindent When the intended client will depend on the cardinal values of the
enumerators during routine use, we can avoid tedious, error-prone, and
repetitive casting by instead employing a classic, C-style
\texttt{enum}, possibly nested within a \texttt{struct} to achieve
explicit scoping of its enumerators. The subsections that follow
highlight specific cases in which classic, C-style, C++03
\texttt{enum}s are appropriate.

\subsubsection[Misuse of {\tt enum} {\tt class} for collections of named constants]{Misuse of {\SubsubsecCode enum} {\SubsubsecCode class} for collections of named constants}\label{misuse-of-enum-class-for-collections-of-named-constants}

When constants are truly independent, we are often encouraged to avoid
enumerations altogether, preferring instead individual constants; see
``\titleref{Default-Member-Initializers}" on page~\pageref{Default-Member-Initializers}. On the other hand, when the constants all participate within a coherent
theme, the expressiveness achieved using a \emph{classic} \texttt{enum}
to aggregate those values is compelling. Another
advantage of an enumerator over an individual constant is that the
enumerator is guaranteed to be a \textbf{compile-time constant} (see
  ``\titleref{constexprvar}" on page~\pageref{constexprvar}) and a \textbf{prvalue} (see
  ``\titleref{Rvalue-References}" on page~\pageref{Rvalue-References}), which never needs static storage
  and cannot have its address taken.

For example, suppose we want to collect the coefficients for various
numerical suffixes representing \emph{thousands}, \emph{millions}, and
\emph{billions} using an enumeration:

\begin{emcppslisting}[emcppsbatch=e13]
enum class S0 { e_K = 1000, e_M = e_K * e_K, e_G = e_M * e_K };  // (BAD IDEA)
\end{emcppslisting}

\noindent A client trying to access one of these enumerated values would need to
cast it explicitly:

\begin{emcppslisting}[emcppsbatch=e13]
void client0()
{
    int distance = 5 * static_cast<int>(S0::e_K);  // casting is error-prone
    // ...
}
\end{emcppslisting}

\noindent By instead making the enumeration an explicitly scoped, \emph{classic}
\texttt{enum} nested within a \texttt{struct}, no casting is needed
during typical use:

\begin{emcppslisting}
struct S1  // scoped
{
    enum Enum { e_K = 1000, e_M = e_K * e_K, e_G = e_M * e_K };
        // *classic* (ù{\codeincomments{enum}}ù) (GOOD IDEA)
};

void client1()
{
    int distance = 5 * S1::e_K;  // no casting required during typical use
    // ...
}
\end{emcppslisting}

\noindent If the intent is that these constants will be specified and used in a
purely local context, we might choose to drop the enclosing scope, along
with the name of the enumeration itself; see ``\titleref{local-types-as-template-arguments}" on page~\pageref{local-types-as-template-arguments}:

\begin{emcppslisting}
void client2()
{
    enum { e_K = 1000, e_M = e_K * e_K, e_G = e_M * e_K };  // function scoped

    double salary = 95 * e_K;
    double netWorth = 0.62 * e_M;
    double companyRevenue = 47.2 * e_G;
    // ...
}
\end{emcppslisting}

\noindent We sometimes use the
lowercase prefix \texttt{k\_} instead of \texttt{e\_} to indicate
salient \textbf{compile-time constants} that are not considered part
of an enumerated set, irrespective of whether they are implemented as
enumerators:

\begin{emcppslisting}[basicstyle={\ttfamily\footnotesize}]
enum { k_NUM_PORTS = 500, k_PAGE_SIZE = 512 };     // compile-time constants
static const double k_PRICING_THRESHOLD = 0.03125; // compile-time constant
\end{emcppslisting}

\subsubsection[Misuse of {\tt enum} {\tt class} in association with bit flags]{Misuse of {\SubsubsecCode enum} {\SubsubsecCode class} in association with bit flags}\label{misuse-of-enum-class-in-association-with-bit-flags}

Using \texttt{enum}~\texttt{class} to implement enumerators that are
intended to interact closely with arithmetic types will typically
require the definition of arithmetic and bitwise operator overloads
between values of the same enumeration and between the enumeration and
arithmetic types, leading to yet more code that needs to be written,
tested, and maintained. This is often the case for bit flags. Consider, for example, an
enumeration used to control a file system:

\begin{emcppslisting}[emcppsbatch=e14]
enum class Ctrl { e_READ = 0x1, e_WRITE = 0x2, e_EXEC = 0x4 };  // (BAD IDEA)
    // low-level bit flags used to control file system

void chmodFile(int fd, int access);
    // low-level function used to change privileges on a file
\end{emcppslisting}

\noindent We could conceivably write a series of functions to combine the
individual flags in a type-safe manner:

\begin{emcppslisting}[emcppsbatch=e14]
#include <type_traits>  // (ù{\codeincomments{std::underlying\_type}}ù)

int flags() { return 0; }
int flags(Ctrl a) { return static_cast<std::underlying_type<Ctrl>::type>(a); }
int flags(Ctrl a, Ctrl b) { return flags(a) | flags(b); }
int flags(Ctrl a, Ctrl b, Ctrl c) { return flags(a, b) | flags(c); }

void setRW(int fd)
{
    chmodFile(fd, flags(Ctrl::e_READ, Ctrl::e_WRITE));  // (BAD IDEA)
}
\end{emcppslisting}

\noindent Alternatively, a \emph{classic}, C-style \texttt{enum} nested within a
\texttt{struct} achieves what's needed:

\begin{emcppslisting}
struct Ctrl // scoped
{
    enum Enum { e_READ = 0x1, e_WRITE = 0x2, e_EXEC = 0x4 };  // classic (ù{\codeincomments{enum}}ù)
        // low-level bit flags used to control file system (GOOD IDEA)
};

void chmodFile(int fd, int access);
    // low-level function used to change privileges on a file

void setRW(int fd)
{
    chmodFile(fd, Ctrl::e_READ | Ctrl::e_WRITE);  // (GOOD IDEA)
}
\end{emcppslisting}


\subsubsection[Misuse of {\tt enum} {\tt class} in association with iteration]{Misuse of {\SubsubsecCode enum} {\SubsubsecCode class} in association with iteration}\label{misuse-of-enum-class-in-association-with-iteration}

Sometimes the relative values of enumerators are considered important as
well. For example, let's again consider enumerating the months of the
year:

\begin{emcppslisting}[emcppsbatch={e15,e16}]
enum class MonthOfYear  // modern, strongly typed enumeration
{
    e_JAN, e_FEB, e_MAR,  // winter
    e_APR, e_MAY, e_JUN,  // spring
    e_JUL, e_AUG, e_SEP,  // summer
    e_OCT, e_NOV, e_DEC,  // autumn
};
\end{emcppslisting}

\noindent If all we need to do is compare the ordinal values of the enumerators,
there's no problem:

\begin{emcppslisting}[emcppsbatch=e15]
bool isSummer(MonthOfYear month)
{
    return MonthOfYear::e_JUL <= month && month <= MonthOfYear::e_AUG;
}
\end{emcppslisting}

\noindent Although the \texttt{enum}~\texttt{class} features allow for relational
and equality operations between like-typed enumerators, no arithmetic
operations are supported directly, which becomes problematic when we
need to iterate over the enumerated values:

\begin{emcppslisting}[emcppsbatch=e15,emcppsignore={Entire for block is marked Error}]
void doSomethingWithEachMonth()
{
    for (MonthOfYear i =  MonthOfYear::e_JAN;
                     i <= MonthOfYear::e_DEC;
                   ++i)  // Error, no match for (ù{\codeincomments{++}}ù)
    {
        // ...
    }
}
\end{emcppslisting}

\noindent To make this code compile, an explicit cast from and to the enumerated
type will be required:

\begin{emcppslisting}[emcppsbatch=e15]
void doSomethingWithEachMonth()
{
    for (MonthOfYear i =  MonthOfYear::e_JAN;
                     i <= MonthOfYear::e_DEC;
                     i = static_cast<MonthOfYear>(static_cast<int>(i) + 1))
    {
        // ...
    }
}
\end{emcppslisting}

\noindent Alternatively, an auxiliary, helper function could be supplied to allow
clients to bump the enumerator:

\begin{emcppslisting}[emcppsbatch=e16]
MonthOfYear nextMonth(MonthOfYear value)
{
    return static_cast<MonthOfYear>(static_cast<int>(value) + 1);
}

void doSomethingWithEachMonth()
{
    for (MonthOfYear i =  MonthOfYear::e_JAN;
                     i <= MonthOfYear::e_DEC;
                     i = nextMonth(i))
    {
        // ...
    }
}
\end{emcppslisting}

\noindent If, however, the cardinal value of the \texttt{MonthOfYear} enumerators
is likely to be relevant to clients, an explicitly scoped \emph{classic}
\texttt{enum} should be considered as a viable alternative:

\begin{emcppslisting}
struct MonthOfYear  // explicit scoping for (ù{\codeincomments{enum}}ù)
{
    enum Enum
    {
        e_JAN, e_FEB, e_MAR,  // winter
        e_APR, e_MAY, e_JUN,  // spring
        e_JUL, e_AUG, e_SEP,  // summer
        e_OCT, e_NOV, e_DEC,  // autumn
    };
};

bool isSummer(MonthOfYear::Enum month)  // must now pass nested (ù{\codeincomments{Enum}}ù) type
{
    return MonthOfYear::e_JUL <= month && month <= MonthOfYear::e_AUG;
}

void doSomethingWithEachMonth()
{
    for (int i =  MonthOfYear::e_JAN;  // iteration variable is now an (ù{\codeincomments{int}}ù)
             i <= MonthOfYear::e_DEC;
           ++i)  // OK, convert to underlying type
    {
        // ... (might require cast back to enumerated type)
    }
}
\end{emcppslisting}

\noindent Note that such code presumes that the enumerated values will (1) remain
in the same order and (2) have contiguous numerical values irrespective
of the implementation choice.

\subsubsection[External use of opaque enumerators]{External use of opaque enumerators}\label{external-use-of-opaque-enumerators-enumclass}

Since \texttt{enum}~\texttt{class} types have an \textbf{underlying
type} of \texttt{int} by default, clients are always able to (re)declare
it, as a \textbf{complete type}, without its enumerators. Unless the
opaque form of an \texttt{enum}~\texttt{class}'s definition is exported
in a header file separate from the one implementing the publicly
accessible full definition, external clients wishing to exploit the
opaque version will experience an \emph{attractive nuisance} in that
they can provide it locally, along with its \textbf{underlying type}, if
any.

If the underlying type of the full definition were to subsequently
change, any program incorporating the original elided definition locally
and also the new, full one from the header would become silently \textbf{ill formed, no diagnostic required (IFNDR)}; see ``\titleref{enumopaque}" on page~\pageref{enumopaque}.

\subsection[Annoyances]{Annoyances}\label{annoyances-enumclass}

\subsubsection[Scoped enumerations do not necessarily add value]{Scoped enumerations do not necessarily add value}\label{scoped-enumerations-do-not-necessarily-add-value}

When the enumeration is local, say, within the scope of a given
function, forcing an additional scope on the enumerators is superfluous.
For example, consider a function that returns an integer status
\texttt{0} on success and a nonzero value otherwise:

\begin{emcppslisting}[emcppserrorlines={5,7}]
int f()
{
    enum { e_ERROR = -1, e_OK = 0 } result = e_OK;
    // ...
    if (/* error 1 */) { result = e_ERROR; }
    // ...
    if (/* error 2 */) { result = e_ERROR; }
    // ...
    return result;
}
\end{emcppslisting}

\noindent Use of \texttt{enum}~\texttt{class} in this context would require
potentially needless qualification --- and perhaps even casting ---
where it might not be warranted:

\begin{emcppslisting}[emcppserrorlines={5,7}]
int f()
{
    enum class RC { e_ERROR = -1, e_OK = 0 } result = RC::e_OK;
    // ...
    if (/* error 1 */) { result = RC::e_ERROR; } // undesirable qualification
    // ...
    if (/* error 2 */) { result = RC::e_ERROR; } // undesirable qualification
    // ...
    return static_cast<int>(result);  // undesirable explicit cast
}
\end{emcppslisting}

\subsection[See Also]{See Also}\label{see-also}

\begin{itemize}
\item{\seealsoref{explicit-enumeration-underlying-type}{\locationc}Absent implicit conversion to integrals, \texttt{enum} class values may use \texttt{static\_cast} in conjunction with their underlying type.}
\item{\seealsoref{enumopaque}{\locationc}Sometimes it is useful to entirely insulate individual enumerators from clients.}
\end{itemize}

\subsection[Further Reading]{Further Reading}\label{further-reading}

\begin{itemize}
\item{\cite{miller19}}
\item{\cite{meyers15}}
\end{itemize}




\newpage
\section[Opaque {\tt enum}s]{Opaque Enumeration Declarations}\label{enumopaque}\label{opaque-enumeration-declarations}
\input{opaqueenum}

\newpage
\section[Underlying Type '11]{Explicit Enumeration Underlying Type}\label{explicit-enumeration-underlying-type}
\input{underlyingenum}

\newpage
\section[{\tt friend} '11]{Extended {\SecCode friend} Declarations}\label{extended-friend-declarations}
\input{extendedfriend}

\newpage
\section[Forwarding References]{Forwarding {\SecCode \&\&} References}\label{forwardingref}\label{forwarding-references}
\input{forwardingref}

\newpage
\section[Generalized PODs]{generalized Plain Old Data Types}%\label{}
%\input{} 
PLACEHOLDER

\newpage
\section[{\tt initializer\_list}]{List Initialization: {\SecCode std::initializer\_list<T>}}%\label{}
%\input{}
PLACEHOLDER

\newpage
\section[Lambdas]{Unnamed Local Function Objects (Closures)}\label{lambda}
\input{lambda-placeholder}

\newpage
\section[{\tt noexcept} Operator]{The {\SecCode noexcept} Operator}%\label{noexcept}
%\input{noexcept-placeholder}
PLACEHOLDER

\newpage
\section[Range {\tt for}]{Range-Based {\SecCode for} Loops}%\label{}
%\input{}
PLACEHOLDER

\newpage
\section[{\itshape rvalue} References]{Rvalue References: {\SecCode \&\&}}\label{Rvalue-References}
\input{rvalueref-placeholder}

\newpage
\section[{\tt union} '11]{Unions Having Non-Trivial Members\sectionmark{{\RHCode union}~'11}}\label{unrestricted-unions}\sectionmark{{\RHCode union}~'11}


Any nonreference type is permitted to be a member of a \texttt{union}.

\subsection[Description]{Description}\label{unrestrictedunion-description}

Prior to C++11, only \textbf{trivial types}\glossary{trivial types} --- e.g.,
\textbf{fundamental types}, such as \texttt{int} and \texttt{double},
enumerated or pointer types, or a C-style array or \texttt{struct}
(a.k.a. a \textbf{POD}\glossary{POD}) --- were allowed to be members of a
\texttt{union}. This limitation prevented any (user-defined) type having
a \textbf{non-trivial special member function}\glossary{non-trivial special member function} from being a member of a
\texttt{union}:

\begin{lstlisting}[language=C++]
union U0
{
    int         d_i;  // OK
    std::string d_s;  // compile-time error in C++03 (OK as of C++11)
};
\end{lstlisting}
    
\noindent C++11 relaxes such restrictions on \texttt{union} members, such as
\texttt{d\_s} above, allowing any type other than a \textbf{reference
type}\glossary{reference
type} to be a member of a \texttt{union}.

A \texttt{union} type is permitted to have user-defined special member
functions but --- by design --- does not initialize any of its members
automatically. Any member of a \texttt{union} having a
\textbf{non-trivial constructor}\glossary{non-trivial constructor}, such as \texttt{struct}~\texttt{Nt}
(below), must be constructed manually (e.g., via \textbf{placement
\texttt{new}}\glossary{placement
\texttt{new}} implemented within the body of a constructor of the union
itself) before it can be used:

\begin{lstlisting}[language=C++]
struct Nt  // used as part of a (ù{\codeincomments{union}}ù) (below)
{
    Nt();   // non-trivial default constructor
    ~Nt();  // non-trivial destructor

    // Copy construction and assignment are implicitly defaulted.
    // Move construction and assignment are implicitly deleted.
};
\end{lstlisting}
    
\noindent As an added safety measure, any non-trivial \textbf{special member
function}\glossary{special member
function} defined --- either implicitly or explicitly --- for any
\emph{member} of a \texttt{union} results in the compiler implicitly
deleting (see Section~\ref{deleted-functions}, ``\titleref{deleted-functions}") the corresponding \textbf{special
member function}\glossary{special
member function} of the \texttt{union} itself:

\begin{lstlisting}[language=C++]
union U1
{
    int d_i;   // fundamental type having all trivial special member functions
    Nt  d_nt;  // user-defined type having non-trivial special member functions

    // Implicitly deleted special member functions of (ù{\codeincomments{U1}}ù):
    /*
        U1()                     = delete; // due to explicit (ù{\codeincomments{Nt::Nt()}}ù)
        U1(const U1&)            = delete; // due to implicit (ù{\codeincomments{Nt::Nt(const Nt\&)}}ù)
        ~U1()                    = delete; // due to explicit (ù{\codeincomments{Nt::}}ù)~(ù{\codeincomments{Nt()}}ù)
        U1& operator=(const U1&) = delete; // due to implicit 
                                           // (ù{\codeincomments{Nt::operator=(const Nt\&)}}ù)
    */
};
\end{lstlisting}
    
\noindent This same sort of precautionary deletion also occurs for any class
containing such a union as a data member (see {\it\titleref{unrestrictedunion-use-cases}: \titleref{implementing-a-sum-type-as-a-discriminating-(or-tagged)-union}} on page~\pageref{implementing-a-sum-type-as-a-discriminating-(or-tagged)-union}).

A special member function of a \texttt{union} that is implicitly deleted
can be restored via explicit declaration, thereby forcing a programmer
to think about how non-trivial members should be managed. For example,
we can start providing a \emph{value constructor} and corresponding
\emph{destructor}:

\begin{lstlisting}[language=C++]
struct U2
{
    union
    {
        int  d_i;   // fundamental type (trivial)
        Nt   d_nt;  // non-trivial user-defined type
    };

    bool d_useInt;  // discriminator

    U2(bool useInt) : d_useInt(useInt)       // value constructor
    {
        if (d_useInt) { new (&d_i) int(); }  // value initialized (to (ù{\codeincomments{0}}ù))
        else          { new (&d_nt) Nt(); }  // default constructed in place
    }

    ~U2()  // destructor
    {
        if (!d_useInt) { d_nt.~Nt(); }
    }
};
\end{lstlisting}
    
\noindent Notice that we have employed \textbf{placement \texttt{new}}\glossary{placement \texttt{new}} syntax to
control the lifetime of both member objects. Although assignment would
be permitted for the (trivial) \texttt{int} type, it would be
\textbf{undefined behavior}\glossary{undefined behavior} for the (non-trivial) \texttt{Nt} type:

\begin{lstlisting}[language=C++]
U2(bool useInt) : d_useInt(useInt)  // value constructor
{
    if (d_useInt) { d_i = int(); }  // value initialized (to (ù{\codeincomments{0}}ù))
    else          { d_nt = Nt(); }  // undefined behavior
}
\end{lstlisting}
    
\noindent Now if we were to try to copy-construct or assign an object of type
\texttt{U2} to another, the operation would fail because we have not
(yet) specifically addressed those \textbf{special member functions}\glossary{special member functions}:

\begin{lstlisting}[language=C++]
void f()
{
    U2 a(false), b(true);  // OK (construct both instances of (ù{\codeincomments{U2}}ù))
    U2 c(a);               // compile-time error: no (ù{\codeincomments{U2(const U2\&)}}ù)
    a = b;                 // compile-time error: no (ù{\codeincomments{U2\& operator=(const U2\&)}}ù)
}
\end{lstlisting}
    
\noindent We can restore these implicitly deleted special member functions too,
simply by adding appropriate copy-constructor and assignment-operator
definitions for \texttt{U2} explicitly{\cprotect\footnote{Attempting to
restore a \texttt{union}'s implicitly deleted special member
functions by using the \texttt{=}~\texttt{default} syntax (see
Section~\ref{Defaulted-Special-Member-Functions}, ``\titleref{Defaulted-Special-Member-Functions}") will still result in their being deleted because
the compiler cannot know which member of the union is active without a
  discriminator.}}:

\begin{lstlisting}[language=C++]
union U2
{
    // ... (everything in (ù{\codeincomments{U2}}ù) above)

    U2(const U2& original) : d_useInt(original.d_useInt)
    {
        if (d_useInt) { new (&d_i) int(original.d_i);  }
        else          { new (&d_nt) Nt(original.d_nt); }
    }

    U2& operator=(const U2& rhs)
    {
        if (this == &rhs) // Prevent self-assignment.
        {
            return *this;
        }

        // Resolve all possible combinations of active types between the 
        // left-hand side and right-hand side of the assignment:

        if (d_useInt)
        {
            if (rhs.d_useInt) { d_i = rhs.d_i; }
            else              { new (&d_nt) Nt(rhs.d_nt); }
        }
        else
        {
            if (rhs.d_useInt) { d_nt.~Nt(); new (&d_i) int(rhs.d_i); }
            else              { d_nt = rhs.d_nt; }
        }

        return *this;
    }
};
\end{lstlisting}
    

\subsection[Use Cases]{Use Cases}\label{unrestrictedunion-use-cases}

\subsubsection[Implementing a \textbf{sum type}\glossary{sum type} as a discriminating (or tagged) \tt{union}]{Implementing a \textbf{sum type}\glossary{sum type} as a discriminating (or tagged) {\ParaCode union}}\label{implementing-a-sum-type-as-a-discriminating-(or-tagged)-union}

A \textbf{sum type}\glossary{sum type} is an abstract data type that provides a choice
among a fixed set of specific types. Although other implementations are
possible, using the storage of a single object to accommodate one out of
a set of types along with a (typically integral) discriminator enables
programmers to implement a \textbf{sum type}\glossary{sum type} (a.k.a.
\emph{discriminating} or \emph{tagged} union) efficiently (e.g., without
necessarily involving memory allocation or virtual dispatch) and
nonintrusively (i.e., the individual types comprised need not be related
in any way). A C++ \texttt{union} can serve as a convenient and
efficient way to define storage for a \textbf{sum type}\glossary{sum type} as alignment and
size calculations are performed (by the compiler) automatically.

As an example, consider writing a parsing function \texttt{parseInteger}
that, given a\linebreak[4] \texttt{std::string} \texttt{input}, will return, as a
\textbf{sum type}\glossary{sum type} \texttt{ParseResult} (see below), either an
\texttt{int} result (on success) or an informative error message
(on failure):

\begin{lstlisting}[language=C++]
ParseResult parseInteger(const std::string& input)  // Return a sum type.
{
    int result;     // accumulate (ù{\codeincomments{result}}ù) as we go
    std::size_t i;  // current character index

    // ...

    if (/* Failure case (1). */)
    {
        std::ostringstream oss;
        oss << "Found non-numerical character '" << input[i]
            << "' at index '" << i << "'.";

        return ParseResult(oss.str());
    }

    if (/* Failure case (2). */)
    {
        std::ostringstream oss;
        oss << "Accumulating '" << input[i]
            << "' at index '" << i
            << "' into the current running total '" << result
            << "' would result in integer overflow.";

        return ParseResult(oss.str());
    }

    // ...

    return ParseResult(result);  // Success!
}
\end{lstlisting}
    
\noindent The implementation above relies on \texttt{ParseResult} being able to
hold a value of type either \texttt{int} or \texttt{std::string}. By
encapsulating a C++ \texttt{union} and a Boolean{\cprotect\footnote{For
\textbf{sum types}\glossary{sum types} comprising more than two types, a larger integral
  or enumerated type may be used instead.}} \emph{discriminator} as part
of the \texttt{ParseResult} \textbf{sum type}\glossary{sum type}, we can achieve the
desired semantics:

\begin{lstlisting}[language=C++]
class ParseResult
{
    union  // storage for either the result or the error
    {
        int         d_value;  // trivial result type
        std::string d_error;  // non-trivial error type
    };

    bool d_isError;  // discriminator

public:
    explicit ParseResult(int value);                 // value constructor (1)
    explicit ParseResult(const std::string& error);  // value constructor (2)

    ParseResult(const ParseResult& rhs);             // copy constructor
    ParseResult& operator=(const ParseResult& rhs);  // copy assignment

    ~ParseResult();                                  // destructor
};
\end{lstlisting}
    
\noindent As discussed in {\it\titleref{unrestrictedunion-description}} on page~\pageref{unrestrictedunion-description}, having a non-trivial
type within a \texttt{union} forces the programmer to provide each
desired special member function and define it manually; note, although,
that the use of placement \texttt{new} is not required for either of the
two \emph{value constructors} (above) because the initializer syntax
(below) is sufficient to begin the lifetime of even a non-trivial
object:

\begin{lstlisting}[language=C++]
ParseResult::ParseResult(double value) : d_value(value), d_isError(false)
{
}

ParseResult::ParseResult(const std::string& error) 
    : d_error(error), d_isError(true)
    // Note that placement (ù{\codeincomments{new}}ù) was not necessary here because a new 
    // (ù{\codeincomments{std::string}}ù) object will be created as part of the initialization of 
    // (ù{\codeincomments{d\_error}}ù).
{
}
\end{lstlisting}
    
\noindent Placement \texttt{new} and explicit destructor calls are, however,
required for destruction and both copy operations{\cprotect\footnote{For
more information on initiating the lifetime of an object, see \textbf{iso14}, section 3.8, ``Object Lifetime," pp. 66--69.}}:

\begin{lstlisting}[language=C++]
ParseResult::~ParseResult()
{
    if(d_isError)
    {
        d_error.std::string::~string();
            // An explicit destructor call is required for (ù{\codeincomments{d\_error}}ù) because its
            // destructor is non-trivial.
    }
}

ParseResult::ParseResult(const ParseResult& rhs) : d_isError(rhs.d_isError)
{
    if (d_isError)
    {
        new (&d_error) std::string(rhs.d_error);
            // Placement (ù{\codeincomments{new}}ù) is necessary here to begin the lifetime of a
            // (ù{\codeincomments{std::string}}ù) object at the address of (ù{\codeincomments{d\_error}}ù).
    }
    else
    {
        d_value = rhs.d_value;
            // Placement (ù{\codeincomments{new}}ù) is not necessary here as (ù{\codeincomments{int}}ù) is a trivial type.
    }
}

ParseResult& ParseResult::operator=(const ParseResult& rhs)
{
    // Destroy (ù{\codeincomments{lhs}}ù)'s error string if existent:
    if (d_isError) { d_error.std::string::~string(); }

    // Copy (ù{\codeincomments{rhs}}ù)'s object:
    if (rhs.d_isError) { new (&d_error) std::string(rhs.d_error); }
    else               { d_value = rhs.d_value; }

    d_isError = rhs.d_isError;
    return *this;
}
\end{lstlisting}
    
\noindent In practice, \texttt{ParseResult} would typically be defined as a
template and renamed to allow any arbitrary result type \texttt{T} to be
returned or else implemented in terms of a more general \textbf{sum
type}\glossary{sum type} abstraction.{\cprotect\footnote{\texttt{std::variant}, introduced
in C++17, is the standard construct used to represent a \textbf{sum
type}\glossary{sum type} as a \emph{discriminating union}. Prior to C++17,
\texttt{boost::variant} was the most widely used \emph{tagged} union
  implementation of a \textbf{sum type}\glossary{sum type}.}}

\subsection[Potential Pitfalls]{Potential Pitfalls}\label{potential-pitfalls}

\subsubsection[Inadvertent misuse can lead to latent \textbf{undefined behavior}\glossary{undefined behavior} at runtime]{Inadvertent misuse can lead to latent \textbf{undefined behavior}\glossary{undefined behavior} at runtime}\label{inadvertent-misuse-can-lead-to-latent-undefined-behavior-at-runtime}

When implementing a type that makes use of an unrestricted union,
forgetting to initialize a non-trivial object (using either a
\emph{member initialization list} or \textbf{placement \texttt{new}}\glossary{placement \texttt{new}}) or
accessing a different object than the one that was actually initialized
can result in tacit \textbf{undefined behavior}\glossary{undefined behavior}. Although forgetting to
destroy an object does not necessarily result in \textbf{undefined
behavior}\glossary{undefined behavior}, failing to do so for any object that manages a resource (such
as dynamic memory) will result in a \emph{resource leak} and/or lead to
unintended behavior. Note that destroying an
object having a trivial destructor is never necessary; there are, however, rare cases where
we may choose not to destroy an object having a non-trivial
one.{\cprotect\footnote{A specific example of where one might
deliberately choose \emph{not} to destroy an object occurs when a
collection of related objects are allocated from the same local memory
resource and then deallocated unilaterally by releasing the memory
back to the resource. No issue arises if the only resource that is ``leaked''
by not invoking each individual destructor is the memory allocated
from that memory resource, and that memory can be
reused without resulting in \textbf{undefined behavior}\glossary{undefined behavior} if it
is not subsequently referenced in the context of the deallocated
  objects.}}

\subsection[Annoyances]{Annoyances}\label{annoyances}

None so far

\subsection[See Also]{See Also}\label{see-also}

\begin{itemize}
\item{Section~\ref{deleted-functions}, ``\titleref{deleted-functions}" — Safe C++11 feature that forbids the invocation of a particular function. Similar effects to deleting a function happen when we specify a special function within a subobject of a union or when a class has such a union as a data member.}
\end{itemize}

\subsection[Further Reading]{Further Reading}\label{further-reading}

None so far




\newpage
\section[User-Defined Literals]{User-Defined Literal Operators}%\label{}
%\input{}
PLACEHOLDER

\newpage
\section[Variadic Templates]{Variable-Argument-Count Templates}\label{variadictemplate}
\input{variadictemplate-placeholder}


%%%%%%%%%%%%%%% C++14
\stepcounter{cppxx}
\addtocontents{toc}{\protect\renewcommand*\protect\cftsecindent{1.25em}}
\cftaddnumtitleline{toc}{section}{\thechapter.\thecppxx}{C++14}{\thepage}
% to add to the TOC at the section level
\renewcommand{\cppxx}{C++14}


\newpage
\addtocontents{toc}{\protect\renewcommand*\protect\cftsecindent{4.5em}}
\section[{\itshape Generic} Lambdas]{Lambdas Having a Templated Call Operator}\label{genericlambda}



\emcppsFeature{
    short={{\itshape Generic} Lambdas},
    long={Lambdas Having a Templated Call Operator},
}{genericlambda}
\setcounter{table}{0}
\setcounter{footnote}{0}
\setcounter{lstlisting}{0}
%\section[{\itshape Generic} Lambdas]{Lambdas Having a Templated Call Operator}\label{genericlambda}



placeholder text.........
 
 


\newpage
\addtocontents{toc}{\protect\renewcommand*\protect\cftsecindent{4.5em}}
\section[Short Title]{Longer Wordier Title for Test}\label{justfortest}

Section text; just to test where/how the C++ change (from C++11 to C++14) in the RH must be placed in relation to sectioning commands such that the correct one displays

\begin{table}[h!] 
\begin{center}
\begin{threeparttable}
\caption{Testing Table Numbering}\label{test-table3}\vspace{1.5ex} 
{\small \begin{tabular}{c|c|c}\thickhline 
\rowcolor[gray]{.9}   {\sffamily\bfseries Compiler} & {\sffamily\bfseries Compiler-Specific} &
{\sffamily\bfseries Standard-Conforming} \\ \hline 
GCC &\texttt{\_\_attribute\_\_((pure))} & \texttt{[[gnu::pure]]} \\ \hline
Clang & \texttt{\_\_attribute\_\_((no\_sanitize))} &\texttt{[[clang::no\_sanitize]]} \\ \hline 
MSVC & \texttt{declspec(deprecated)} & \texttt{[[deprecated]]} \\ \thickhline
\end{tabular}
}
\end{threeparttable} 
\end{center}
\end{table}

Donec mattis, ex a ornare auctor, tortor felis pulvinar odio, eu mollis tortor orci id ligula. Vestibulum dignissim magna vitae lectus cursus, ac sollicitudin mi pulvinar. Praesent vulputate lorem lectus, eu pulvinar mi finibus eu. Donec sit amet massa massa. Sed eget molestie quam.\footnote{testing a footnote} Praesent iaculis diam ut nunc condimentum auctor. Morbi porttitor varius ante vel venenatis. Integer ut porttitor elit, id lacinia turpis. Integer id iaculis ante, placerat ornare tortor. Sed et consectetur est, sit amet rhoncus turpis.

\begin{table}[h!] 
\begin{center}
\begin{threeparttable}
\caption{Testing Table Numbering}\label{test-table4}\vspace{1.5ex} 
{\small \begin{tabular}{c|c|c}\thickhline 
\rowcolor[gray]{.9}   {\sffamily\bfseries Compiler} & {\sffamily\bfseries Compiler-Specific} &
{\sffamily\bfseries Standard-Conforming} \\ \hline 
GCC &\texttt{\_\_attribute\_\_((pure))} & \texttt{[[gnu::pure]]} \\ \hline
Clang & \texttt{\_\_attribute\_\_((no\_sanitize))} &\texttt{[[clang::no\_sanitize]]} \\ \hline 
MSVC & \texttt{declspec(deprecated)} & \texttt{[[deprecated]]} \\ \thickhline
\end{tabular}
}
\end{threeparttable} 
\end{center}
\end{table}

This feature, when used in conjunction with \emph{explicit instantiation
definitions}, can significantly improve compilation times for a set of
translation units that often instantiate common templates:

\noindent\begin{minipage}{.45\textwidth}
\begin{lstlisting}[language=C++, caption=code 1,frame=tb]{Name}
void code()
{

}
\end{lstlisting}
\end{minipage}\hfill
\begin{minipage}{.45\textwidth}
\begin{lstlisting}[language=C++, caption=code 2,frame=tb]{Name}
void code()
{

}
\end{lstlisting}
\end{minipage}

Nullam nibh tortor, finibus ut lectus eu, convallis vehicula libero. Cras maximus ligula nisl, a eleifend mauris venenatis nec. Suspendisse potenti. Maecenas viverra laoreet mauris ut laoreet. Donec non felis risus. Ut faucibus, justo id luctus lobortis, nisi lacus pharetra est, id pretium arcu ligula ac magna. Morbi a nisl sit amet lacus vehicula tempor nec sit amet nibh. 

\noindent\begin{minipage}{.45\textwidth}
\begin{lstlisting}[language=C++, caption=code 1 with a long and wrapping title,frame=tb]{Name}
void code()
{

}
\end{lstlisting}
\end{minipage}\hfill
\begin{minipage}{.45\textwidth}
\begin{lstlisting}[language=C++, caption=code 2 with a long and wrapping title,frame=tb]{Name}
void code()
{

}
\end{lstlisting}
\end{minipage}

Pellentesque id lorem ac sem ullamcorper semper. Donec imperdiet sapien in nisi faucibus, et tempus lacus hendrerit. Nulla laoreet risus eu tortor ultrices, sed bibendum neque sodales. Ut faucibus ipsum id convallis convallis. Vivamus vehicula rutrum metus in semper. Morbi fringilla ex vel vulputate vehicula. Maecenas ut porttitor massa.





  % Chapter 2
\cleardoublepage
%\addtocontents{toc}{\protect\renewcommand\cftsecindent{1.25em}} %%%%%
%\addtocontents{toc}{\protect\renewcommand*\protect\cftsubsecindent{4em}}

\chapter[Unsafe Features]{Unsafe Features}\label{ch-unsafe}

%%%%%%%%%%%%%%% C++11
\stepcounter{cppxx}
\addtocontents{toc}{\protect\renewcommand*\protect\cftsecindent{1.25em}}
\cftaddnumtitleline{toc}{section}{\thechapter.\thecppxx}{C++11}{\thepage}
% to add to the TOC at the section level
\renewcommand{\cppxx}{C++11}

\newpage
\addtocontents{toc}{\protect\renewcommand*\protect\cftsecindent{4.5em}}
%\cftaddtitleline{toc}{section}{\protect\hspace*{-5em}\framebox{\sffamily\itshape\cppxx}\vspace*{-3ex}}{}
%\section[The {\tt [[carries\_dependency]]} Attribute]{The {\SecCode [[carries\_dependency]]} Attribute\sectionmark{{\tt [[carries\_dependency]]}}\checkoddpage\Cbox}\label{carriesdependency}
 \section[{\tt [[carries\_dependency]]}\hspace*{1em}(The {\tt [[carries\_dependency]]} Attribute)]{The {\SecCode [[carries\_dependency]]} Attribute\sectionmark{{\tt [[carries\_dependency]]}}}\label{carriesdependency}\sectionmark{{\tt [[carries\_dependency]]}}
\input{carriesdependency-placeholder}



 %%%%%%%%%%%%%%% C++14 starts
\stepcounter{cppxx}
\addtocontents{toc}{\protect\renewcommand*\protect\cftsecindent{1.25em}}
\cftaddnumtitleline{toc}{section}{\thechapter.\thecppxx}{C++14}{\thepage}
% to add to the TOC at the section level
\renewcommand{\cppxx}{C++14}
%%%%%%%%%%%%%%

\newpage
\addtocontents{toc}{\protect\renewcommand*\protect\cftsecindent{4.5em}}
%\cftaddtitleline{toc}{section}{\protect\hspace*{-5em}\framebox{\sffamily\itshape\cppxx}\vspace*{-3ex}}{}
%\section[Function Return Type Deduction]{Function Return Type Deduction\sectionmark{Deduced Return Types}\checkoddpage\Cbox}\label{Function-Return-Type-Deduction}
 \section[Deduced Return Types\hspace*{1em}(Function Return Type Deduction)]{Function Return Type Deduction\sectionmark{Deduced Return Types}}\label{Function-Return-Type-Deduction}
\renewcommand{\cppxx}{C++14}

\input{autoreturn-placeholder}

  % Chapter 3
\cleardoublepage

%
%%%%%% RHs revert for Ch 4 + BM
\fancyhead[LO]{}
%\renewcommand{\sectionmark}[1]{\markright{{\sfnine\thesection\ #1}}}
%\renewcommand{\chaptermark}[1]{\markboth{{\sfb\thechapter\quad {\sfb #1}}}{}}
%%%%%%%%%%


\chapter[Parting Thoughts]{Parting Thoughts}\label{ch-parting}
\emcppschapterstart{}

\addtocontents{toc}{\protect\renewcommand*\protect\cftsecindent{1.25em}}



  % Chapter 4
\cleardoublepage
%
%
\cleardoublepage
%\bibliographystyle{ieeetr-hacked}
\bibliographystyle{lakos}
%\bibliographystyle{abstract}
\bibliography{emcs}
\nocite{*}
%
\chapter*{Bibliography}

\noindent\textbf{ballman}\\
\hspace*{2em}\parbox[t]{\textwidth}{Aaron Ballman. ``Rule 03. Integers (INT)," \textit{ SEI CERT C++ Coding Standard} (Pittsburgh, PA: Carnegie Mellon University Software Engineering Institute)\\
https://wiki.sei.cmu.edu/confluence/pages/viewpage.action?pageId=88046333}\\

\noindent\textbf{calabrese20}\\
\hspace*{2em}\parbox[t]{\textwidth}{Matt Calabrese and Ryan McDougall. ``{\ttfamily any\_invocable}."  Technical Report P0288R6. (Geneva, Switzerland: International Standards Organization, 2020)\\
http://www.open-std.org/jtc1/sc22/wg21/docs/papers/2020/p0288r6.html}\\

\noindent\textbf{cpp11}\\
\hspace*{2em}\parbox[t]{\textwidth}{\textit{Information Technology — Programming
Languages — C++} (Geneva, Switzerland: International Standards Organization, 2011)\\
http://www.open-std.org/jtc1/sc22/wg21/docs/papers/2011/n3242.pdf}\\

\noindent\textbf{cppreferencea}\\
\hspace*{2em}\parbox[t]{\textwidth}{``std::all\_of, std::any\_of, std::none\_of," \textit{C++ Algorithm Library}, cppreference.com\\
https://en.cppreference.com/w/cpp/algorithm/all\_any\_none\_of}\\

\noindent\textbf{cppreferenceb}\\
\hspace*{2em}\parbox[t]{\textwidth}{``std::unordered\_map," \textit{C++ Containers Library}, cppreference.com\\
https://en.cppreference.com/w/cpp/container/unordered\_map}\\

\noindent\textbf{freesoftwarefdn20}\\
\hspace*{2em}\parbox[t]{\textwidth}{\textit{Using the GNU Compiler Collection (GCC)} (Boston, MA: Free Software Foundation, Inc., 2020)\\
https://gcc.gnu.org/onlinedocs/gcc/\\
\textbf{RETAIN IN CASE WE NEED THIS; USED IN ATTRIBUTES FEATURE}\qquad  section~6.33.1, ``Common Function Attributes," \\
https://gcc.gnu.org/onlinedocs/gcc/Common-Function-Attributes.html\#Common-Function-Attribut}\\

\noindent\textbf{hinnant02}\\
\hspace*{2em}\parbox[t]{\textwidth}{Howard Hinnant. \textit{A Proposal to Add Move Semantics Support to the C++ Language} (Geneva, Switzerland: International Standards Organization, 2002)\\
http://www.open-std.org/jtc1/sc22/wg21/docs/papers/2002/n1377.htm}\\

\noindent\textbf{hinnant14}\\
\hspace*{2em}\parbox[t]{\textwidth}{Howard Hinnant. ``Everything You Ever Wanted To Know About Move Semantics (and Then Some)" \textit{Conference of the ACCU}, Bristol, England, April 8--12, 2014\\
https://accu.org/content/conf2014/Howard\_Hinnant\_Accu\_2014.pdf}\\

\noindent\textbf{hinnant16}\\
\hspace*{2em}\parbox[t]{\textwidth}{Howard Hinnant. ``Everything You Ever Wanted to Know About Move Semantics" \textit{Bloomberg Engineering Distinguished Speaker Series}, New York, NY, 2016\\
https://www.youtube.com/watch?v=vLinb2fgkHk\&t=28s }\\

\noindent\textbf{hruska20}\\
\hspace*{2em}\parbox[t]{\textwidth}{Joel Hruska. ``How L1 and L2 CPU Caches Work, and Why They’re an Essential Part of Modern Chips," \textit{Extreme Tech}. April 14, 2020.\\
https://www.extremetech.com/extreme/188776-how-l1-and-l2-cpu-caches-work-and-why-theyre-an-essential-part-of-modern-chips}\\

\noindent\textbf{ieee19}\\
\hspace*{2em}\parbox[t]{\textwidth}{\textit{IEEE Standard for Floating-Point Arithmetic} (New York, NY: Institute of Electrical and Electronics Engineers, Inc., 2019)}\\

\noindent\textbf{intel16}\\
\hspace*{2em}\parbox[t]{\textwidth}{Intel Corporation. \textit{Intel 64 and IA-32 Architectures Optimization Reference Manual}, Publication Number 248966-033. (Santa Clara, CA: Intel Corporation, 2016)\\
https://www.intel.com/content/dam/www/public/us/en/documents/manuals/64-ia-32-architectures-optimization-manual.pdf}\\

\noindent\textbf{inteliig}\\
\hspace*{2em}\parbox[t]{\textwidth}{Intel. \textit{The Intel Intrinsics Guide} (Santa Clara, CA: Intel Corporation)\\
https://software.intel.com/sites/landingpage/IntrinsicsGuide/\#}\\

\noindent\textbf{iso14}\\
\hspace*{2em}\parbox[t]{\textwidth}{\textit{ISO/IEC 14882:2014 Programming Language C++} (Geneva, Switzerland: International Standards Organization, 2014)\\ 
http://www.open-std.org/JTC1/SC22/WG21/docs/papers/2013/n3797.pdf}\\

\noindent\textbf{iso18a}\\
\hspace*{2em}\parbox[t]{\textwidth}{\textit{C++ Standard Core Language Active Issues, Revision 100} (Geneva, Switzerland: International Standards Organization, 2018)\\
http://www.open-std.org/jtc1/sc22/wg21/docs/cwg\_active.html}\\

\noindent\textbf{iso18b}\\
\hspace*{2em}\parbox[t]{\textwidth}{\textit{ISO/IEC 9899:2018 XXXXXXXXXXXXXXXXX} (Geneva, Switzerland: International Standards Organization, 2018)\\ 
URL}\\

\noindent\textbf{iso20a}\\
\hspace*{2em}\parbox[t]{\textwidth}{\textit{ISO/IEC P2156R1:2020 Allow Duplicate Attributes} (Geneva, Switzerland: International Standards Organization, 2020\\
http://www.open-std.org/jtc1/sc22/wg21/docs/papers/2020/p2156r1.pdf}\\

\noindent\textbf{iso20b}\\
\hspace*{2em}\parbox[t]{\textwidth}{\textit{ISO/IEC 14882:2020 XXXXXXXXXXXXXXXXX} (Geneva, Switzerland: International Standards Organization, 2020)\\ 
URL}\\

\noindent\textbf{kahan97}\\
\hspace*{2em}\parbox[t]{\textwidth}{W. Kahan, \textit{Lecture Notes on the Status of IEEE Standard 754 for Binary Floating-Point Arithmetic} (Berkeley, CA: Electrical Engineering and Computer Science Department, University of California, 1997)\\
https://people.eecs.berkeley.edu/~wkahan/ieee754status/IEEE754.PDF }\\

\noindent\textbf{lakos16}\\
\hspace*{2em}\parbox[t]{\textwidth}{John Lakos, Jeffrey Mendelsohn, Alisdair Meredith, and Nathan Myers, ``On Quantifying Memory-Allocation Strategies (Revision 2),"  P0089R1, February 12, 2016\\
http://www.open-std.org/jtc1/sc22/wg21/docs/papers/2016/p0089r1.pdf}\\

\noindent\textbf{lakos17}\\
\hspace*{2em}\parbox[t]{\textwidth}{John Lakos, ``Local (`Arena') Memory Allocators - Part II," \textit{Meeting C++}, 2017 \\
https://www.youtube.com/watch?v=fN7nVzbRiEk}\\

\noindent\textbf{lakos19}\\
\hspace*{2em}\parbox[t]{\textwidth}{John Lakos, ``Value Proposition: Allocator-Aware (AA) Software," \textit{C++Now}, 2019\\
https://www.youtube.com/watch?v=dDR93TfacHc}\\

\noindent\textbf{lakos20}\\
\hspace*{2em}\parbox[t]{\textwidth}{John Lakos, \textit{Large-Scale C++ --- Volume I: Process and Architecture} (Reading, MA: Addison-Wesley, 2019)}\\

\noindent\textbf{lakos22}\\
\hspace*{2em}\parbox[t]{\textwidth}{John Lakos and Joshua Berne, \textit{C++ Allocators for the Working Programmer}, manuscript} \\

\noindent\textbf{lakos23}\\
\hspace*{2em}\parbox[t]{\textwidth}{John Lakos and Rostislava Klebnikhov. \textit{title}, manuscript} \\

\noindent\textbf{lakos2a}\\
\hspace*{2em}\parbox[t]{\textwidth}{John Lakos. \textit{Large-Scale C++ --- Volume II: XXXXXXXXXX}, manuscript} \\

\noindent\textbf{lakos2b}\\
\hspace*{2em}\parbox[t]{\textwidth}{John Lakos. \textit{Large-Scale C++ --- Volume II: XXXXXXXXXX} manuscript} \\

\noindent\textbf{lakos96}\\
\hspace*{2em}\parbox[t]{\textwidth}{John Lakos, \textit{Large-Scale C++ Software Design} (Reading, MA: Addison-Wesley, 1996)} \\

\noindent\textbf{meredith08}\\
\hspace*{2em}\parbox[t]{\textwidth}{A. Meredith, M. Wong, and J. Maurer. ``Inheriting Constructors (Revision 4),” N2512, 02 April 2008,  (Geneva, Switzerland: International Standards Organization, 2008)\\
http://www.open-std.org/jtc1/sc22/wg21/docs/papers/2008/n2512.html}\\

\noindent\textbf{meyers97}\\
\hspace*{2em}\parbox[t]{\textwidth}{Scott Meyers. \textit{Effective C++}, second ed. (Boston, MA: Addison Wesley, 1997)} \\

\noindent\textbf{microsoft}\\
\hspace*{2em}\parbox[t]{\textwidth}{Microsoft. \textit{Guidelines Support Library}.\\
https://github.com/Microsoft/GSL}\\

\noindent\textbf{microsoftC26481}\\
\hspace*{2em}\parbox[t]{\textwidth}{Microsoft. ``C26481 NO\_POINTER\_ARITHMETIC," \textit{Guidelines Support Library}, 2020.
https://docs.microsoft.com/en-us/cpp/code-quality/c26481?view=vs-2019}\\

\noindent\textbf{miller13}\\
\hspace*{2em}\parbox[t]{\textwidth}{Mike Miller. ``Issue 1655: Line Endings in Raw String Literals,” 26 April 2013, included within \textit{C++ Standard Core Language Active Issues, Revision 100} (Geneva, Switzerland: International Standards Organization, 2018)\\
http://www.open-std.org/jtc1/sc22/wg21/docs/cwg\_active.html\#1655}

\noindent\textbf{miller17}\\
\hspace*{2em}\parbox[t]{\textwidth}{Mike Miller. ``Issue 2354: Extended Alignment and Object Representation,” 05 September 2017, included within \textit{Core Language Working Group Tentatively Ready Issues for the February, 2019 (Kona) Meeting} (Geneva, Switzerland: International Standards Organization, 2019)\\
http://www.open-std.org/jtc1/sc22/wg21/docs/papers/2019/p1359r0.html\#2354}\\

\noindent\textbf{meyers05}\\
\hspace*{2em}\parbox[t]{\textwidth}{Scott Meyers. \textit{Effective C++} (Boston, MA: Addison-Wesley, 2005)}\\

\noindent\textbf{xxxxxxxxxx}\\
\hspace*{2em}\parbox[t]{\textwidth}{\textit{}p2156r0\\
https://wg21.link/p2156r0}\\

\noindent\textbf{pusz20}\\
\hspace*{2em}\parbox[t]{\textwidth}{Mateusz Pusz.\textit{Enable variable template template Parameters}, Technical Report P2008R0 (Geneva, Switzerland: C++ Standards Committee Working Group, 2005)}\\

\noindent\textbf{solihin15}\\
\hspace*{2em}\parbox[t]{\textwidth}{Yan Solihin. \textit{Fundamentals of Parallel Multi-Core Architecture} (Boca Raton, FL: CRC Press, 2015)}\\

\noindent\textbf{seacord13}\\
\hspace*{2em}\parbox[t]{\textwidth}{Robert C. Seacord. \textit{Secure Coding in C and C++}, 2nd ed. (Boston, MA: Addison Wesley, 2013)}\\

\noindent\textbf{sharpe13}\\
\hspace*{2em}\parbox[t]{\textwidth}{Chris Sharpe. ``Contextually converted to bool," \textit{Chris's C++ Thoughts} (blog). (Blogspot, July 28, 2013)\\
http://chris-sharpe.blogspot.com/2013/07/contextually-converted-to-bool.html\#:{$\sim$}:\linebreak[4]text$=$Certain$\%$20language$\%$20constructs$\%$20require$\%$20that,temporary$\%$20variable\linebreak[4]$\%$20t$\%$20(8.5)}\\

\noindent\textbf{smith15a}\\
\hspace*{2em}\parbox[t]{\textwidth}{Richard Smith. \textit{Attributes for namespaces and enumerators},  Technical Report N4196 (Geneva, Switzerland: C++ Standards Committee Working Group, 2015)}\\

\noindent\textbf{smith18}\\
\hspace*{2em}\parbox[t]{\textwidth}{Richard Smith. \textit{Base Class Copy and Move Constructors Should Not Be Inherited}, 26 Feb 2018, Technical Report CWG2356 (Geneva, Switzerland: C++ Standards Committee Working Group, 2018)}\\

\noindent\textbf{smith15b}\\
\hspace*{2em}\parbox[t]{\textwidth}{Richard Smith. \textit{Rewording Inheriting Constructors (Core Issue 1941 et al.)},  19 Oct 2020, Technical Report P0136R1 (Geneva, Switzerland: C++ Standards Committee Working Group, 2015)}\\

\noindent\textbf{stevens93}\\
\hspace*{2em}\parbox[t]{\textwidth}{W. Richard Stevens. \textit{Advanced Programming in the UNIX Environment} (Boston, MA: Addison-Wesley, 1993)} \\

\noindent\textbf{stroustrup13}\\
\hspace*{2em}\parbox[t]{\textwidth}{Bjarne Stroustrup. \textit{The C++ Programming Language}, 4th ed. (Boston, MA: Addison-Wesley, 2013)} \\

\noindent\textbf{stroustrup20}\\
\hspace*{2em}\parbox[t]{\textwidth}{Bjarne Stroustrup and Herb Sutter, Eds. \textit{C++ Core Guidelines} (\textbf{CITY}, WA: Standard C++ Foundation), 2020\\
https://isocpp.github.io/CppCoreGuidelines/CppCoreGuidelines} \\

\noindent\textbf{vandevoorde05}\\
\hspace*{2em}\parbox[t]{\textwidth}{Daveed Vandevoorde. \textit{Right Angle Brackets}, Technical Report N1757, Revision 2, (Geneva, Switzerland: C++ Standards Committee Working Group, 2005)\\
http://www.open-std.org/jtc1/sc22/wg21/docs/papers/2005/n1757.html}\\

\noindent\textbf{wight}\\
\hspace*{2em}\parbox[t]{\textwidth}{Hyrum Wight. \textit{Hyrum's Law}.\\
https://www.hyrumslaw.com/}\\






\cleardoublepage
\cleardoublepage % This command ensures that each chapter starts on a right-hand page 

\chapter*{Glossary}\label{Glossary}

% A new command, gloss, handles the glossary entry. \gloss{word to be defined}{definition}

\gloss{ABI}{}

\gloss{access modifier}{}

\gloss{access specifiers}{}

\gloss{acquire/release memory barrier}{}

\gloss{aggregate}{}

\gloss{aggregate initialization}{}

\gloss{algebra}{}

\gloss{alias template}{}

\gloss{alignment}{}

\gloss{alignment requirements}{}

\gloss{API}{}

\gloss{argument-dependent lookup (ADL)}{}

\gloss{arithmetic type}{}

\gloss{array type}{}

\gloss{atomic}{}

\gloss{automatic variables}{}

\gloss{barriers}{}

\gloss{basic source character set}{}

\gloss{benchmark test}{}

\gloss{B0}{}

\gloss{boilerplate code}{}

\gloss{brace elision}{}

\gloss{bytes}{}

\gloss{C-style functions}{}

\gloss{cache hit}{}

\gloss{cache line}{}

\gloss{cache miss}{}

\gloss{callback functions}{}

\gloss{callable object}{}

\gloss{cast}{}

\gloss{CI}{}

\gloss{Class member access expression}{TODO (include any expression that is used to refer to a class member, such as \texttt{object.member}, \texttt{object->member}, \texttt{object.*member}.)}

\gloss{closures}{}

\gloss{code bloat}{}

\gloss{compile-time constant}{}

\gloss{compile-time coupling}{}

\gloss{Compile-time dispatch}{TODO}

\gloss{complete type}{}

\gloss{component}{}

\gloss{component-local}{}

\gloss{Concepts}{}

\gloss{concrete class}{}

\gloss{conditionally supported}{}

\gloss{Constant expression}{An expression that can be evaluated at compile time.  Mention \texttt{constexpr} and state that \texttt{const} variables that are initialized from a compile-time constants are themselves required to be compile-time constants.  New info to me in June 2020, worthwhile to have. }

\gloss{constant initialization}{}

\gloss{contextual convertibility to {\tt bool}}{}

\gloss{Contextual keyword}{A \textit{contextual keyword} is a special identifier that acts like a \textit{keyword} when used in particular contexts. \texttt{override} is an example as it can be used as a regular identifier outside of member-function declarators.}

\gloss{continuous refactoring}{}

\gloss{contract}{}

\gloss{conventional string literals}{}

\gloss{converting constructors}{}

\gloss{conversion operators}{}

\gloss{copy elision}{}

\gloss{copy initialization}

\gloss{Copy semantics}{TODO}

\gloss{critical section}{}

\gloss{\textbf{curiously recurring template pattern (CRTP)}}{}

\gloss{cv-qualifiers}{}

\gloss{cyclically dependent}{}

\gloss{data member}{}

\gloss{data races}{}

\gloss{Declaration}{TODO}

\gloss{declared interface}{}

\gloss{Declared type}{TODO The type of the *entity* named by the given expression.}

\gloss{declaring}{}

\gloss{defensive}{}

\gloss{defining}{}

\gloss{definition}{}

\gloss{delegating constructor}{}

\gloss{deleted function}{}

\gloss{dependent type}{}

\gloss{design pattern}{}

\gloss{diffusion}{}

\gloss{direct initialization}{}

\gloss{direct mapped}{}

\gloss{duck typing}{}

\gloss{dumb data}{}

\gloss{emplacement}{}

\gloss{encapsulation}{}

\gloss{Entity}{TODO}

\gloss{excess-\textit{n}}{}

\gloss{Expression SFINAE}{TODO}

\gloss{explicit}{}

\gloss{{\tt explicit} specifier}{}

\gloss{extended alignment}{}

\gloss{false sharing}{}

\gloss{fences}{}

\gloss{flow of control}{}

\gloss{forward declaration}{}

\gloss{forward declared}{}

\gloss{forwarding reference}{}

\gloss{fragmentation}{}

\gloss{full specialization}{}

\gloss{fully associative}{}

\gloss{fully constructed}{fully constructed}

\gloss{function object}{}

\gloss{function scope}{}

\gloss{fundamental alignment}{}

\gloss{fundamental integral types}{}

\gloss{Fundamental type}{TODO}

\gloss{garbage value}{}

\gloss{golden output}{}

\gloss{header-only library}{}

\gloss{hide or hiding}{Function-name \textbf{hiding} occurs when a member function in a derived class has the same name as one in the base class, but it is not overriding it due to a difference in the function signature or because the member function in the base class is not \texttt{virtual}. The hidden member function will \textbf{not} participate in dynamic dispatch; the member function of the base class will be invoked instead when invoked via a pointer or reference to the base class . The same code would have invoked the derived class's implementation had the member function of the base class had been \textbf{overridden} rather than \textbf{hidden}.}

\gloss{higher-order function}{}

\gloss{Hyrum's law}{}

\gloss{Id-expression}{TODO are most commonly \textbf{Identifiers}; other forms include overloaded operator names (in function notation), names of user-defined-conversion or literal operators, and destructor names, and ??template names followed by their argument lists??.}

\gloss{ill formed}{TODO (\texttt{[temp.res]p8})}

\gloss{ill formed, no diagnostic required (IFNDR)}{}

\gloss{imperative}{}

\gloss{implementation-defined}{}

\gloss{implementation inheritance}{}

\gloss{incomplete type}{}

\gloss{inheriting constructors}{}

\gloss{{\tt inline}~{\tt namespace}}{}

\gloss{in-process}{}

\gloss{instantiation time}{}

\gloss{insulate}{}

\gloss{insulating}

\gloss{Insulation}{TODO}

\gloss{Integer literal}{TODO}

\gloss{integral constant expression}{}

\gloss{integral promotion}{}

\gloss{integral type}{}

\gloss{interface inheritance}{}

\gloss{invocable}{}

\gloss{join}{}

\gloss{lambda expression}{}

\gloss{lambda-introducer (adj)}{}

\gloss{leaks}{}

\gloss{linkage}{}

\gloss{list initialization}{}

\gloss{literal type}{}

\gloss{local class}{}

\gloss{local declaration}{}

\gloss{local scope}{}

\gloss{locality of reference}{}

\gloss{long-distance friendship}{}

\gloss{lvalue}{}

\gloss{lvalue reference}{}

\gloss{mangled names}{}

\gloss{mantissa}{}

\gloss{maximal alignment}{}

\gloss{maximally aligned}{}

\gloss{mechanism}{}

\gloss{member initialization lists}{}

\gloss{member initializer list}{}

\gloss{memory leak}{}

\gloss{metafunction}{}

\gloss{mixed-mode builds}{}

\gloss{mix-in}{}

\gloss{modules}{}

\gloss{monotonic allocator}{}

\gloss{most vexing parse}{}

\gloss{move operations}{}

\gloss{Move semantics}{TODO}

\gloss{multithreading contexts}{}

\gloss{Meyers Singleton}{}

\gloss{name mangling}{}

\gloss{natural alignment}{}

\gloss{naturally aligned}{}

\gloss{{\tt new} handler}{}

\gloss{nibbles}{}

\gloss{nonprimitive functionality}{}

\gloss{non-trivial constructor}{}

\gloss{non-trivial special member function}{TODO}

\gloss{null address}{}

\gloss{object invariants}{}

\gloss{ODR}{}

\gloss{ODR-used}{}

\gloss{one definition rule}{}

\gloss{opaque declarations}{}

\gloss{opaque enumeration}{}

\gloss{Overriding}{TODO}

\gloss{parameter pack}{}

\gloss{partial class template specialization}{}

\gloss{partial implementation}{}

\gloss{partially constructed}{}

\gloss{perfectly forwarded}{}

\gloss{physical design}{}

\gloss{Placement}{TODO}

\gloss{placement {\tt new}}{TODO}

\gloss{POD type}{TODO}

\gloss{pointer to member}{}

\gloss{polymorphic memory resource}{}

\gloss{precondition}{}

\gloss{predicate function}{}

\gloss{protocol}{}

\gloss{\romeovalue{prvalue}}{}

\gloss{qualified name}{}

\gloss{RAII}{``Resource Acquisition is Initialization"}

\gloss{Range}{TODO}

\gloss{raw string literals}{}

\gloss{Redundant check}{TODO}

\gloss{reference collapsing}{}

\gloss{Reference type}{TODO}

\gloss{ref-qualifiers}{}

\gloss{regular type}{}

\gloss{return value optimization (RVO)}{}

\gloss{\romeovalue{rvalue}}{}

\gloss{safe-bool-idiom}{}

\gloss{set associative}{}

\gloss{SFINAE}{TODO}

\gloss{SHA}{}

\gloss{shadowed}{}

\gloss{side effects}{}

\gloss{Signature}{TODO}

\gloss{signed integer overflow}{}

\gloss{slicing}{}

\gloss{Special member function}{TODO}

\gloss{standard conversion}{}

\gloss{standard-layout types}{}

\gloss{static data space}{}

\gloss{\lstinline!static! duration}{}

\gloss{{\ttfamily std::unique\_ptr}}{}

\gloss{String literal}{TODO}

\gloss{strong \lstinline!typedef!}{}

\gloss{strong-\lstinline!typedef! idiom}{}

\gloss{Structural inheritance}{TODO}

\gloss{Sum type}{Abstract data type allowing the representation of one of multiple possible alternative types. Each alternative has its own type (and state), and only one alternative can be ``active" at any given point in time. Sum types automatically keep track of which choice is ``active," and properly implement value-sematic special member functions (even for non-trivial types). They can be implemented efficiently as a C++ \texttt{class} using a C++ \texttt{union} and a separate (integral) discriminator. This sort of implementation is commonly referred to as a discriminating (or ``tagged") union.}

\gloss{synthetization}{}

\gloss{template-head}{}

\gloss{Template instantiation time}{TODO}

\gloss{Template instantiation}{TODO}

\gloss{template template parameter}{}

\gloss{test driver}{}

\gloss{thrashing}{}

\gloss{thread pool}{}

\gloss{TLB}{}

\gloss{translation unit (TU)}{}

\gloss{trivial copy constructor}{}

\gloss{trivial operation}{}

\gloss{trivially-constructible}{}

\gloss{trivially copyable}{}

\gloss{Trivial type}{TODO}

\gloss{type alias}{}

\gloss{type deduction}{}

\gloss{typedef}{}

\gloss{type expression}{}

\gloss{type inference}{}

\gloss{type list}{}

\gloss{Type trait}{TODO}

\gloss{UDT}{}

\gloss{Undefined behavior}{TODO}

\gloss{underlying type}{}

\gloss{user-defined type (UDT)}{}

\gloss{user-provided}{}

\gloss{user-provided special member function}{}

\gloss{{\tt using}-declaration}{}

\gloss{{\tt using} directive}{}

\gloss{value}{}

\gloss{value category}{}

\gloss{value constructor}{}

\gloss{value-semantic}{}

\gloss{value-semantic type (VST)}{}

\gloss{variable}{}

\gloss{variadic pack}{}

\gloss{vocabulary type}{}

\gloss{well formed}{}

\gloss{working set}{}

\gloss{zero initialized}{}

%
\cleardoublepage

\newcommand\see[2]{\emph{\seename} #1}  
\newcommand\also[2]{\emph{See also} #1} 
%\newcommand\seename{See} 
\renewcommand\seename{See} % changed becasue of glossary package
\newcommand\seealso[2]{\emph{See also} #1}

\begin{theindex}
\footnotesize
{\sffamily\bfseries{Symbols}}\nopagebreak

  \item \$, 31
  \item \%, 12
  \item \&, 31

  \indexspace
{\sffamily\bfseries\sffamily\bfseries{A}}\nopagebreak

  \item a very long entry to test the column width, 31
  \item anise, 31
  \item apple
    \subitem braebrun, 31
    \subitem cameo, 31
    \subitem fuji, 31
    \subitem gala, 31
    \subitem granny smith, 31
    \subitem red delicious, 31
  \item apricots, 31
  \item avocado, 31

  \indexspace
{\sffamily\bfseries{B}}\nopagebreak

  \item banana, 31
  \item basil, 31
  \item bibendum, 12
    \subitem dapibus, 12
  \item blackberry, 31

  \indexspace
{\sffamily\bfseries{C}}\nopagebreak

  \item cabbage, 31
  \item celery, 31
  \item chervil, 31
  \item chives, 31
  \item cilantro, 31
  \item \protect{\tt codeword}, 12
  \item corn, 31
  \item cucumber, 31

  \indexspace
{\sffamily\bfseries{D}}\nopagebreak

  \item dates, 31
  \item dill, 31

  \indexspace
{\sffamily\bfseries{E}}\nopagebreak

  \item eggplant, 31
  \item endive, 31

  \indexspace
{\sffamily\bfseries{F}}\nopagebreak

  \item fennel, 31
  \item fig, 31
  \item \protect{\tt function}, 31

  \indexspace
{\sffamily\bfseries{G}}\nopagebreak

  \item garlic, 31
  \item grapes, 31

  \indexspace
{\sffamily\bfseries{H}}\nopagebreak

  \item horseradish, 31
  \item how about another long entry, 31
  \item huckleberry, 31

  \indexspace
{\sffamily\bfseries{J}}\nopagebreak

  \item jicama, 31

  \indexspace
{\sffamily\bfseries{K}}\nopagebreak

  \item kale, 31
  \item kiwi, 31

  \indexspace
{\sffamily\bfseries{L}}\nopagebreak

  \item leeks, 31
  \item lemon, 31
  \item lettuce
    \subitem boston bibb, 31
    \subitem iceberg, 31
    \subitem mesclun, 31
    \subitem red leaf, 31
  \item lime, 31
  \item lorem, \see{lobortis}{12}

  \indexspace
{\sffamily\bfseries{M}}\nopagebreak

  \item majoram, 31
  \item mango, 31
  \item maybe another long entry for this test, 31
  \item melon
    \subitem canary, 31
    \subitem cantaloupe, 31
    \subitem honeydew, 31
    \subitem watermelon, 31
  \item mushrooms
    \subitem button, 31
    \subitem porcini, 31
    \subitem portabella, 31
    \subitem shitake, 31

  \indexspace
{\sffamily\bfseries{N}}\nopagebreak

  \item nectarine, 31
  \item nutmeg, 31

  \indexspace
{\sffamily\bfseries{O}}\nopagebreak

  \item okra, 31\pagebreak % you'll need to manually break pages or columns at times to get a nice layout
  \item onion
    \subitem red, 31
    \subitem vidalia, 31
    \subitem yellow, 31
  \item orange, 31

  \indexspace
{\sffamily\bfseries{P}}\nopagebreak

  \item papaya, 31
  \item parsley, 31
  \item peaches, 31
  \item peppers
    \subitem ancho, 31
    \subitem bell, 31
    \subitem habe\~{n}eros, 31
    \subitem jalape\~{n}os, 31
    \subitem pabla\~{n}os, 31
  \item perhaps yet another long entry for this test, 31
  \item plantains, 31
  \item plums, 31
  \item potatoes
    \subitem red-skinned, 31
    \subitem russet, 31
    \subitem yukon gold, 31
  \item pumpkin, 31

  \indexspace
{\sffamily\bfseries{Q}}\nopagebreak

  \item quince, 31

  \indexspace
{\sffamily\bfseries{R}}\nopagebreak

  \item radicchio, 31
  \item radish, 31
  \item raspberry, 31
  \item rosemary, 31
  \item rutabaga, 31\newpage % to end this column short so that it and the last column are balanced

  \indexspace
{\sffamily\bfseries{S}}\nopagebreak

  \item shallots, 31
  \item spinach, 31
  \item squash, 31
  \item still another long entry for column width testing, 31

  \indexspace
{\sffamily\bfseries{T}}\nopagebreak

  \item thyme, 31
  \item tomatillo, 31
  \item tomatoes
    \subitem cherry, 31
    \subitem grape, 31
    \subitem heirloom, 31
    \subitem hybrid, 31
    \subitem roma, 31
  \item \protect{\tt typeof}, 31

  \indexspace
{\sffamily\bfseries{U}}\nopagebreak

  \item ugly fruit, 31

  \indexspace
{\sffamily\bfseries{V}}\nopagebreak

  \item verbena, 31
  \item viverra, \seealso{neque}{12}

  \indexspace
{\sffamily\bfseries{X}}\nopagebreak

  \item xacuti masala, 31

  \indexspace
{\sffamily\bfseries{Y}}\nopagebreak

  \item yams, 31
  \item yet another very long entry for column width test, 31

  \indexspace
{\sffamily\bfseries{Z}}\nopagebreak

  \item zucchini, 31

\end{theindex}


% 
% %
% 
% 
\end{document}
